% ============================================================
%  Plan de Área 2026 — Tecnología e Informática
%  I.E. Sor María Juliana · Cartago, Valle del Cauca
%  PhD. Álvaro Cárdenas Orozco
%  Compilar con: xelatex plan-area-2026.tex (2 pasadas)
% ============================================================
\documentclass[12pt, a4paper]{article}

% --- Codificación y lengua ---
\usepackage{fontspec}
\usepackage{polyglossia}
\setdefaultlanguage{spanish}
\setmainfont[
  BoldFont       = {Helvetica Neue},
  BoldFeatures   = {FakeBold=2.5},
  ItalicFont     = {Helvetica Neue},
  ItalicFeatures = {FakeSlant=0.15}
]{Helvetica Neue}
\setsansfont{Helvetica Neue}

% --- Diseño de página ---
\usepackage[
  top=2.5cm, bottom=2.5cm,
  left=2.5cm, right=2.5cm,
  headheight=15pt
]{geometry}

% --- Color (xcolor antes de tcolorbox) ---
\usepackage[table, dvipsnames]{xcolor}
% Paleta institucional — vino
\definecolor{vinoInst}{RGB}{107, 18, 55}
\definecolor{vinoClaro}{RGB}{245, 228, 238}
% Alias para compatibilidad con seccion8-tablas.tex (generado)
\colorlet{azulInst}{vinoInst}
\colorlet{azulClaro}{vinoClaro}
% Ejes curriculares
\definecolor{grisClaro}{RGB}{242, 242, 242}
\definecolor{verdePC}{RGB}{0, 100, 60}
\definecolor{verdePCclaro}{RGB}{215, 240, 225}
\definecolor{moradoSEL}{RGB}{90, 30, 120}
\definecolor{moradoSELclaro}{RGB}{237, 220, 248}
% Anclajes ancestrales — paleta sobria armónica con el vino institucional
\definecolor{p1bg}{RGB}{244, 232, 237}  % rosa vino muy claro
\definecolor{p1fg}{RGB}{120, 35, 70}
\definecolor{p2bg}{RGB}{245, 238, 226}  % arena / dorado claro
\definecolor{p2fg}{RGB}{130, 95, 45}
\definecolor{p3bg}{RGB}{232, 231, 239}  % malva gris claro
\definecolor{p3fg}{RGB}{85, 70, 110}
% Triángulo filosófico
\definecolor{verdeTeal}{RGB}{0, 120, 100}
\definecolor{azulFloridi}{RGB}{20, 50, 120}
\definecolor{dorado}{RGB}{180, 140, 0}

% --- Tcolorbox + TikZ (grafismo de portada) ---
\usepackage[most]{tcolorbox}
\usepackage{tikz}
\usetikzlibrary{calc}

% --- Tablas ---
\usepackage{longtable}
\usepackage{multirow}
\usepackage{multicol}
\usepackage{booktabs}
\usepackage{array}
\usepackage{tabularx}
\usepackage{makecell}
\usepackage{colortbl}

% --- Landscape para mallas ---
\usepackage{pdflscape}

% --- Encabezados y pies ---
\usepackage{fancyhdr}
\pagestyle{fancy}
\fancyhf{}
\fancyhead[L]{\small\color{vinoInst}\textbf{Plan de Área · Tecnología e Informática 2026}}
\fancyhead[R]{\small\color{vinoInst}I.E. Sor María Juliana · Cartago}
\fancyfoot[L]{\footnotesize\color{black!42}Área de Tecnología e Informática}
\fancyfoot[C]{\small\color{vinoInst}\textbf{\thepage}}
\fancyfoot[R]{\footnotesize\color{black!42}Cartago · Valle del Cauca}
% Reglas finas en gris claro (V6), consistentes con las tablas
\renewcommand{\headrule}{{\color{lineGris}\hrule height 0.5pt}}
\renewcommand{\footrule}{{\color{lineGris}\hrule height 0.4pt}}
\renewcommand{\headrulewidth}{0.5pt}
\renewcommand{\footrulewidth}{0.4pt}

% --- Secciones con color institucional ---
\usepackage{titlesec}
% Sección: número en badge vino + título + regla gris fina (V5)
\titleformat{\section}
  {\Large\bfseries\color{vinoInst}}
  {\colorbox{vinoInst}{\textcolor{white}{\rule[-2.5pt]{0pt}{13pt}\,\thesection\,}}}
  {0.6em}{}
  [{\color{lineGris}\titlerule[1.2pt]}]
\titlespacing*{\section}{0pt}{2.4ex plus 1ex minus .2ex}{1.4ex plus .2ex}
\titleformat{\subsection}
  {\color{vinoInst}\normalsize\bfseries}{\color{vinoInst!75!black}\thesubsection}{0.5em}{}
\titleformat{\subsubsection}
  {\normalsize\bfseries}{\color{vinoInst!70!black}\thesubsubsection}{0.5em}{}

% --- Índice con jerarquía visual (V4) ---
\usepackage{tocloft}
\renewcommand{\cftsecfont}{\bfseries\color{vinoInst}}
\renewcommand{\cftsecpagefont}{\bfseries\color{vinoInst}}
\renewcommand{\cftsubsecfont}{\color{black!72}}
\renewcommand{\cftsubsecpagefont}{\color{black!55}}
\setlength{\cftbeforesecskip}{5pt}
\setlength{\cftbeforesubsecskip}{1.6pt}

% --- Listas y espaciado ---
\usepackage{enumitem}
\setlist[itemize]{leftmargin=1.5em, itemsep=2pt, topsep=4pt}
\setlist[enumerate]{leftmargin=1.5em, itemsep=2pt, topsep=4pt}
\usepackage{parskip}
\usepackage{needspace}
\setlength{\parindent}{0pt}
\setlength{\parskip}{6pt}

% --- Hipervínculos ---
\usepackage[hidelinks, colorlinks=false]{hyperref}

% --- Comandos utilitarios ---
\newcommand{\encabezadoMalla}{%
  \rowcolor{azulInst}
  \color{white}\footnotesize\textbf{\makecell{ESTÁNDARES,\\LINEAMIENTOS\\U ORIENTACIONES}} &
  \color{white}\footnotesize\textbf{COMPETENCIAS} &
  \color{white}\footnotesize\textbf{APRENDIZAJES} &
  \multicolumn{3}{c|}{\color{white}\textbf{DESCRIPTORES DE DESEMPEÑO}%
    \newline\color{white}\textit{\small Metas del aprendizaje por periodo}}\\
  \cline{4-6}
  \rowcolor{azulClaro}
  & & &
  \textbf{COGNITIVOS} &
  \textbf{PROCEDIMENTALES} &
  \textbf{ACTITUDINALES}\\
  \hline
}

\newcommand{\filaPC}[4]{%
  \rowcolor{verdePCclaro}
  \makecell[l]{%
    \textbf{\color{verdePC}Pensamiento}\\\textbf{\color{verdePC}Computacional}\\
    \footnotesize Wing (2006)\\
    \footnotesize MEN Or. 2022\\
    \footnotesize\textit{#1}%
  } & #2 & #3 & \multicolumn{3}{p{8.5cm}|}{#4}\\
  \hline
}

\newcommand{\filaSEL}[4]{%
  \rowcolor{moradoSELclaro}
  \makecell[l]{%
    \textbf{\color{moradoSEL}Habilidades}\\\textbf{\color{moradoSEL}Socioemocionales}\\
    \footnotesize Ley 2383/2024\\
    \footnotesize Cátedra de la Paz\\
    \footnotesize\textit{#1}%
  } & #2 & #3 & \multicolumn{3}{p{8.5cm}|}{#4}\\
  \hline
}

\newcolumntype{M}[1]{>{\centering\arraybackslash}p{#1}}
\newcolumntype{L}[1]{>{\raggedright\arraybackslash}p{#1}}

% --- Estilo editorial de tablas (V1): líneas finas grises + aire + zebra ---
\definecolor{lineGris}{RGB}{212, 210, 218}
\definecolor{zebraGris}{RGB}{247, 245, 249}
\arrayrulecolor{lineGris}
\setlength{\arrayrulewidth}{0.4pt}
\renewcommand{\arraystretch}{1.4}

% --- Tarjetas y cifras (V3: ritmo visual) ---
\newcommand{\tarjeta}[2]{%
  \begin{minipage}[t]{0.305\textwidth}
  \begin{tcolorbox}[enhanced, colback=vinoClaro, colframe=vinoInst, boxrule=0.8pt,
    arc=6pt, left=7pt, right=7pt, top=7pt, bottom=7pt,
    equal height group=tarjetas, valign=top]
  {\color{vinoInst}\textbf{#1}}\\[3pt]
  {\footnotesize #2}
  \end{tcolorbox}
  \end{minipage}%
}
\newcommand{\cifra}[2]{%
  \begin{minipage}[t]{0.235\textwidth}\centering
  {\fontsize{36}{40}\selectfont\bfseries\color{vinoInst}#1}\\[1pt]
  {\footnotesize\color{black!72}#2}
  \end{minipage}%
}
\newcommand{\callout}[2]{%
  \begin{tcolorbox}[enhanced, colback=vinoClaro, colframe=vinoInst, boxrule=0pt,
    leftrule=3.5pt, arc=2pt, left=11pt, right=11pt, top=8pt, bottom=8pt]
  {\color{vinoInst}\textbf{#1}}\, #2
  \end{tcolorbox}}

% ============================================================
\begin{document}
% ============================================================

% ---- PORTADA (tcolorbox — tarjeta vino sobre fondo blanco) ----
\begin{titlepage}
\newgeometry{top=1.2cm, bottom=1.2cm, left=1.2cm, right=1.2cm}
\noindent
\begin{tcolorbox}[
  enhanced,
  colback=vinoInst,
  colframe=white,
  boxrule=2pt,
  arc=20pt,
  borderline={2pt}{6pt}{white, dashed},
  left=2cm, right=2cm, top=2cm, bottom=2cm,
  width=\linewidth,
  height=26cm,
  underlay={%
    \colorlet{net}{white!18!vinoInst}%
    \tikzset{nd/.style={fill=net, circle, inner sep=1.1pt}}%
    % --- constelación superior derecha ---
    \begin{scope}[shift={($(interior.north east)+(-0.6,-0.6)$)}]
      \node[nd](a)at(0,0){}; \node[nd](b)at(-2,-0.4){}; \node[nd](c)at(-3.6,0.2){};
      \node[nd](d)at(-1.3,-1.9){}; \node[nd](e)at(-3.1,-1.5){}; \node[nd](f)at(-0.4,-2.7){};
      \node[nd](g)at(-2.4,-3.1){};
      \draw[net, line width=0.35pt] (a)--(b)--(c) (b)--(d)--(e) (d)--(f) (a)--(d) (e)--(g)--(d);
    \end{scope}
    % --- constelación inferior izquierda ---
    \begin{scope}[shift={($(interior.south west)+(0.6,0.6)$)}]
      \node[nd](p)at(0,0){}; \node[nd](q)at(2,0.4){}; \node[nd](r)at(3.5,-0.1){};
      \node[nd](s)at(1.2,1.9){}; \node[nd](t)at(3,1.5){}; \node[nd](u)at(0.5,2.7){};
      \node[nd](v)at(2.5,3.0){};
      \draw[net, line width=0.35pt] (p)--(q)--(r) (q)--(s)--(t) (s)--(u) (p)--(s) (t)--(v)--(s);
    \end{scope}
  }
]
\centering
\vspace{0.3cm}
{\fontsize{11}{13}\selectfont\color{white!55!vinoInst}\textbf{D O C U M E N T O \quad R E C T O R \quad I N S T I T U C I O N A L}}\\[0.55cm]
{\color{dorado!85!white}\rule{0.46\linewidth}{1.3pt}}\\[1.0cm]
{\fontsize{38}{44}\bfseries\color{white}Plan de Área}\\[0.15cm]
{\fontsize{84}{90}\bfseries\color{dorado!85!white}2026}\\[0.8cm]
{\LARGE\color{white}\textbf{TECNOLOGÍA E INFORMÁTICA}}\\[0.35cm]
{\normalsize\color{white!75!vinoInst}Grados 6° a 11° · Básica Secundaria y Media}\\[1.2cm]
{\color{white!30!vinoInst}\rule{0.30\linewidth}{0.4pt}}\\[0.6cm]
{\large\color{white}\textbf{Institución Educativa Sor María Juliana}}\\[0.22cm]
{\normalsize\color{white}Cartago · Valle del Cauca · Colombia}\\[0.22cm]
{\small\color{white!70!vinoInst}Área de Tecnología e Informática}\\[1.0cm]
{\normalsize\color{white}Modelo pedagógico \textbf{MILC v3}}\\[0.12cm]
{\small\itshape\color{white!72!vinoInst}Modelo de Investigación Liberadora y Científica}
\vfill
\begin{tcolorbox}[
  colback=white!10!vinoInst,
  colframe=dorado!55!vinoInst,
  boxrule=1pt,
  arc=10pt,
  width=0.84\linewidth,
  halign=center,
  top=10pt, bottom=10pt
]
\centering
{\color{dorado!88!white}\textbf{216 piezas curriculares activas}}\\[0.3cm]
{\small\color{white}180 guías \,\textperiodcentered\, 18 proyectos integradores \,\textperiodcentered\, 18 exámenes}\\[0.16cm]
{\footnotesize\color{white!72!vinoInst}Plataforma Conéctate · alguarito.github.io/plataformaconectate}
\end{tcolorbox}

\vspace{0.7cm}
{\color{white!30!vinoInst}\rule{0.46\linewidth}{0.4pt}}\\[0.5cm]
{\normalsize\color{white}\textbf{Equipo del Área de Tecnología e Informática}}\\[0.15cm]
{\small\color{white!72!vinoInst}Docentes del área · I.E. Sor María Juliana · Cartago, 2026}
\end{tcolorbox}
\restoregeometry
\end{titlepage}
\nopagecolor

% ---- APERTURA / MANIFIESTO ----
\newpage
\thispagestyle{empty}
\vspace*{0.4cm}

% --- Cabecera de acento ---
\noindent{\color{dorado}\rule{3cm}{3pt}}\hfill
{\footnotesize\color{vinoInst!55!white}\textbf{D O C U M E N T O\, \,R E C T O R}}

\vspace{0.7cm}
\noindent{\fontsize{32}{36}\selectfont\bfseries\color{vinoInst}Por qué existe\\[2pt]este plan}

\vfill

% --- Lead / hook emotivo (cursiva, voz) ---
\noindent\begin{minipage}{0.95\linewidth}\raggedright
{\fontsize{17}{26}\selectfont\itshape\color{vinoInst!88!black}Cada año, decenas de
jóvenes de Cartago encienden por primera vez una pantalla que no fabricaron,
gobernada por reglas que no escribieron y alimentada con sus datos sin que lo
sepan.}
\end{minipage}

\vfill

% --- Cuerpo (justificado) ---
\noindent\begin{minipage}{0.93\linewidth}\setlength{\parskip}{10pt}\normalsize\color{black!82}
La pregunta que abre este plan no es \emph{qué herramientas enseñarles}; es más
honda: ¿los preparamos para obedecer ese mundo, o para entenderlo, cuestionarlo y
rehacerlo?

{\large\color{vinoInst}\bfseries Elegimos lo segundo. Sin matices.}

Por eso no empezamos por el computador, sino por lo que cada estudiante ya sabe sin
saber que lo sabe: el oficio del relojero que mide el tiempo con las manos, el
cálculo callado de la tejedora, la voz del pregonero que ordena un barrio entero.
Sobre ese saber —el suyo, el de los suyos— se levanta todo lo demás; y la
tecnología deja de ser algo que se \emph{padece} para volverse algo que se
\emph{piensa}.

Y no es una promesa escrita al aire: detrás de cada página hay \textbf{\color{vinoInst}216
piezas ya construidas, probadas y publicadas} —180 guías, 18 proyectos y 18
exámenes— hechas para abrirse incluso \emph{sin señal}, porque ningún muchacho
debería quedar afuera por no tener datos.
\end{minipage}

\vfill

% --- Cierre (banner) ---
\noindent{\color{dorado}\rule{2.6cm}{2.5pt}}

\vspace{0.45cm}
\noindent{\fontsize{27}{31}\selectfont\bfseries\color{vinoInst}Formar autores,\\[2pt]no usuarios.}

\vspace{0.5cm}
\noindent\begin{minipage}{0.86\linewidth}\raggedright
{\large\color{black!75}No es un lema: es lo que está en juego en cada clase —que
estos jóvenes dejen de ser el \emph{público} de la tecnología para convertirse en
sus \textbf{\color{vinoInst}autores}.}
\end{minipage}

\vspace*{2cm}
\newpage
\setcounter{tocdepth}{2}
\tableofcontents
\newpage

% ---- CÓMO LEER ESTE PLAN ----
\thispagestyle{empty}
\vspace*{0.5cm}
\noindent{\Large\bfseries\color{vinoInst}Cómo leer este plan}

\vspace{0.4cm}
\noindent Este documento sirve a varios lectores. Según quién seas, este es un
buen punto de entrada:

\medskip
\begin{center}
\begin{tabular}{|L{4cm}|L{10cm}|}
\hline
\rowcolor{vinoInst}
\color{white}\textbf{Si eres\ldots} & \color{white}\textbf{Empieza por} \\
\hline
\textbf{Docente del área} & El Marco teórico (sección\,5) y la Metodología (sección\,11) para el «cómo»; el Anexo A para la malla de tu grado. \\
\hline
\textbf{Directivo o coordinador} & La Presentación (sección\,3), los Propósitos y metas (sección\,9) y el Sistema de evaluación (sección\,15). \\
\hline
\textbf{Par evaluador o visitante} & El Marco legal (sección\,2) con su síntesis norma–malla (sección\,2.5), y el Diseño curricular (sección\,10). \\
\hline
\textbf{Familia o acudiente} & La Presentación (sección\,3) y la Comunicación con familias (sección\,20). \\
\hline
\end{tabular}
\end{center}
\newpage

% Zebra striping activo en el cuerpo (se desactiva antes del Anexo A)
\rowcolors{2}{white}{zebraGris}

% ============================================================
\section{Identificación institucional}
% ============================================================

\begin{center}
\begin{tabular}{|L{5cm}|L{9cm}|}
\hline
\rowcolor{vinoInst}\color{white}\textbf{Campo} & \color{white}\textbf{Dato} \\
\hline
Institución educativa & I.E. Sor María Juliana \\
\hline
Municipio & Cartago, Valle del Cauca \\
\hline
Área & Tecnología e Informática \\
\hline
Grados & 6° a 11° (Básica Secundaria y Media) \\
\hline
Docentes responsables & Equipo del Área de Tecnología e Informática \\
\hline
Modelo pedagógico & MILC v3 — Modelo de Investigación Liberadora y Científica \\
\hline
Año lectivo & 2026 \\
\hline
Periodos & 3 periodos por grado (enero–noviembre) \\
\hline
Total piezas curriculares & 216 (180 guías + 18 proyectos integradores + 18 exámenes) \\
\hline
Plataforma digital & alguarito.github.io/plataformaconectate \\
\hline
\end{tabular}
\end{center}

% ============================================================
\section{Marco legal y referentes normativos}
% ============================================================

\subsection{Normativa nacional de base}

\begin{longtable}{|L{6.5cm}|L{7.5cm}|}
\hline
\rowcolor{vinoInst}
\color{white}\textbf{Norma / Documento} &
\color{white}\textbf{Articulación con el área} \\
\hline
\endhead
\textbf{Ley 115 de 1994 — Ley General de Educación} &
Arts. 5, 22, 23 y 30: tecnología e informática como área obligatoria
y fundamental; formación en uso de la ciencia y la técnica. \\
\hline
\textbf{Decreto 1860 de 1994} &
Reglamenta la Ley 115; estructura del PEI, áreas obligatorias y plan de estudios. \\
\hline
\textbf{Decreto 1290 de 2009} &
Sistema Institucional de Evaluación (SIE); escala de valoración y criterios
de desempeño. \\
\hline
\textbf{Ley 1341 de 2009 (TIC)} &
Acceso universal y uso responsable de las TIC; articulación con la formación
ciudadana digital. \\
\hline
\textbf{Ley 1581 de 2012 — Habeas Data} &
Protección de datos personales; competencia ciudadana digital transversal. \\
\hline
\textbf{Ley 23 de 1982 — Derechos de Autor} &
Propiedad intelectual en producción de contenidos digitales; eje del
componente de autoría responsable con IA (G10°–G11°). \\
\hline
\textbf{Ley 1014 de 2006 — Fomento al Emprendimiento} &
Base legal del componente de microemprendimiento digital en G10° P3
y G11° P3. \\
\hline
\end{longtable}

\subsection{Lineamientos curriculares MEN — Tecnología e Informática}

\begin{longtable}{|L{6.5cm}|L{7.5cm}|}
\hline
\rowcolor{vinoInst}
\color{white}\textbf{Documento} &
\color{white}\textbf{Qué aporta al área} \\
\hline
\endhead
\textbf{MEN (1996). Lineamientos Curriculares — Tecnología e Informática.} &
Primer marco oficial del área en Colombia. Define la tecnología
como proceso de diseño e innovación humana. \\
\hline
\textbf{MEN (2008). Guía No. 30: \textit{Ser competente en tecnología.}} &
Documento rector vigente. Define los cuatro componentes del área:
(1) Naturaleza y evolución, (2) Apropiación y uso,
(3) Solución de problemas, (4) Tecnología y sociedad.
Este plan los toma como base y agrega PC y SEL como componentes 5 y 6. \\
\hline
\textbf{MEN (2022). Orientaciones Curriculares — Tecnología e Informática.} &
Actualiza la Guía 30. Integra el Pensamiento Computacional como
estrategia didáctica transversal con énfasis en programación (Scratch,
algoritmos) y enfoques maker, STEM+ y aprendizaje basado en juegos.
Este plan lo materializa como quinto componente explícito en cada malla. \\
\hline
\end{longtable}

\subsection{Marco normativo de Inteligencia Artificial}

\begin{longtable}{|L{6.5cm}|L{7.5cm}|}
\hline
\rowcolor{vinoInst}
\color{white}\textbf{Documento / Política} &
\color{white}\textbf{Qué aporta al área} \\
\hline
\endhead
\textbf{CONPES 4144 de 2025 — Política Nacional de IA.} &
Hoja de ruta vigente de Colombia en IA. Seis ejes: ética y gobernanza,
datos e infraestructura, I+D+I, capacidades, riesgos y adopción.
Sustenta el componente de ciudadanía digital crítica en todos los grados. \\
\hline
\textbf{Día de la IA en tu Colegio — CPE / MEN / MinTIC, 2025.} &
Primera edición: 2 de octubre de 2025; segunda edición: 2026.
Este plan lo articula como cierre de sustentación del proyecto integrador
de G10° P1. \\
\hline
\textbf{UNESCO (2021). Recomendación sobre Ética de la IA.} &
11 principios de ética en IA. Obliga a garantizar alfabetización en IA
e integrar ética de la IA en currículos escolares. El CONPES 4144/2025
los incorpora. \\
\hline
\textbf{UNESCO (2024). Marco de competencias de IA para estudiantes.} &
4 dimensiones × 3 niveles de progresión. Este plan mapea las dimensiones
técnicas al eje PC y las éticas al eje SEL en los grados de Media. \\
\hline
\end{longtable}

\subsection{Lineamientos para la Educación Socioemocional (SEL)}

Las habilidades socioemocionales son el \textbf{sexto componente curricular
explícito} de este plan. Su inclusión responde a un mandato legal colombiano
actual y a evidencia científica sólida sobre su impacto en el aprendizaje.

\begin{longtable}{|L{6.5cm}|L{7.5cm}|}
\hline
\rowcolor{vinoInst}
\color{white}\textbf{Documento / Política} &
\color{white}\textbf{Qué aporta al área} \\
\hline
\endhead
\textbf{Ley 1620 de 2013 — Convivencia Escolar.} &
Crea el Sistema Nacional de Convivencia Escolar. Sustenta la formación
en derechos humanos y ciudadanía digital como eje transversal. \\
\hline
\textbf{Ley 1732 de 2014 — Cátedra de la Paz.} &
Establece la Cátedra de la Paz como obligatoria. Sus temáticas
(dilemas morales, resolución de conflictos, proyecto de vida)
se articulan con el componente \textit{Tecnología y sociedad} y
el cierre filosófico MILC. \\
\hline
\textbf{Ley 2294 de 2023 — PND «Colombia Potencia Mundial de la Vida».} &
Establece la estrategia CRESE (ciudadana, reconciliación, socioemocional,
antirracista, cambio climático). Eleva la educación socioemocional a
política de Estado 2022–2026. \\
\hline
\textbf{Ley 2383 de 19 de julio de 2024.} \textit{Por la cual se promueve
la educación socioemocional de NNA en instituciones educativas.} &
\textbf{Primera ley colombiana que denomina y mandata explícitamente
la educación socioemocional.} Obliga al MEN a formular lineamientos
pedagógicos. Este plan anticipa ese mandato con la fila SEL en cada malla. \\
\hline
\textbf{Ley 2491 de 23 de julio de 2025.} \textit{Por la cual se incorporan
competencias socioemocionales al PEI.} &
Obliga a incorporar competencias socioemocionales al PEI de cada institución.
Consolida el marco legal SEL como obligación institucional. \\
\hline
\textbf{CASEL (2020). \textit{SEL Framework.}} &
Marco internacional: 5 competencias (autoconciencia, autogestión,
conciencia social, habilidades relacionales, toma de decisiones responsable).
Adoptado como referente técnico, subordinado a la legislación colombiana. \\
\hline
\end{longtable}

\subsection{Síntesis: de la norma a la malla curricular}

Esta tabla cierra el marco legal mostrando \textbf{dónde se materializa cada
referente}: conecta cada eje curricular con las normas que lo sustentan y con la
fila exacta de la malla donde el docente lo encuentra. Es el puente entre la
norma y el aula.

\begin{center}
\begin{tabular}{|L{3.5cm}|L{5cm}|L{5cm}|}
\hline
\rowcolor{vinoInst}
\color{white}\textbf{Eje curricular} &
\color{white}\textbf{Referentes normativos principales} &
\color{white}\textbf{Componente en la malla} \\
\hline
4 Componentes MEN &
Guía 30 (2008) · Ley 115/1994 &
Filas 1–4 de cada tabla \\
\hline
Pensamiento Computacional &
MEN Or. 2022 · Wing (2006) · CONPES 4144 · UNESCO Frameworks &
Fila 5 (verde) \\
\hline
Habilidades Socioemocionales &
Ley 2383/2024 · Ley 1732/2014 · CASEL 2020 &
Fila 6 (violeta) \\
\hline
IA Ética y Día IA &
UNESCO Rec. 2021 · CONPES 4144 · Día IA 2026 &
Transversal G10°–G11° \\
\hline
MILC v3 &
Modelo institucional propio · Dussel · Marco Aurelio · Floridi &
Metodología + Evaluación liberadora \\
\hline
\end{tabular}
\end{center}

% ============================================================
\section{Presentación del área}
% ============================================================

El área de Tecnología e Informática de la I.E. Sor María Juliana forma ciudadanos
críticos, creativos y responsables en el uso de la tecnología como herramienta de
transformación personal y social. No se reduce a la enseñanza de herramientas
digitales: es un espacio de pensamiento, producción y liberación donde el
estudiante pasa de \textbf{usuario} a \textbf{autor} de tecnología.

La identidad del área descansa en tres rasgos que la distinguen:

\medskip
\noindent
\tarjeta{Modelo propio · MILC v3}{Parte del saber del estudiante y su comunidad,
lo confronta con un problema real y lo acompaña a crear una solución con criterio.}%
\hfill
\tarjeta{Dos ejes transversales}{Pensamiento Computacional y Habilidades
Socioemocionales atraviesan todos los grados y periodos (ver Marco conceptual, sección\,7).}%
\hfill
\tarjeta{Despliegue \textit{offline-first}}{Plataforma Conéctate: 180 guías, 18
proyectos y 18 exámenes, accesibles sin conexión permanente.}

\bigskip
\noindent\begin{tcolorbox}[enhanced, colback=white, colframe=vinoInst, boxrule=0.8pt,
  arc=6pt, left=10pt, right=10pt, top=12pt, bottom=12pt]
\cifra{180}{guías de aprendizaje}\hfill%
\cifra{18}{proyectos integradores}\hfill%
\cifra{18}{exámenes de periodo}\hfill%
\cifra{216}{piezas activas}
\end{tcolorbox}

% ============================================================
\section{Justificación}
% ============================================================

La sociedad del siglo XXI exige ciudadanos capaces de operar, cuestionar y
transformar entornos tecnológicos cada vez más complejos. En Cartago y el Valle
del Cauca, estudiantes de estratos 1, 2 y 3 con acceso limitado a datos necesitan
una formación que no solo los habilite como usuarios, sino como \textbf{autores,
diseñadores y críticos} de los sistemas que los rodean.

Tres urgencias contemporáneas hacen impostergable esta formación hoy:

\medskip
\noindent
\tarjeta{1 · IA generativa}{Sin criterio para usarla, cuestionarla y declararla,
los jóvenes quedan a merced de sistemas que no comprenden ni controlan.}%
\hfill
\tarjeta{2 · Brecha digital}{En un territorio de recursos y conectividad
limitados, formar solo usuarios —y no autores— reproduce la desigualdad en lugar
de cerrarla.}%
\hfill
\tarjeta{3 · Bienestar digital}{La salud mental, la convivencia en línea y la
decisión responsable son hoy un mandato legal (Ley 2383 de 2024) que el área
atiende desde el aula.}

% ============================================================
\section{Marco teórico y pedagógico}
% ============================================================

\subsection{MILC v3 — Modelo de Investigación Liberadora y Científica}

\textbf{¿Qué es?} El \textbf{MILC v3} (Modelo de Investigación Liberadora y
Científica) es la especificidad del modelo pedagógico institucional interactivo
y crítico de la I.E. Sor María Juliana, aplicada al área de Tecnología e
Informática. No es un marco importado: es una pedagogía propia, construida desde
el territorio de Cartago y el Valle del Cauca, que integra la pedagogía crítica
(Freire), la filosofía de la liberación (Dussel), el estoicismo aplicado
(Marco Aurelio) y la filosofía de la información (Floridi).

\textbf{¿Por qué este modelo?} La enseñanza tradicional de la informática reduce
al estudiante a un \textit{usuario} que reproduce instrucciones: abrir un
programa, copiar un formato, seguir un tutorial. En un municipio donde la mayoría
de los jóvenes accede a la tecnología desde un celular de gama media y con datos
limitados, formar solo usuarios reproduce la brecha digital y los deja a merced
de plataformas que no controlan ni comprenden. El MILC responde a esa realidad:
parte de un \textbf{saber propio} del estudiante, lo confronta con un problema
real y lo acompaña a producir una solución con criterio. La tecnología deja de
ser algo que se padece para convertirse en algo que se piensa, se cuestiona y
se crea.

\textbf{¿Para qué?} El modelo persigue tres propósitos concretos: (1) que el
estudiante pase de \textit{consumidor} a \textbf{autor} de tecnología; (2) que
cada aprendizaje técnico quede \textbf{anclado} en su realidad cultural y ética,
y no en abstracto; y (3) que desarrolle un \textbf{criterio propio} para decidir,
en un mundo de inteligencia artificial y algoritmos, qué usar, cómo y para qué.

\textbf{¿Qué impacto tiene?} En el aula, el MILC se traduce en estudiantes que
escriben, argumentan y producen —no solo que manipulan software—; en guías donde
cada actividad termina en un producto verificable del cuaderno; y en una
evaluación que mide criterio y transformación, no memorización. Su efecto se
hace visible en los \textbf{proyectos integradores} de cierre de periodo, donde
el estudiante resuelve un problema de su comunidad con una solución tecnológica
propia y la sustenta ante sus pares.

Las \textbf{Orientaciones Curriculares MEN (2022)} establecen que los procesos
pedagógicos del área deben orientarse hacia la resolución de problemas
auténticos, el diseño de soluciones con tecnología y la participación crítica y
ética en entornos digitales. El MILC v3 materializa estas directrices mediante
cuatro fases que articulan la acción del estudiante, la del docente, los cuatro
componentes MEN y el Pensamiento Computacional como estrategia transversal:

\begin{center}
\setlength{\tabcolsep}{4pt}
\begin{tabular}{|L{3cm}|L{4cm}|L{4cm}|L{3cm}|}
\hline
\rowcolor{vinoInst}
\color{white}\textbf{Fase MILC} &
\color{white}\textbf{Qué hace el estudiante} &
\color{white}\textbf{Qué hace el docente} &
\color{white}\textbf{Herramienta} \\
\hline
\textbf{1. Escucha} &
Conecta con el saber previo y el anclaje ancestral &
Activa el saber local; formula la pregunta-puente &
Estación 1 — Saber ancestral \\
\hline
\textbf{2. Sistematización} &
Lee, observa, explora; construye el marco conceptual &
Modela con ejemplos contextualizados de Cartago &
Bloques LEE · OBSERVA · EXPLORA \\
\hline
\textbf{3. Praxis} &
Produce, aplica, transforma; lleva al cuaderno &
Orienta el ciclo IDENTIFICA $\rightarrow$ CREA &
Bloque ACTIVIDAD (6 verbos Bloom) \\
\hline
\textbf{4. Evaluación liberadora} &
Juzga con criterio propio; reflexiona metacognitivamente &
Cierre filosófico y evaluación en 5 dimensiones &
Bloques REFLEXIONA · Evaluación 5D \\
\hline
\end{tabular}
\end{center}

\medskip
\noindent\textit{En síntesis: el MILC convierte cada clase en un pequeño acto de
investigación —escuchar el saber propio, sistematizarlo, transformarlo y juzgarlo
con criterio—. No se enseña tecnología \emph{sobre} el estudiante; se construye
\emph{con} él.}

\subsection{Triángulo filosófico del área}

El área opera sobre tres pilares filosóficos complementarios que enmarcan cada
guía de aprendizaje y cada cierre de periodo:

\medskip
\noindent
\begin{minipage}[t]{0.31\textwidth}
\begin{tcolorbox}[
  enhanced,
  colback=vinoClaro,
  colframe=vinoInst,
  coltitle=white,
  colbacktitle=vinoInst,
  title={\small\textbf{Enrique Dussel}},
  fonttitle=\bfseries,
  arc=8pt,
  boxrule=2pt,
  left=4pt, right=4pt, top=4pt, bottom=4pt
]
{\small\textbf{Filosofía de la Liberación}

La tecnología no es neutral: puede liberar o dominar según quién la controla.
El estudiante de Cartago es sujeto con voz propia, no objeto de la modernidad
tecnológica.}
\end{tcolorbox}
\end{minipage}
\hfill
\begin{minipage}[t]{0.31\textwidth}
\begin{tcolorbox}[
  enhanced,
  colback=white!90!verdeTeal,
  colframe=verdeTeal,
  coltitle=white,
  colbacktitle=verdeTeal,
  title={\small\textbf{Marco Aurelio}},
  fonttitle=\bfseries,
  arc=8pt,
  boxrule=2pt,
  left=4pt, right=4pt, top=4pt, bottom=4pt
]
{\small\textbf{Estoicismo Aplicado}

Disciplina del tiempo y la atención. En el trabajo técnico, lo que no depende de ti
(la red, el dispositivo) no te perturba; lo que sí depende de ti
(tu proceso, tu criterio) lo perfeccionas.}
\end{tcolorbox}
\end{minipage}
\hfill
\begin{minipage}[t]{0.31\textwidth}
\begin{tcolorbox}[
  enhanced,
  colback=white!90!azulFloridi,
  colframe=azulFloridi,
  coltitle=white,
  colbacktitle=azulFloridi,
  title={\small\textbf{Luciano Floridi}},
  fonttitle=\bfseries,
  arc=8pt,
  boxrule=2pt,
  left=4pt, right=4pt, top=4pt, bottom=4pt
]
{\small\textbf{Infoesfera}

Habitas un ambiente simbólico-informacional que moldea tu identidad, tu
privacidad y tus decisiones. Ser ciudadano digital es ser consciente de
ese ambiente.}
\end{tcolorbox}
\end{minipage}
\medskip

\textbf{Del concepto al aula — el cierre filosófico.} El triángulo no es decorativo:
se traduce en la \textbf{estación de cierre} de cada guía, donde el estudiante
piensa el tema técnico desde los tres pilares. Las preguntas se \textbf{adaptan
por grado}: en 6°–8° son cotidianas y concretas; en 9°–11° incorporan los autores
y exigen argumentación. Ejemplo sobre un mismo tema (\textit{huella digital}):

\begin{center}
\begin{tabular}{|L{2.6cm}|L{5.5cm}|L{5.5cm}|}
\hline
\rowcolor{vinoInst}
\color{white}\textbf{Pilar} &
\color{white}\textbf{Cierre en 6°–8° (cotidiano)} &
\color{white}\textbf{Cierre en 9°–11° (con autor)} \\
\hline
\textbf{Dussel} (liberación) &
¿Quién gana cuando regalas tus datos sin darte cuenta? &
¿Tu huella digital te hace sujeto o producto? Discute desde la idea de dominación
tecnológica de Dussel. \\
\hline
\textbf{Marco Aurelio} (estoicismo) &
De tu huella, ¿qué sí depende de ti y qué no? &
¿Qué control real tienes sobre tu rastro? Aplica la dicotomía estoica de lo que
depende y no depende de ti. \\
\hline
\textbf{Floridi} (infoesfera) &
Si tus datos viven en internet para siempre, ¿quién eres en línea? &
¿Cómo moldea la infoesfera tu identidad y privacidad? Argumenta desde Floridi. \\
\hline
\end{tabular}
\end{center}

\subsection{Referentes integradores}

\begin{center}
\begin{tabular}{|L{3cm}|L{3cm}|L{8cm}|}
\hline
\rowcolor{vinoInst}
\color{white}\textbf{Referente} &
\color{white}\textbf{Eje} &
\color{white}\textbf{Contribución al área de Tecnología e Informática} \\
\hline
Paulo Freire (1970) & Pedagogía crítica &
El estudiante como sujeto político que problematiza su realidad tecnológica. \\
\hline
Enrique Dussel (1977) & Filosofía de la liberación &
La tecnología como herramienta de emancipación o dominación; reconocimiento del
«Otro» como sujeto del aprendizaje. \\
\hline
Luciano Floridi (2014) & Filosofía de la información &
La infoesfera como ambiente de vida; ética de los datos, identidad digital y
responsabilidad en la red. \\
\hline
Jeannette Wing (2006) & Pensamiento Computacional &
5 habilidades cognitivas universales (descomposición, reconocimiento de patrones,
abstracción, diseño de algoritmos, depuración / evaluación) introducidas
progresivamente G6°–G11°. \\
\hline
CASEL (2020) & Habilidades Socioemocionales &
5 competencias SEL para la ciudadanía digital activa; adoptadas en la fila 6
de cada malla curricular. \\
\hline
UNESCO (2024) & Ética de la IA &
Marco de competencias de IA para estudiantes y docentes; mapeado a los ejes
PC (dimensiones técnicas) y SEL (dimensiones éticas) en Media. \\
\hline
MEN (2008, 2022) & Referente curricular &
Guía No. 30 (4 componentes) + Orientaciones Curriculares (PC como eje);
documento rector de la estructura de este plan. \\
\hline
\end{tabular}
\end{center}

% ============================================================
\section{Marco contextual}
% ============================================================

\subsection{Evolución normativa del área}

Mientras el Marco legal (sección\,2) reúne la normativa \textit{vigente}, esta sección
traza el \textbf{recorrido histórico} del área: cómo cada hito fue transformando
su identidad, de un enfoque instrumental al actual enfoque crítico y computacional.

\medskip
\newcommand{\hitonorm}[2]{%
  \noindent\colorbox{vinoInst}{\textcolor{white}{\textbf{~#1~}}}\hspace{0.5em}%
  \begin{minipage}[t]{0.85\linewidth}#2\end{minipage}\par\smallskip
}
\hitonorm{1994}{\textbf{Ley 115 — Ley General de Educación.} Tecnología e
Informática queda como área obligatoria y fundamental (arts. 23 y 31); la
formación en ciencia y técnica se reconoce como derecho.}
\hitonorm{1996}{\textbf{Lineamientos Curriculares MEN.} La tecnología pasa a ser
objeto de estudio independiente; se supera el enfoque de «uso de computadores» y
se incorpora su dimensión cultural y social.}
\hitonorm{2008}{\textbf{Guía No. 30 — Ser competente en tecnología.} Estructura
los 4 componentes (Naturaleza y evolución, Apropiación y uso, Solución de
problemas, Tecnología y sociedad) que articulan todas las mallas de este plan.}
\hitonorm{2022}{\textbf{Orientaciones Curriculares MEN.} Incorpora el Pensamiento
Computacional como estrategia transversal, junto a programación, robótica, maker,
STEM+ y narrativas transmedia: fundamento del eje PC de cada malla.}
\hitonorm{2024}{\textbf{Ley 2383 — Educación Socioemocional.} Formaliza la
enseñanza de competencias SEL en todos los grados: base de la fila 6 de cada malla.}
\hitonorm{2025}{\textbf{CONPES 4144 — Política Nacional de IA.} Lleva la formación
en IA ética y la alfabetización digital a la educación básica; orienta los
contenidos de 10°–11°.}

\subsection{Contexto socioeconómico y tecnológico del estudiantado}

La I.E. Sor María Juliana atiende principalmente a estudiantes de estratos 1, 2
y 3 del municipio de Cartago. La mayoría tiene acceso a internet móvil pero
limitado, y predominan dispositivos Android de gama media (360–414 px). Esta
realidad define el enfoque \textbf{offline-first} de la Plataforma Conéctate
y la prioridad de herramientas gratuitas, de bajo consumo de datos y sin
requerimiento de cuenta personal para menores de edad.

Esta realidad determina decisiones curriculares concretas: la \textbf{Plataforma
Conéctate} opera en modo \textit{offline-first} con Service Worker activo; las
guías evitan herramientas que requieran cuenta personal para menores de edad;
y los \textbf{anclajes ancestrales} de cada periodo parten de saberes locales
de Cartago, el Valle del Cauca y el Pacífico, no de contextos genéricos o
importados.

\subsection{Anclajes ancestrales del área (18 periodos)}

\textbf{¿Por qué abrir cada periodo con un saber ancestral?} El anclaje
ancestral no es un adorno cultural ni una anécdota de apertura: es una
\textbf{decisión epistemológica} del modelo MILC. Antes de presentar cualquier
herramienta digital, el área reconoce que el estudiante y su comunidad ya poseen
saberes tecnológicos propios —el oficio del relojero, el cálculo de la tejedora,
la lógica del pregonero, la memoria del campesino—. Partir de ese saber cumple
tres funciones curriculares:

\begin{itemize}
  \item \textbf{Puente cognitivo:} conecta lo nuevo (un concepto técnico) con
    algo que el estudiante ya entiende y respeta, reduciendo la barrera de entrada.
  \item \textbf{Valoración identitaria:} afirma que la tecnología del Valle y el
    Pacífico es tan legítima como la global, en línea con la filosofía de la
    liberación de Dussel; el estudiante es sujeto, no objeto, de la modernidad
    tecnológica.
  \item \textbf{Sentido ético:} enmarca cada aprendizaje técnico en una pregunta
    sobre \textit{para qué} y \textit{para quién} sirve la tecnología.
\end{itemize}

Así, el saber ancestral abre la fase de \textit{Escucha} del MILC y sostiene toda
la secuencia: el tema técnico del periodo se desarrolla como una \textbf{evolución}
de ese saber, no como su reemplazo. La siguiente tabla presenta los 18 anclajes
que abren los periodos del año 2026 (uno por grado y periodo):

\medskip
\setlength{\tabcolsep}{3pt}
\renewcommand{\arraystretch}{1.5}
\begin{tabularx}{\textwidth}{|X|X|X|}
\hline
\rowcolor{vinoInst}
\multicolumn{3}{|c|}{\color{white}\textbf{GRADO 6° — Comunicación e identidad digital}}\\
\hline
\cellcolor{p1bg}\textcolor{p1fg}{\small\textbf{P1 · Ene--Abr}}

El pregonero del barrio

\textit{\small Don Aurelio · Cartago}
&
\cellcolor{p2bg}\textcolor{p2fg}{\small\textbf{P2 · Abr--Ago}}

Relojero, costurera y panadero

\textit{\small Cartago}
&
\cellcolor{p3bg}\textcolor{p3fg}{\small\textbf{P3 · Ago--Nov}}

Doña Mercedes, maestra rural

\textit{\small La Plata, Cartago}
\\
\hline
\rowcolor{vinoInst}
\multicolumn{3}{|c|}{\color{white}\textbf{GRADO 7° — Producción multimedia y narrativa digital}}\\
\hline
\cellcolor{p1bg}\textcolor{p1fg}{\small\textbf{P1 · Ene--Abr}}

La minga del Valle

\textit{\small Valle del Cauca}
&
\cellcolor{p2bg}\textcolor{p2fg}{\small\textbf{P2 · Abr--Ago}}

Doña Sofía, la tejedora

\textit{\small Valle del Cauca}
&
\cellcolor{p3bg}\textcolor{p3fg}{\small\textbf{P3 · Ago--Nov}}

El consejero del barrio

\textit{\small Cartago}
\\
\hline
\rowcolor{vinoInst}
\multicolumn{3}{|c|}{\color{white}\textbf{GRADO 8° — Diseño de presentaciones y pensamiento visual}}\\
\hline
\cellcolor{p1bg}\textcolor{p1fg}{\small\textbf{P1 · Ene--Abr}}

El gesto de la cocina del Valle

\textit{\small Valle del Cauca}
&
\cellcolor{p2bg}\textcolor{p2fg}{\small\textbf{P2 · Abr--Ago}}

Decisiones del campesino

\textit{\small Campo caucano}
&
\cellcolor{p3bg}\textcolor{p3fg}{\small\textbf{P3 · Ago--Nov}}

Diseño visual del Pacífico

\textit{\small Chocó--Pacífico}
\\
\hline
\rowcolor{vinoInst}
\multicolumn{3}{|c|}{\color{white}\textbf{GRADO 9° — Diseño editorial y bases de datos}}\\
\hline
\cellcolor{p1bg}\textcolor{p1fg}{\small\textbf{P1 · Ene--Abr}}

La técnica en la mano

\textit{\small Artesanos del Valle}
&
\cellcolor{p2bg}\textcolor{p2fg}{\small\textbf{P2 · Abr--Ago}}

Pasquines, almanaques, esquelas

\textit{\small Cartago, siglo XX}
&
\cellcolor{p3bg}\textcolor{p3fg}{\small\textbf{P3 · Ago--Nov}}

La libreta de la abuela

\textit{\small Registros domésticos}
\\
\hline
\rowcolor{vinoInst}
\multicolumn{3}{|c|}{\color{white}\textbf{GRADO 10° — Automatización, IA y microemprendimiento}}\\
\hline
\cellcolor{p1bg}\textcolor{p1fg}{\small\textbf{P1 · Ene--Abr}}

El personaje silencioso del barrio

\textit{\small Cartago}
&
\cellcolor{p2bg}\textcolor{p2fg}{\small\textbf{P2 · Abr--Ago}}

Los 4 oficios de la palabra

\textit{\small Barrio Obrero, Cartago}
&
\cellcolor{p3bg}\textcolor{p3fg}{\small\textbf{P3 · Ago--Nov}}

Tiendas y panaderías del centro

\textit{\small Cartago / Quindío}
\\
\hline
\rowcolor{vinoInst}
\multicolumn{3}{|c|}{\color{white}\textbf{GRADO 11° — SEO, BPMN, identidad digital y legado}}\\
\hline
\cellcolor{p1bg}\textcolor{p1fg}{\small\textbf{P1 · Ene--Abr}}

Oficio y palabra

\textit{\small Saberes del barrio}
&
\cellcolor{p2bg}\textcolor{p2fg}{\small\textbf{P2 · Abr--Ago}}

El maestro carpintero

\textit{\small Artesano de Cartago}
&
\cellcolor{p3bg}\textcolor{p3fg}{\small\textbf{P3 · Ago--Nov}}

Oficios que escuchan al pueblo

\textit{\small Comunidad, Cartago}
\\
\hline
\end{tabularx}
\setlength{\tabcolsep}{6pt}
\renewcommand{\arraystretch}{1}

% ============================================================
\section{Marco conceptual del área}
% ============================================================

\subsection{Los 4 componentes MEN — Guía No. 30 / Orientaciones Curriculares 2022}

\begin{center}
\setlength{\tabcolsep}{4pt}
\begin{tabular}{|L{3.5cm}|L{5cm}|L{5cm}|}
\hline
\rowcolor{vinoInst}
\color{white}\textbf{Componente} &
\color{white}\textbf{Guía No. 30 (2008)} &
\color{white}\textbf{Orientaciones Curriculares MEN (2022)} \\
\hline
\textbf{Naturaleza y evolución de la tecnología} &
Historia, impacto y transformación social de la tecnología;
análisis crítico del cambio tecnológico. &
Agrega la dimensión de la IA, la robótica y las narrativas
transmedia como fenómenos tecnológicos actuales. \\
\hline
\textbf{Apropiación y uso de la tecnología} &
Manejo competente de herramientas digitales con criterio y
propósito; ofimática y producción. &
Enfatiza programación (Scratch, Code.org), robótica, maker y
creación con herramientas STEM+ como formas de apropiación
activa. \\
\hline
\textbf{Solución de problemas con tecnología} &
Diseño de soluciones tecnológicas a problemas reales;
proyectos integradores. &
Incorpora el Pensamiento Computacional como estrategia
transversal de resolución de problemas; depuración iterativa
y validación con datos reales. \\
\hline
\textbf{Tecnología y sociedad} &
Impacto ético, social y cultural de la tecnología;
ciudadanía digital, privacidad. &
Suma el pensamiento crítico ante la IA, la ética de
algoritmos, la brecha de género en STEM y la participación
ciudadana en entornos digitales. \\
\hline
\end{tabular}
\end{center}

\subsection{Énfasis de las Orientaciones Curriculares MEN (2022)}

Las Orientaciones Curriculares MEN 2022 para Tecnología e Informática no
sustituyen la Guía No. 30: la actualizan y amplían. Sus énfasis más relevantes
para este plan son:

\begin{center}
\setlength{\tabcolsep}{4pt}
\begin{tabular}{|L{4cm}|L{9.5cm}|}
\hline
\rowcolor{vinoInst}
\color{white}\textbf{Énfasis OC 2022} &
\color{white}\textbf{Implementación en la I.E. Sor María Juliana} \\
\hline
\textbf{Pensamiento Computacional (PC) como eje transversal} &
Las 5 habilidades Wing (2006) —descomposición, reconocimiento de patrones,
abstracción, diseño de algoritmos, depuración / evaluación— aparecen en la fila 5
(verde) de cada malla curricular, distribuidas progresivamente de G6° a G11°. \\
\hline
\textbf{Programación y creación digital} &
Scratch (G6°–G7°), Code.org y Python básico (G9°–G10°), automatización
con datos (G11°). Herramientas gratuitas sin registro de menores. \\
\hline
\textbf{Robótica y maker} &
Prototipos con materiales del entorno (cartón, madera, sensores Arduino)
como respuesta al contexto de recursos limitados del municipio. \\
\hline
\textbf{STEM+ y proyectos transversales} &
Articulación documentada con Ciencias Naturales (PRAE), Matemáticas
(estadística y algoritmos) y Ciudadanía (derechos digitales). \\
\hline
\textbf{Narrativas transmedia y game-based learning} &
Uso de plataformas de gamificación y narrativas multimodales en los
proyectos integradores de G7° y G8°. \\
\hline
\textbf{Inteligencia Artificial (IA) ética} &
G10°–G11°: declaración explícita de uso de IA, sesgo algorítmico,
privacidad y marco UNESCO 2024 de competencias en IA. \\
\hline
\end{tabular}
\end{center}

\subsection{Las 4 fases MILC en el diseño curricular}

La malla curricular de cada periodo organiza los 4 componentes MEN con las
4 fases del MILC y añade los ejes PC y SEL:

\begin{center}
\begin{tabular}{|L{3cm}|L{4cm}|L{3.5cm}|L{3.5cm}|}
\hline
\rowcolor{vinoInst}
\color{white}\textbf{Fase MILC} &
\color{white}\textbf{Componente MEN dominante} &
\color{white}\textbf{Habilidad PC} &
\color{white}\textbf{Competencia SEL} \\
\hline
\textbf{1. Escucha} &
Naturaleza y evolución &
Descomposición &
Autoconciencia \\
\hline
\textbf{2. Sistematización} &
Apropiación y uso &
Abstracción · Patrones &
Conciencia social \\
\hline
\textbf{3. Praxis} &
Solución de problemas &
Algoritmos · Depuración &
Autogestión · Habilidades relacionales \\
\hline
\textbf{4. Evaluación liberadora} &
Tecnología y sociedad &
Evaluación &
Toma de decisiones responsable \\
\hline
\end{tabular}
\end{center}

\subsection{Pensamiento Computacional como eje transversal}

\begin{center}
\begin{tabular}{|L{3cm}|L{4.5cm}|L{6cm}|}
\hline
\rowcolor{vinoInst}
\color{white}\textbf{Habilidad PC} &
\color{white}\textbf{Definición (Wing, 2006)} &
\color{white}\textbf{Cómo atraviesa todo el currículo} \\
\hline
Descomposición &
Dividir un problema en partes manejables. &
Presente en cada guía y proyecto: el estudiante divide toda tarea —de organizar
su cuaderno a diseñar un sistema— en partes, con exigencia creciente año a año. \\
\hline
Reconocimiento de patrones &
Identificar similitudes y regularidades. &
Se ejercita en todo trabajo digital: regularidades en textos, datos, interfaces,
lenguajes y procesos; la mirada se afina en cada periodo. \\
\hline
Abstracción &
Identificar lo esencial; ignorar detalles irrelevantes. &
Aparece siempre que se modela algo: representar lo esencial de un ícono, una red
o una base de datos, ocultando el detalle innecesario. \\
\hline
Diseño de algoritmos &
Crear soluciones paso a paso. &
En toda solución secuencial: instrucciones, recetas y automatizaciones; la
secuencia gana rigor lógico conforme avanza el grado. \\
\hline
Depuración / Evaluación &
Probar, identificar errores y mejorar. &
Cierra cada producto del área: probar, hallar el error y mejorar es el hábito que
acompaña todas las actividades, no una etapa final aislada. \\
\hline
\end{tabular}
\end{center}

\subsection{Habilidades Socioemocionales como eje transversal}

\begin{center}
\begin{tabular}{|L{3cm}|L{4.5cm}|L{6cm}|}
\hline
\rowcolor{vinoInst}
\color{white}\textbf{Competencia SEL} &
\color{white}\textbf{Definición (MEN / CASEL)} &
\color{white}\textbf{Anclaje en Tecnología e Informática} \\
\hline
Autoconciencia &
Reconocer emociones, fortalezas, valores e identidad propia. &
Identidad digital, huella en línea, autoría de contenidos. \\
\hline
Autogestión &
Regular emociones, establecer metas, perseverar. &
Gestión del tiempo en proyectos, depuración iterativa. \\
\hline
Conciencia social &
Ponerse en el lugar del otro; reconocer perspectivas diversas. &
Diseño accesible, privacidad del otro, impacto de la tecnología. \\
\hline
Habilidades relacionales &
Comunicarse, colaborar y resolver conflictos con empatía. &
Trabajo colaborativo en proyectos; comunicación profesional digital. \\
\hline
Toma de decisiones responsable &
Evaluar consecuencias éticas y tomar decisiones fundamentadas. &
Uso ético de IA, derechos de autor, declaración de autoría. \\
\hline
\end{tabular}
\end{center}

% ============================================================
\section{Política de uso de Inteligencia Artificial (IA)}
% ============================================================

La irrupción de la IA generativa (asistentes de texto, imagen y código) cambió
las reglas del aula. El área no la prohíbe ni la idealiza: la \textbf{integra con
criterio}. Esta política regula su uso por estudiantes y docentes, en
concordancia con el CONPES 4144 de 2025 (Política Nacional de IA), el marco
UNESCO de competencias en IA (2024) y el principio rector del MILC: la tecnología
se piensa y se cuestiona, no solo se usa.

\subsection{Principio rector}

La IA es una \textbf{herramienta de apoyo al pensamiento, nunca su reemplazo}.
Su uso es legítimo cuando \emph{amplía} la capacidad del estudiante de crear,
comprender y decidir; es ilegítimo cuando \emph{sustituye} el proceso de
aprendizaje que la actividad busca formar. El criterio no es la herramienta, sino
si el estudiante sigue siendo el autor de su pensamiento.

\subsection{Usos permitidos y no permitidos}

\begin{center}
\begin{tabular}{|L{6.8cm}|L{6.8cm}|}
\hline
\rowcolor{vinoInst}
\color{white}\textbf{Uso permitido (con declaración)} &
\color{white}\textbf{Uso no permitido} \\
\hline
Consultar dudas, pedir ejemplos o explicaciones alternativas de un concepto. &
Entregar como propio un texto, imagen o código generado por IA sin declararlo. \\
\hline
Generar borradores o lluvias de ideas que el estudiante luego revisa, corrige y
hace suyos. &
Usar IA para evadir el proceso de aprendizaje que la actividad busca formar
(p. ej. resolver el examen). \\
\hline
Apoyar la depuración de un trabajo (revisar ortografía, encontrar un error). &
Suplantar la voz, autoría o criterio propio del estudiante. \\
\hline
Producir recursos de práctica adicionales o accesibilidad (audio, resúmenes). &
Compartir datos personales, de terceros o sensibles con servicios de IA. \\
\hline
\end{tabular}
\end{center}

\subsection{Declaración obligatoria de uso}

Todo producto en el que se haya usado IA debe incluir una \textbf{declaración
breve}: qué herramienta se usó, para qué, y qué aportó el estudiante. Declarar el
uso \textbf{no penaliza}: lo que se evalúa es el criterio con que se usó. Ocultarlo,
en cambio, se considera falta de honestidad académica según el SIE institucional.

\subsection{Progresión por nivel}

\begin{center}
\begin{tabular}{|L{3cm}|L{11cm}|}
\hline
\rowcolor{vinoInst}
\color{white}\textbf{Nivel} & \color{white}\textbf{Uso esperado de la IA} \\
\hline
\textbf{6° -- 8°} &
Uso \textbf{guiado y supervisado}: la IA se usa en clase, con el docente, para
comprender cómo funciona, reconocer sus errores y cuidar la privacidad. Énfasis
en alfabetización: qué es, qué no es y por qué se equivoca. \\
\hline
\textbf{9° -- 11°} &
Uso \textbf{autónomo con criterio}: el estudiante decide cuándo la IA aporta,
declara su uso, evalúa el sesgo y la veracidad de lo generado, y conserva la
autoría de su producto, especialmente en los proyectos integradores. \\
\hline
\end{tabular}
\end{center}

\subsection{Privacidad, sesgo y rol docente}

El área enseña que los modelos de IA \textbf{pueden equivocarse y tener sesgos}, y
que requieren verificación humana. No se comparten datos de menores ni de terceros
con servicios externos. El docente modela el uso ético, verifica las
declaraciones de uso y orienta la reflexión sobre el impacto social de la IA
—componente \textit{Tecnología y sociedad} de la Guía No. 30—.

% ============================================================
\section{Propósitos y metas del área 2026}
% ============================================================

\noindent Definidos el marco, el contexto y los conceptos del área, este apartado
fija el \textbf{horizonte}: qué se busca formar y cómo se sabrá si se logró.

\subsection{Propósito general}

Formar ciudadanos \textbf{tecnológicamente letrados, computacionalmente
pensantes, socioemocional\-mente competentes y éticamente responsables},
capaces de comprender, usar y transformar la tecnología al servicio de
su comunidad y de sus proyectos de vida, desde el territorio de Cartago
y el Valle del Cauca.

\subsection{Metas 2026}

Cada meta se enuncia con su indicador y el instrumento que la verifica, para que
sea evaluable y no quede como mera declaración de intención:

\begin{center}
\footnotesize
\setlength{\tabcolsep}{4pt}
\renewcommand{\arraystretch}{1.25}
\begin{tabular}{|L{5.6cm}|L{1.7cm}|L{5.8cm}|}
\hline
\rowcolor{vinoInst}
\color{white}\textbf{Meta 2026} &
\color{white}\textbf{Indicador} &
\color{white}\textbf{Instrumento de verificación} \\
\hline
Cobertura: guías y proyectos con PC y SEL explícitos & 100\% & Revisión editorial + linter MILC v3 \\
\hline
Aprendizaje en PC (G8°–G11°): demuestra las 5 habilidades & 80\% & Rúbrica del proyecto integrador (sección\,15.6) \\
\hline
Aprendizaje en SEL: nivel \textit{En proceso} o superior en ≥3 competencias & 85\% & Rúbrica de habilidades socioemocionales (sección\,15.7) \\
\hline
IA ética (G10°–G11°): uso declarado y con criterio & 100\% & Declaración de uso + criterio IA de la rúbrica del proyecto \\
\hline
Participación en Bebras u olimpiada de PC (G8°–G10°) & ≥50\% & Registro de inscripción del concurso \\
\hline
Desempeño general Básico o superior & ≥90\% & Escala de valoración institucional (boletín) \\
\hline
Guías publicadas antes de cada periodo, accesibles offline & 100\% & Estado de publicación en la Plataforma Conéctate \\
\hline
Articulación con otras áreas en proyectos documentados & ≥3 áreas & Evidencias de proyectos transversales (PRAE, ciudadanía, emprendimiento) \\
\hline
Calidad didáctica: guías sin errores de estructura & 100\% & Validación \texttt{make guia-lint} \\
\hline
\end{tabular}
\end{center}

% ============================================================
\section{Diseño curricular — Malla curricular (G6°--G11°)}
% ============================================================

\noindent Hasta aquí, el \emph{qué} y el \emph{para qué} del área. La malla
curricular es el \emph{cómo} se organiza todo ello, grado a grado y periodo a
periodo.

\medskip
\noindent Cada periodo de cada grado se organiza en una tabla con \textbf{6 filas
de componentes}:

\begin{enumerate}
  \item Naturaleza y evolución de la tecnología
  \item Apropiación y uso de la tecnología
  \item Solución de problemas con tecnología
  \item Tecnología y sociedad
  \item \textcolor{verdePC}{\textbf{Pensamiento Computacional}} (fondo verde)
  \item \textcolor{moradoSEL}{\textbf{Habilidades Socioemocionales}} (fondo violeta)
\end{enumerate}

Las \textbf{18 mallas completas} (G6° a G11°, Periodos 1 a 3), en formato
apaisado, se presentan en el \textbf{Anexo A} al final del documento. Cada
periodo incluye su tabla con las 6 filas de componentes y los descriptores
cognitivo, procedimental y actitudinal.


% ============================================================
\section{Metodología}
% ============================================================

El área de Tecnología e Informática opera bajo el \textbf{Modelo de Investigación
Liberadora y Científica (MILC v3)}, especificidad del modelo pedagógico institucional
interactivo y crítico de la I.E. Sor María Juliana. El MILC no es un marco externo
adoptado: es la pedagogía propia del área, construida desde las realidades del
territorio de Cartago y el Valle del Cauca.

Sus \textbf{cuatro fases} —Escucha, Sistematización, Praxis y Evaluación
liberadora, con su matriz completa de estudiante, docente y herramienta en el
Marco teórico (sección\,5.1)— no se quedan en el papel. Las secciones siguientes muestran
cómo el Pensamiento Computacional las potencia, cómo se estructura cada guía y
cómo transcurre una sesión de clase.

\subsection{Pensamiento Computacional como gramática cognitiva transversal}

El Pensamiento Computacional (Wing, 2006; MEN OC 2022) atraviesa las 4 fases del
MILC como una gramática cognitiva que potencia cada una:

\begin{itemize}
  \item \textbf{En la Escucha:} el estudiante descompone el problema o la situación
    ancestral en partes comprensibles.
  \item \textbf{En la Sistematización:} abstrae patrones y reconoce estructuras
    generalizables desde los casos concretos.
  \item \textbf{En la Praxis:} diseña algoritmos de solución y los depura con
    criterio iterativo.
  \item \textbf{En la Evaluación:} evalúa la solución con criterios explícitos y
    propone mejoras.
\end{itemize}

\subsection{Modelo pedagógico de la guía de aprendizaje}

Cada guía de la Plataforma Conéctate sigue la estructura de 10 estaciones MILC:
(1) Saber ancestral, (2) Tu ruta hoy, (3) Propósito de aprendizaje, (4) Marco
conceptual, (5) Secuencia de actividades, (6) Componente práctico tecnológico,
(7) Evaluación, (8) Cierre filosófico Triángulo Dussel/Marco Aurelio/Floridi,
(9) Evaluación liberadora — 5 dimensiones, (10) Pie institucional.

\callout{Regla de 200 palabras.}{ningún bloque de texto supera 200 palabras sin
una acción intercalada. Es la regla más poderosa para mantener al estudiante
activo durante toda la guía.}

\subsection{Secuencia de una sesión de aula}

Las 4 fases del MILC se traducen en el siguiente ritmo de clase. Los tiempos son
orientativos para un bloque presencial; en \textbf{Media} (1 h/semana) las fases de
lectura y práctica se adelantan en la Plataforma Conéctate y la clase se concentra
en socializar, resolver dudas y cerrar.

\begin{center}
\setlength{\tabcolsep}{4pt}
\begin{tabular}{|L{2.1cm}|L{3.2cm}|L{6cm}|L{1.9cm}|}
\hline
\rowcolor{vinoInst}
\color{white}\textbf{Momento} &
\color{white}\textbf{Fase MILC} &
\color{white}\textbf{Qué ocurre en el aula} &
\color{white}\textbf{Tiempo} \\
\hline
Apertura & 1. Escucha & El docente activa el saber ancestral y lanza la pregunta-puente; el estudiante conecta con lo que ya sabe. & 10–15 min \\
\hline
Desarrollo & 2. Sistematización & Lectura del marco con bloques LEE · OBSERVA · EXPLORA; el docente modela con ejemplos de Cartago. & 25–35 min \\
\hline
Producción & 3. Praxis & El estudiante realiza la ACTIVIDAD (uno de los 6 verbos) en el cuaderno o el computador; el docente acompaña. & 30–50 min \\
\hline
Verificación & Formativa & Bloque VERIFICA o mini-quiz mental; corrección entre pares. & 5–10 min \\
\hline
Cierre & 4. Evaluación liberadora & Reflexión 5D en el cuaderno y, en 9°–11°, cierre filosófico con el triángulo de pensamiento. & 10–15 min \\
\hline
\end{tabular}
\end{center}

% ============================================================
\section{Recursos educativos}
% ============================================================

\subsection{Recursos digitales}

\begin{center}
\begin{tabular}{|L{3cm}|L{5cm}|L{6cm}|}
\hline
\rowcolor{vinoInst}
\color{white}\textbf{Tipo} &
\color{white}\textbf{Recurso} &
\color{white}\textbf{Uso principal} \\
\hline
Plataforma propia &
Plataforma Conéctate &
180 guías interactivas + 18 proyectos, acceso offline \\
\hline
Suite IA &
Claude.ai, ChatGPT, Gemini &
Asistente de escritura, prompting, análisis \\
\hline
Diseño gráfico &
Canva, Adobe Express, Figma (básico) &
Piezas editoriales y visuales \\
\hline
Ofimática &
Google Workspace (Docs, Sheets, Slides, Forms) &
Documentos, datos y presentaciones \\
\hline
Automatización &
Zapier, Make (n8n básico) &
Automatización de procesos (G11°) \\
\hline
Maquetación &
Scribus, DTP en línea &
Diseño editorial (G9°) \\
\hline
Video y podcast &
Audacity, CapCut, DaVinci Resolve básico &
Producción multimedia (G7°) \\
\hline
\end{tabular}
\end{center}

\subsection{Recursos impresos y físicos}

\begin{itemize}
  \item Guías en PDF descargables (generadas desde YAML con pipeline \texttt{make guia-build}).
  \item Cuaderno del estudiante (único repositorio físico de la producción escrita).
  \item Sala de informática de la institución (mínimo 20 equipos operativos).
  \item Conexión a internet compartida (los recursos funcionan offline-first).
\end{itemize}

\subsection{Criterios de selección de herramientas}

\begin{enumerate}
  \item \textbf{Acceso gratuito o freemium} para estudiantes sin costo adicional.
  \item \textbf{Funcionamiento offline o con baja conectividad} como primera opción.
  \item \textbf{Sin requerimiento de cuenta personal} para menores de edad (cuando sea posible).
  \item \textbf{Herramientas con propósito educativo claro,} no entretenimiento disfrazado.
\end{enumerate}

\subsection{Herramienta, grado y acceso de menores}

Dado el contexto (estudiantes menores, datos limitados), cada herramienta se elige
según el grado, si exige cuenta personal y qué alternativa offline o sin registro
existe:

\begin{center}
\footnotesize
\setlength{\tabcolsep}{4pt}
\begin{tabular}{|L{3.4cm}|L{1.8cm}|L{3.2cm}|L{4cm}|}
\hline
\rowcolor{vinoInst}
\color{white}\textbf{Herramienta} &
\color{white}\textbf{Grado(s)} &
\color{white}\textbf{¿Cuenta de menor?} &
\color{white}\textbf{Alternativa offline / sin cuenta} \\
\hline
Plataforma Conéctate & 6°–11° & No & Es la base offline del área \\
\hline
Scratch & 6°–8° & No (modo sin registro) & Scratch Desktop (local) \\
\hline
Google Workspace & 8°–11° & Cuenta institucional & LibreOffice (offline) \\
\hline
Canva / Adobe Express & 7°–11° & Gestionada por el docente & GIMP · Scribus (offline) \\
\hline
Claude.ai / IA generativa & 9°–11° & Uso guiado por el docente & Solo con acompañamiento docente \\
\hline
CapCut / DaVinci & 7° & No (DaVinci local) & DaVinci Resolve (offline) \\
\hline
\end{tabular}
\end{center}

% ============================================================
\section{Plataforma Conéctate}
% ============================================================

\subsection{Descripción general}

La \textbf{Plataforma Conéctate} (alguarito.github.io/plataformaconectate) es el
repositorio digital de todos los recursos del área de Tecnología e Informática
de la I.E. Sor María Juliana. Construida con \textbf{Astro + Tailwind + TypeScript},
opera como \textbf{Progressive Web App (PWA) offline-first} mediante Service Worker,
lo que permite su uso en dispositivos con conectividad intermitente.

\subsection{Alcance actual}

\begin{center}
\begin{tabular}{|L{7cm}|L{3cm}|}
\hline
\rowcolor{vinoInst}
\color{white}\textbf{Contenido} & \color{white}\textbf{Cantidad} \\
\hline
Guías de aprendizaje (G6°–G11°) & 180 \\
\hline
Proyectos integradores & 18 \\
\hline
Exámenes de periodo & 18 \\
\hline
Grados cubiertos & 6° a 11° \\
\hline
Periodos por grado & 3 \\
\hline
\end{tabular}
\end{center}

\subsection{Malla de sesiones — relación de guías por grado y periodo}

La estructura curricular del área se organiza en \textbf{10 sesiones de aprendizaje
(una guía por sesión) en cada periodo}. La siguiente malla muestra la relación
completa de guías publicadas en la Plataforma Conéctate, de sexto a undécimo:

\begin{center}
\setlength{\tabcolsep}{6pt}
\renewcommand{\arraystretch}{1.35}
\begin{tabular}{|L{3.5cm}|M{2.3cm}|M{2.3cm}|M{2.3cm}|M{2cm}|}
\hline
\rowcolor{vinoInst}
\color{white}\textbf{Grado} &
\color{white}\textbf{Periodo 1} &
\color{white}\textbf{Periodo 2} &
\color{white}\textbf{Periodo 3} &
\color{white}\textbf{Total} \\
\hline
Sexto (6°) & 10 & 10 & 10 & \textbf{30} \\
\hline
Séptimo (7°) & 10 & 10 & 10 & \textbf{30} \\
\hline
Octavo (8°) & 10 & 10 & 10 & \textbf{30} \\
\hline
Noveno (9°) & 10 & 10 & 10 & \textbf{30} \\
\hline
Décimo (10°) & 10 & 10 & 10 & \textbf{30} \\
\hline
Undécimo (11°) & 10 & 10 & 10 & \textbf{30} \\
\hline
\rowcolor{vinoClaro}
\textbf{Total por periodo} & \textbf{60} & \textbf{60} & \textbf{60} & \textbf{180} \\
\hline
\end{tabular}
\end{center}

\noindent\textbf{Síntesis del diseño:} 6 grados $\times$ 3 periodos $\times$ 10
sesiones $=$ \textbf{180 guías de aprendizaje}, complementadas con \textbf{18
proyectos integradores} (uno por grado y periodo) y \textbf{18 exámenes de periodo},
para un total de \textbf{216 piezas curriculares} activas en la plataforma. Cada
sesión corresponde a una guía con su estructura MILC de 10 estaciones; cada periodo
cierra con un proyecto integrador y su examen.

\subsection{Ruta de las guías — títulos de sesión por grado y periodo}

Cada sesión tiene un título que la ubica en la ruta de aprendizaje del grado. El
\textbf{temario completo de las 180 guías} —en el orden en que el estudiante las
recorre a lo largo del año— se presenta en el \textbf{Anexo B}. Cada título
corresponde a una guía navegable en línea y descargable en PDF.

\subsection{Pipeline de generación de guías}

Las guías se generan desde una fuente única (YAML) mediante un pipeline automatizado:

\begin{itemize}
  \item \texttt{content/guias/\{grado\}/\{clave\}.yaml} $\rightarrow$ fuente única editable
  \item $\rightarrow$ PDF (xelatex) $\rightarrow$ \texttt{public/guias-mejoras/}
  \item $\rightarrow$ TypeScript $\rightarrow$ \texttt{src/data/guiasContenido/}
\end{itemize}

\callout{Regla de oro.}{nunca editar los archivos generados (\texttt{.tex},
\texttt{.pdf}, \texttt{.ts}); siempre editar el YAML.}

% ============================================================
\section{Intensidad horaria}
% ============================================================

\begin{center}
\begin{tabular}{|L{3.5cm}|L{2cm}|L{3cm}|L{4cm}|}
\hline
\rowcolor{vinoInst}
\color{white}\textbf{Nivel} &
\color{white}\textbf{Grado} &
\color{white}\textbf{Horas semanales} &
\color{white}\textbf{Horas anuales (aprox.)} \\
\hline
Básica Secundaria & 6° & 3 & 120 \\
\hline
Básica Secundaria & 7° & 3 & 120 \\
\hline
Básica Secundaria & 8° & 3 & 120 \\
\hline
Básica Secundaria & 9° & 3 & 120 \\
\hline
Media & 10° & 1 & 40 \\
\hline
Media & 11° & 1 & 40 \\
\hline
\end{tabular}
\end{center}

\textbf{Nota:} las horas anuales se estiman sobre 40 semanas lectivas (Decreto 1850
de 2002). Cada periodo desarrolla aproximadamente 10 sesiones (guías). En
\textbf{Básica Secundaria} (3 h/semana) las sesiones se trabajan principalmente en
el aula; en \textbf{Media} (1 h/semana) la menor intensidad presencial se
complementa con el \textbf{trabajo autónomo en la Plataforma Conéctate}
(\textit{offline-first}): el estudiante avanza las guías por su cuenta y el docente
acompaña, resuelve dudas y evalúa los productos en la hora semanal.

% ============================================================
\section{Sistema de evaluación}
% ============================================================

\subsection{Marco legal y principio rector}

La evaluación en el área sigue el \textbf{Decreto 1290 de 2009} y el
\textbf{Sistema Institucional de Evaluación (SIE)} de la I.E. Sor María Juliana.
El principio rector es la \textbf{evaluación liberadora}: el estudiante no es
objeto de medición sino sujeto que se reconoce, reflexiona y transforma.

\textbf{¿Cómo conviven la evaluación liberadora y la calificación numérica?} El
área no las opone, las articula. La \textbf{calificación} (escala 1.0–5.0 con
pesos 30/40/30) certifica el desempeño ante el SIE y es obligatoria por el
Decreto 1290. La \textbf{evaluación liberadora} es la capa formativa que da
sentido a esa nota: no produce una cifra adicional, sino que (1) alimenta el
componente \textit{actitudinal} (30\%) con evidencia del cuaderno, y (2) es
\textbf{requisito de cierre} de cada guía —sin la reflexión de las 5 dimensiones,
la guía no se considera terminada—. Así el estudiante recibe una nota que cumple
la norma y, a la vez, un espacio donde es sujeto y no objeto de la evaluación.

\subsection{Pesos por componente}

\begin{center}
\begin{tabular}{|L{8cm}|L{3cm}|}
\hline
\rowcolor{vinoInst}
\color{white}\textbf{Componente} & \color{white}\textbf{Peso} \\
\hline
Desempeño cognitivo (saber conceptual) & 30\% \\
\hline
Desempeño procedimental (saber hacer) & 40\% \\
\hline
Desempeño actitudinal (saber ser) & 30\% \\
\hline
\end{tabular}
\end{center}

\subsection{Tipos de evaluación}

\begin{center}
\begin{tabular}{|L{2.5cm}|L{3.5cm}|L{7cm}|}
\hline
\rowcolor{vinoInst}
\color{white}\textbf{Tipo} &
\color{white}\textbf{Cuándo} &
\color{white}\textbf{Instrumento} \\
\hline
Diagnóstica &
Inicio de cada periodo &
Conversación + pregunta-puente del saber ancestral \\
\hline
Formativa &
Durante las guías &
Bloque VERIFICA · Rúbricas de proceso · Auto y coevaluación \\
\hline
Sumativa &
Final de cada sesión / periodo &
Proyecto integrador + Examen de periodo \\
\hline
Liberadora &
Cierre de guía &
5 dimensiones (personal, emocional, ciudadana, local, intergeneracional) \\
\hline
\end{tabular}
\end{center}

\subsection{Escala de valoración institucional}

\begin{center}
\begin{tabular}{|L{4cm}|L{3cm}|}
\hline
\rowcolor{vinoInst}
\color{white}\textbf{Nivel} & \color{white}\textbf{Rango} \\
\hline
Superior & 4.6 – 5.0 \\
\hline
Alto & 4.0 – 4.5 \\
\hline
Básico & 3.0 – 3.9 \\
\hline
Bajo & 1.0 – 2.9 \\
\hline
\end{tabular}
\end{center}

\subsection{Evaluación liberadora — las 5 dimensiones}

La evaluación liberadora cierra cada guía con cinco preguntas que el estudiante
responde en su cuaderno. No buscan una respuesta correcta, sino que el estudiante
se reconozca como sujeto del aprendizaje. El docente las lee como evidencia del
componente actitudinal y como termómetro del sentido que el estudiante construyó:

\begin{center}
\begin{tabular}{|L{3.2cm}|L{10.8cm}|}
\hline
\rowcolor{vinoInst}
\color{white}\textbf{Dimensión} & \color{white}\textbf{Pregunta-guía de cierre (modelo)} \\
\hline
\textbf{Personal} & ¿Qué aprendí hoy que no sabía y para qué me sirve a mí? \\
\hline
\textbf{Emocional} & ¿Cómo me sentí durante la actividad? ¿Qué me costó y cómo lo enfrenté? \\
\hline
\textbf{Ciudadana} & ¿Cómo puede esto que aprendí ayudar —o dañar— a otras personas? \\
\hline
\textbf{Local} & ¿Dónde veo esto en mi barrio, mi familia o mi comunidad de Cartago? \\
\hline
\textbf{Intergeneracional} & ¿Qué le explicaría de esto a alguien mayor? ¿Y qué me podría enseñar esa persona a cambio? \\
\hline
\end{tabular}
\end{center}

\noindent\textbf{Adaptación por ciclo:} en 6°–8° se responden en 1–2 renglones,
en lenguaje cotidiano; en 9°–11° se exige argumentación y conexión con el
triángulo de pensamiento (cierre filosófico del Marco teórico, sección\,5).
\textbf{Registro:} van al cuaderno bajo el rótulo «Reflexión 5D».
\textbf{Peso:} no producen nota propia; alimentan el componente actitudinal
(30\%) y son requisito de cierre de la guía.

\subsection{Rúbrica del proyecto integrador}

Cada periodo cierra con un \textbf{proyecto integrador}: una solución tecnológica
a un problema real de la comunidad, sustentada ante los pares. Se evalúa con esta
rúbrica única —reutilizable en los 18 proyectos— sobre la escala institucional:

\begin{center}
\scriptsize
\setlength{\tabcolsep}{3pt}
\begin{tabular}{|L{2.7cm}|L{2.7cm}|L{2.7cm}|L{2.7cm}|L{2.4cm}|}
\hline
\rowcolor{vinoInst}
\color{white}\textbf{Criterio} &
\color{white}\textbf{Superior} &
\color{white}\textbf{Alto} &
\color{white}\textbf{Básico} &
\color{white}\textbf{Bajo} \\
\hline
\textbf{Producto (40\%)} & Resuelve el problema y supera lo pedido. & Resuelve el problema completo. & Resuelve parcialmente. & No funciona o no resuelve. \\
\hline
\textbf{Pensamiento Computacional (20\%)} & Descompone, abstrae y depura con autonomía. & Aplica las 5 habilidades con apoyo. & Aplica algunas habilidades. & No evidencia proceso PC. \\
\hline
\textbf{Anclaje y pertinencia (15\%)} & Conecta de forma original con el saber local. & Conecta con el contexto. & Conexión débil o forzada. & Sin conexión con el contexto. \\
\hline
\textbf{Uso ético de IA (10\%)} & Declara y justifica el uso con criterio. & Declara el uso de IA. & Declaración incompleta. & Usa IA sin declarar. \\
\hline
\textbf{Sustentación (15\%)} & Explica con claridad y responde preguntas. & Explica el proceso. & Explica con dificultad. & No sustenta. \\
\hline
\end{tabular}
\end{center}

\subsection{Rúbrica de Habilidades Socioemocionales (SEL)}

Las competencias SEL se observan a lo largo del periodo (no en una prueba) y se
registran con esta rúbrica de tres niveles. Sostiene la meta institucional de que
el 85\% de los estudiantes alcance nivel \textit{En proceso} o superior en al
menos tres competencias:

\begin{center}
\scriptsize
\setlength{\tabcolsep}{3pt}
\begin{tabular}{|L{3.1cm}|L{3.5cm}|L{3.5cm}|L{3.4cm}|}
\hline
\rowcolor{vinoInst}
\color{white}\textbf{Competencia} &
\color{white}\textbf{Inicial} &
\color{white}\textbf{En proceso} &
\color{white}\textbf{Consolidado} \\
\hline
Autoconciencia & Reconoce con ayuda sus emociones y autoría. & Identifica sus emociones e identidad digital. & Regula su identidad digital con criterio propio. \\
\hline
Autogestión & Necesita supervisión constante. & Gestiona su tiempo con recordatorios. & Planifica y persevera de forma autónoma. \\
\hline
Conciencia social & Le cuesta reconocer al otro. & Reconoce el impacto en otros cuando se le señala. & Anticipa el impacto de su tecnología en otros. \\
\hline
Habilidades relacionales & Colabora con dificultad. & Colabora cuando se le asigna un rol. & Lidera y resuelve conflictos en el equipo. \\
\hline
Decisiones responsables & Decide sin ver consecuencias. & Considera consecuencias al decidir. & Fundamenta éticamente sus decisiones (IA, autoría). \\
\hline
\end{tabular}
\end{center}

\subsection{Evaluación del Pensamiento Computacional}

El PC se evalúa de forma \textbf{integrada} dentro de cada proyecto integrador
y guía, no como prueba aislada. Los indicadores son:

\begin{enumerate}
  \item ¿El estudiante descompone el problema antes de intentar resolverlo?
  \item ¿Identifica y aplica patrones de solución conocidos?
  \item ¿Abstrae lo esencial y descarta lo irrelevante?
  \item ¿Diseña una solución paso a paso verificable?
  \item ¿Depura su propuesta con base en criterios explícitos?
\end{enumerate}

\subsection{Ítems tipo ICFES}

Cada evaluación de periodo incluye mínimo \textbf{3 ítems de selección múltiple
con única respuesta}, con estructura ICFES: enunciado contextualizado (situación-problema
del entorno del estudiante), 4 opciones (A, B, C, D) y una sola respuesta correcta
con justificación en la clave del docente.

\noindent\textbf{Ítem modelo (9°, componente Tecnología y sociedad):}

\begin{center}
\begin{tabular}{|L{13.5cm}|}
\hline
\rowcolor{vinoClaro}
\textbf{Situación.} Daniela, estudiante de Cartago, publica a diario una foto con
la ubicación activada de la cancha donde entrena, siempre a la misma hora. Un
desconocido comienza a aparecer en sus entrenamientos. \\
\hline
\textbf{Pregunta.} ¿Cuál es la causa \textit{tecnológica} más directa de este riesgo? \\
A. La cámara del celular tiene demasiada resolución. \newline
B. \textbf{Los metadatos de geolocalización revelaron un patrón de lugar y hora.} \newline
C. La foto se publicó en una red social gratuita. \newline
D. El desconocido tiene conocimientos avanzados de informática. \\
\hline
\rowcolor{verdePCclaro}
\textbf{Clave: B.} El riesgo no nace de la calidad de la cámara (A), ni de que la
red sea gratuita (C), ni requiere un experto (D): nace de que los \textit{metadatos}
de ubicación, combinados con la rutina, exponen un patrón predecible. El ítem
evalúa el concepto de huella digital y privacidad, no la anécdota. \\
\hline
\end{tabular}
\end{center}

% ============================================================
\section{Atención a la diversidad y NEE}
% ============================================================

\noindent Las secciones que siguen responden a una convicción: un plan solo es
bueno si \textbf{nadie queda afuera}. Diversidad, recuperación, articulación y
comunicación con las familias son las garantías de equidad del área.

\subsection{Principio de Diseño Universal para el Aprendizaje (DUA)}

El área adopta los principios del \textbf{DUA} como marco de diversificación:

\begin{itemize}
  \item \textbf{Múltiples medios de representación:} texto, imagen, audio, video, interacción digital.
  \item \textbf{Múltiples medios de expresión:} escritura en cuaderno, producto digital, exposición oral.
  \item \textbf{Múltiples medios de compromiso:} elección de tema de proyecto dentro del eje curricular.
\end{itemize}

\subsection{Adaptaciones específicas}

\begin{center}
\begin{tabular}{|L{4cm}|L{10cm}|}
\hline
\rowcolor{vinoInst}
\color{white}\textbf{Necesidad} & \color{white}\textbf{Adaptación} \\
\hline
Dificultad de lectura &
Versiones de guía con tipografía aumentada y audio-descripción. \\
\hline
Déficit atencional &
Actividades de ≤15 min con cierre claro; alternancia leer/hacer (regla 200 palabras). \\
\hline
Movilidad reducida &
Actividades adaptadas sin requerimiento de escritura en tablero; uso de voz. \\
\hline
Superdotación &
Retos de ampliación (nivel CREA avanzado) y proyectos de investigación autónoma. \\
\hline
Acceso limitado a dispositivos &
PDF descargable + actividades adaptables a papel. \\
\hline
\end{tabular}
\end{center}

\subsection{Estudiantes con PIAR}

Para estudiantes con Plan Individual de Ajustes Razonables (PIAR), el docente:
(1) identifica el ajuste en coordinación con el equipo de apoyo,
(2) adapta los descriptores de desempeño al nivel del estudiante,
(3) mantiene el marco del MILC y el espíritu de la evaluación liberadora, y
(4) registra las adaptaciones en el PIAR sin cambiar la identidad del área.

\subsection{Ejemplo de adaptación en una actividad técnica}

La adaptación no baja la meta de aprendizaje: cambia la vía o el nivel de
exigencia, conservando el concepto técnico. Ejemplo sobre una misma actividad
—\textit{«crear una animación en Scratch» (G7°)}—:

\begin{center}
\footnotesize
\setlength{\tabcolsep}{4pt}
\begin{tabular}{|L{3.6cm}|L{5cm}|L{5cm}|}
\hline
\rowcolor{vinoInst}
\color{white}\textbf{Estudiante} &
\color{white}\textbf{Descriptor original} &
\color{white}\textbf{Descriptor adaptado (misma lógica)} \\
\hline
Dificultad motriz fina & Programa 5 bloques arrastrándolos con el mouse. & Programa 3 bloques con teclado o por dictado al docente; conserva la secuencia lógica. \\
\hline
TDAH / déficit atencional & Completa la animación en una sesión. & Completa la animación en 3 micro-metas de 10 min, cada una verificada. \\
\hline
Talento excepcional & Anima un personaje. & Anima una escena con 2 personajes e interacción condicional (\textit{si–entonces}). \\
\hline
\end{tabular}
\end{center}

% ============================================================
\section{Actividades de recuperación}
% ============================================================

\subsection{Recuperación formativa}

Los estudiantes con desempeño Bajo tienen derecho a \textbf{procesos de recuperación}
durante el periodo y al final de éste:

\begin{enumerate}
  \item \textbf{Acompañamiento en guía:} repetir la guía con orientación personalizada del docente.
  \item \textbf{Tutoría entre pares:} el estudiante con desempeño Superior orienta al compañero.
  \item \textbf{Proyecto de recuperación:} producto alternativo que demuestre el desempeño no alcanzado.
\end{enumerate}

\subsection{Criterios operativos de la recuperación}

\begin{center}
\begin{tabular}{|L{3.5cm}|L{10.5cm}|}
\hline
\rowcolor{vinoInst}
\color{white}\textbf{Criterio} & \color{white}\textbf{Definición} \\
\hline
\textbf{Disparador} & Desempeño Bajo (inferior a 3.0) en una o más guías o en el periodo. \\
\hline
\textbf{Momento} & Continua durante el periodo (formativa) y una semana de recuperación al cierre de cada periodo. \\
\hline
\textbf{Evidencia} & El producto pendiente (guía del cuaderno o proyecto), no un examen nuevo. La recuperación reproduce la praxis, no la sustituye por memorización. \\
\hline
\textbf{Nota máxima} & La recuperación de cierre alcanza hasta nivel \textbf{Básico} (3.0–3.9), conforme al SIE; la mejora formativa durante el periodo puede alcanzar la nota plena. \\
\hline
\textbf{Registro} & Se documenta en la planilla del docente con fecha, evidencia y desempeño superado. \\
\hline
\end{tabular}
\end{center}

\subsection{Criterios de promoción}

Sigue el SIE institucional. Un estudiante reprueba el área si al finalizar el año su
desempeño es Bajo (promedio inferior a 3.0) y no supera el proceso de recuperación final.

% ============================================================
\section{Articulación con otras áreas}
% ============================================================

\begin{center}
\begin{tabular}{|L{4cm}|L{10cm}|}
\hline
\rowcolor{vinoInst}
\color{white}\textbf{Área} &
\color{white}\textbf{Punto de articulación} \\
\hline
\textbf{Matemáticas} &
Estadística y datos (G9° P3), algoritmos como secuencia lógica. \\
\hline
\textbf{Lengua Castellana} &
Producción textual (informes, revistas), argumentación digital. \\
\hline
\textbf{Ciencias Naturales} &
Modelado de fenómenos, sensores, bases de datos ambientales (PRAE). \\
\hline
\textbf{Ciencias Sociales} &
Impacto tecnológico, ciudadanía digital, brecha tecnológica. \\
\hline
\textbf{Ética y Valores} &
Ética de la IA, derechos de autor, privacidad digital. \\
\hline
\textbf{Emprendimiento} &
Microemprendimiento digital (G10°–G11°), BMC, plan de negocio. \\
\hline
\textbf{Educación Artística} &
Diseño editorial y visual (G9° P2), identidad digital. \\
\hline
\end{tabular}
\end{center}

\subsection{Proyectos transversales obligatorios}

\begin{itemize}
  \item \textbf{PRAE} (Proyecto Ambiental Escolar): uso de datos y visualización
    para monitoreo ambiental (G9° P3).
  \item \textbf{Ciudadanía digital}: convivencia en entornos digitales y uso
    responsable de redes (G7° P3).
  \item \textbf{Educación sexual}: privacidad, consentimiento digital y relaciones
    en línea (G8° P1, transversal).
  \item \textbf{Emprendimiento} (Ley 1014): microemprendimiento con IA
    (G10° P3, G11° P3).
\end{itemize}

\subsection{Formato de articulación (ejemplo: PRAE, G9° P3)}

Para que la articulación no quede en la intención, cada proyecto transversal se
documenta con un formato que precisa el aporte de cada área, su evidencia y el
momento. Ejemplo con el \textbf{PRAE} (monitoreo ambiental con datos):

\begin{center}
\footnotesize
\setlength{\tabcolsep}{4pt}
\begin{tabular}{|L{3.2cm}|L{5.5cm}|L{4.8cm}|}
\hline
\rowcolor{vinoInst}
\color{white}\textbf{Área} &
\color{white}\textbf{Aporte al proyecto} &
\color{white}\textbf{Evidencia · momento} \\
\hline
Tecnología e Informática & Base de datos y \textit{dashboard} de visualización de los datos ambientales. & Tablero digital · cierre de G9° P3 \\
\hline
Ciencias Naturales & Variables ambientales a medir y su interpretación. & Informe de campo · durante el periodo \\
\hline
Matemáticas & Estadística descriptiva y lectura de los datos. & Análisis estadístico · mitad de periodo \\
\hline
\end{tabular}
\end{center}

\noindent Este formato es la plantilla reutilizable para los demás proyectos
transversales (ciudadanía digital, emprendimiento, educación sexual).

% ============================================================
\section{Seguimiento y mejoramiento continuo}
% ============================================================

\subsection{Ciclo PHVA del área}

\begin{center}
\begin{tabular}{|L{2cm}|L{5cm}|L{6.5cm}|}
\hline
\rowcolor{vinoInst}
\color{white}\textbf{Fase} &
\color{white}\textbf{Acción} &
\color{white}\textbf{Momento} \\
\hline
\textbf{Planear} &
Actualizar mallas y guías según resultados previos &
Antes de cada año escolar \\
\hline
\textbf{Hacer} &
Ejecutar guías con el MILC y registrar observaciones &
Durante el periodo \\
\hline
\textbf{Verificar} &
Linter de guías (\texttt{make guia-lint}) + análisis de progreso en Plataforma Conéctate &
Al cierre de cada periodo \\
\hline
\textbf{Actuar} &
Ajustar YAMLs de guías fallidas; actualizar el plan si cambian lineamientos MEN &
Antes del siguiente periodo \\
\hline
\end{tabular}
\end{center}

\subsection{Indicadores de calidad del área}

\begin{center}
\begin{tabular}{|L{8cm}|L{4cm}|}
\hline
\rowcolor{vinoInst}
\color{white}\textbf{Indicador} &
\color{white}\textbf{Meta 2026} \\
\hline
\rowcolor{vinoClaro}\multicolumn{2}{|l|}{\textit{De aprendizaje}} \\
\hline
\% estudiantes con desempeño ≥ Básico & ≥ 90\% \\
\hline
\% estudiantes que mejora del diagnóstico a la evaluación final & ≥ 70\% \\
\hline
\% estudiantes en nivel \textit{En proceso} o superior en ≥3 competencias SEL & ≥ 85\% \\
\hline
\rowcolor{vinoClaro}\multicolumn{2}{|l|}{\textit{De gestión y cobertura}} \\
\hline
\% guías publicadas en Plataforma Conéctate & 100\% \\
\hline
\% guías que pasan \texttt{make guia-lint} sin errores & 100\% \\
\hline
\% proyectos integradores sustentados & 100\% \\
\hline
Participación en Bebras u olimpiada de PC & ≥ 50\% de G8°–G10° \\
\hline
\end{tabular}
\end{center}

% ============================================================
\section{Comunicación con familias y acudientes}
% ============================================================

\subsection{Informe mensual automático}

El área usa el sistema de informes automáticos de la Plataforma Conéctate para
enviar mensualmente a los acudientes un reporte con:

\begin{itemize}
  \item Progreso por guía del estudiante (completo / en proceso / sin iniciar).
  \item Desempeños alcanzados en el periodo.
  \item Indicaciones de apoyo desde casa.
\end{itemize}

\subsection{Canales de comunicación}

\begin{center}
\begin{tabular}{|L{5cm}|L{9cm}|}
\hline
\rowcolor{vinoInst}
\color{white}\textbf{Canal} &
\color{white}\textbf{Propósito} \\
\hline
Plataforma Conéctate (panel acudiente) &
Progreso en tiempo real \\
\hline
Informe de periodo (documento institucional) &
Nota y descriptores de desempeño \\
\hline
Reunión de padres (cada periodo) &
Contexto pedagógico del área \\
\hline
WhatsApp institucional &
Comunicados urgentes \\
\hline
\end{tabular}
\end{center}

% ============================================================
\section{Referencias bibliográficas y normativas}
% ============================================================

\subsection{Normativa colombiana}

\begin{itemize}
  \item MEN. (2008). \textit{Guía No. 30: Ser competente en tecnología.}
    Ministerio de Educación Nacional.
  \item MEN. (2022). \textit{Orientaciones curriculares para el área de Tecnología
    e Informática: Pensamiento Computacional como eje transversal.} MEN.
  \item Ley 115 de 1994. Ley General de Educación. Congreso de la República.
  \item Decreto 1290 de 2009. Sistema Institucional de Evaluación. MEN.
  \item Ley 1341 de 2009. Ley de TIC. Congreso de la República.
  \item Ley 23 de 1982. Derechos de autor. Congreso de la República.
  \item Ley 1014 de 2006. Fomento a la cultura del emprendimiento. Congreso.
  \item Ley 2383 de 2024. Educación socioemocional de NNA. Congreso.
  \item Ley 2491 de 2025. Competencias socioemocionales en el PEI. Congreso.
\end{itemize}

\subsection{Referencias pedagógicas y del MILC}

\begin{itemize}
  \item Cárdenas Orozco, Á. (2024). \textit{Modelo de Investigación Liberadora
    y Científica (MILC) — Versión 3.} Documento institucional, I.E. Sor María
    Juliana, Cartago.
  \item Freire, P. (1970). \textit{Pedagogía del oprimido.} Siglo XXI Editores.
  \item Dussel, E. (1977). \textit{Filosofía de la liberación.} Edicol.
  \item Marco Aurelio Antonino. (s. II d. C.). \textit{Meditaciones.}
    (Trad. moderna: J. Cortés, 1994).
  \item Floridi, L. (2014). \textit{The Fourth Revolution.} Oxford University Press.
  \item Bloom, B. S. (Ed.). (1956). \textit{Taxonomy of Educational Objectives.} McKay.
  \item Biggs, J., \& Collis, K. (1982). \textit{Evaluating the Quality of Learning.}
    Academic Press. (SOLO Taxonomy.)
\end{itemize}

\subsection{Referencias de Pensamiento Computacional}

\begin{itemize}
  \item Wing, J. M. (2006). Computational thinking. \textit{Communications of the
    ACM, 49}(3), 33–35.
  \item Brennan, K., \& Resnick, M. (2012). New frameworks for studying computational
    thinking. \textit{AERA 2012.}
  \item Zapata-Ros, M. (2015). Pensamiento computacional: Una nueva alfabetización
    digital. \textit{RED — Revista de Educación a Distancia, 46.}
  \item CSTA \& ISTE. (2011). \textit{Computational Thinking Leadership Toolkit.} ISTE.
\end{itemize}

\subsection{Referencias de diseño instruccional y tecnología educativa}

\begin{itemize}
  \item Mayer, R. E. (2009). \textit{Multimedia Learning} (2.ª ed.).
    Cambridge University Press.
  \item Papert, S. (1980). \textit{Mindstorms.} Basic Books.
  \item Resnick, M. (2017). \textit{Lifelong Kindergarten.} MIT Press.
  \item UNESCO. (2024). \textit{AI Competency Framework for Students.}
    UNESDOC: pf0000391105.
\end{itemize}

% Desactivar zebra: las mallas del anexo tienen su propio esquema de color
\rowcolors{1}{}{}

\clearpage
\appendix
% ============================================================
\section{Anexo A — Mallas curriculares por grado y periodo}
% ============================================================

\noindent Las 18 mallas (6 grados $\times$ 3 periodos) desarrollan, para
cada periodo, las 6 filas de componentes descritas en la sección «Diseño
curricular», en formato apaisado para la correcta lectura de las columnas.

\clearpage
\begin{landscape}
\setlength{\tabcolsep}{4pt}
\renewcommand{\arraystretch}{1.3}

% ------- GRADO 6° · PERIODO 1 -------
\subsection{Grado Sexto}
\subsubsection{Periodo 1 · Comunicación e identidad digital}
\noindent\begin{tabular}{ll}
\textbf{Grado:} Sexto & \textbf{Periodo:} Primero \\
\textbf{Proyecto integrador:} Mi identidad en la red & \textbf{Anclaje:} El pregonero del barrio · Cartago \\
\end{tabular}
\medskip

\begin{longtable}{|L{3.8cm}|L{3.3cm}|L{3.3cm}|L{2.5cm}|L{2.5cm}|L{2.5cm}|}
\caption{Malla curricular — G6° Periodo 1 · Comunicación e identidad digital}
\label{tab:g6p1}\\
\hline
\rowcolor{vinoInst}
\multirow{2}{*}{\color{white}\scriptsize\textbf{\makecell{ESTÁNDARES, LINEAMIENTOS\\U ORIENTACIONES}}} &
\multirow{2}{*}{\color{white}\scriptsize\textbf{\makecell{COMPETENCIAS}}} &
\multirow{2}{*}{\color{white}\scriptsize\textbf{\makecell{APRENDIZAJES}}} &
\multicolumn{3}{c|}{%
  \color{white}\makecell{\textbf{\footnotesize DESCRIPTORES DE DESEMPEÑO}\\[1pt]\textit{\scriptsize Metas del aprendizaje por periodo}}%
}\\
\cline{4-6}
\rowcolor{vinoClaro}
\cellcolor{vinoInst} & \cellcolor{vinoInst} & \cellcolor{vinoInst} &
\textbf{\scriptsize COGNITIVOS} &
\textbf{\scriptsize PROCEDIMENTALES} &
\textbf{\scriptsize ACTITUDINALES} \\
\hline
\endfirsthead
\hline
\rowcolor{vinoInst}
\multirow{2}{*}{\color{white}\scriptsize\textbf{\makecell{ESTÁNDARES, LINEAMIENTOS\\U ORIENTACIONES}}} &
\multirow{2}{*}{\color{white}\scriptsize\textbf{\makecell{COMPETENCIAS}}} &
\multirow{2}{*}{\color{white}\scriptsize\textbf{\makecell{APRENDIZAJES}}} &
\multicolumn{3}{c|}{%
  \color{white}\makecell{\textbf{\footnotesize DESCRIPTORES DE DESEMPEÑO}\\[1pt]\textit{\scriptsize Metas del aprendizaje por periodo}}%
}\\
\cline{4-6}
\rowcolor{vinoClaro}
\cellcolor{vinoInst} & \cellcolor{vinoInst} & \cellcolor{vinoInst} &
\textbf{\scriptsize COGNITIVOS} &
\textbf{\scriptsize PROCEDIMENTALES} &
\textbf{\scriptsize ACTITUDINALES} \\
\hline
\endhead
% -- FILA 1 --
\textbf{Naturaleza y evolución de la tecnología} &
Comprender la evolución de la comunicación digital y la identidad en red. &
Construye una línea de tiempo de la comunicación (carta $\rightarrow$ correo $\rightarrow$ red social). &
Define identidad digital, huella en línea, red y comunicación asincrónica. &
Diseña su línea de tiempo con 5 hitos verificados en fuentes. &
Reconoce que su identidad digital es tan valiosa como su nombre. \\
\hline
% -- FILA 2 --
\textbf{Apropiación y uso de la tecnología} &
Usar herramientas digitales de comunicación con criterio y propósito. &
Maneja correo electrónico institucional con protocolo profesional. &
Conoce la anatomía del correo: asunto, saludo, cuerpo, despedida, firma. &
Redacta y envía 3 correos con protocolo profesional verificado por el docente. &
Cuida el tono y la ortografía como lo haría en una carta impresa. \\
\hline
% -- FILA 3 --
\textbf{Solución de problemas con tecnología} &
Resolver problemas de identidad y privacidad en redes sociales. &
Diagnostica su presencia digital y aplica configuraciones de privacidad. &
Identifica los 5 riesgos principales de privacidad en redes sociales. &
Revisa y ajusta la configuración de privacidad de al menos una red. &
Asume responsabilidad por lo que publica y comparte. \\
\hline
% -- FILA 4 --
\textbf{Tecnología y sociedad} &
Reconocer el impacto social de la comunicación digital en la comunidad. &
Analiza un caso de impacto positivo y uno negativo de las redes en Cartago. &
Conoce casos reales de ciberacoso, desinformación y comunidades positivas. &
Presenta análisis de 2 casos con fuentes verificadas. &
Construye su criterio sobre el uso ético de redes desde su propia experiencia. \\
\hline
% -- FILA 5 (PC) --
\rowcolor{verdePCclaro}
\makecell[l]{\textbf{\color{verdePC}Pensamiento}\\\textbf{\color{verdePC}Computacional}\\
\footnotesize Wing (2006)\\
\footnotesize MEN Or. 2022\\
\footnotesize\textit{Descomposición + Abstracción}} &
Aplicar descomposición y abstracción para analizar la identidad digital como sistema de componentes interrelacionados. &
Descompone su identidad digital en componentes (foto, bio, publicaciones, red de contactos) y los evalúa por separado. &
Reconoce la identidad digital como sistema con módulos independientes. &
Construye un mapa de su identidad digital con 5+ componentes evaluados. &
Acepta que su identidad digital es una construcción consciente; asume responsabilidad por cada componente. \\
\hline
% -- FILA 6 (SEL) --
\rowcolor{moradoSELclaro}
\makecell[l]{\textbf{\color{moradoSEL}Habilidades}\\\textbf{\color{moradoSEL}Socioemocionales}\\
\footnotesize Ley 2383/2024\\
\footnotesize CASEL 2020\\
\footnotesize\textit{Autoconciencia + Toma de decisiones}} &
Desarrollar autoconciencia sobre la propia identidad digital y tomar decisiones responsables. &
Reflexiona sobre cómo su presencia en línea refleja o contradice sus valores personales. &
Identifica sus emociones ante situaciones de presión social en línea. &
Escribe en el cuaderno un protocolo personal de publicación. &
Defiende sus decisiones de privacidad ante la presión del grupo. \\
\hline
\end{longtable}

\medskip
\noindent\textbf{Nota:} Las filas 5 (\textcolor{verdePC}{verde}) y 6 (\textcolor{moradoSEL}{violeta})
son componentes adicionales propios de este plan, complementarios a los cuatro
componentes de la Guía No. 30. Las 17 tablas siguientes (G6P2–G11P3) siguen el
mismo esquema.

% ============================================================
% Sección 8 — Malla curricular por grado y periodo
% Generado automáticamente por gen_tablas.py
% ============================================================

\subsubsection{Periodo 2 · Hardware, software y la máquina por dentro}
\noindent\begin{tabular}{ll}
\textbf{Grado:} Sexto & \textbf{Periodo:} Segundo \\
\textbf{Proyecto integrador:} Maestros del oficio digital: abrir la máquina & \textbf{Anclaje:} Relojero, costurera y panadero de Cartago \\
\end{tabular}
\medskip

\begin{longtable}{|L{3.8cm}|L{3.3cm}|L{3.3cm}|L{2.5cm}|L{2.5cm}|L{2.5cm}|}
\caption{Estructura de malla curricular — G6° Periodo 2 · Hardware, software y la máquina por dentro}
\label{tab:g6p2}\\
\hline
\rowcolor{azulInst}
\multirow{2}{*}{\color{white}\scriptsize\textbf{\makecell{ESTÁNDARES, LINEAMIENTOS\\U ORIENTACIONES}}} &
\multirow{2}{*}{\color{white}\scriptsize\textbf{\makecell{COMPETENCIAS}}} &
\multirow{2}{*}{\color{white}\scriptsize\textbf{\makecell{APRENDIZAJES}}} &
\multicolumn{3}{c|}{%
  \color{white}\makecell{\textbf{\footnotesize DESCRIPTORES DE DESEMPEÑO}\\[1pt]\textit{\scriptsize Metas del aprendizaje por periodo}}%
}\\
\cline{4-6}
\rowcolor{azulClaro}
\cellcolor{azulInst} & \cellcolor{azulInst} & \cellcolor{azulInst} &
\textbf{\scriptsize COGNITIVOS} &
\textbf{\scriptsize PROCEDIMENTALES} &
\textbf{\scriptsize ACTITUDINALES} \\
\hline
\endfirsthead
\hline
\rowcolor{azulInst}
\multirow{2}{*}{\color{white}\scriptsize\textbf{\makecell{ESTÁNDARES, LINEAMIENTOS\\U ORIENTACIONES}}} &
\multirow{2}{*}{\color{white}\scriptsize\textbf{\makecell{COMPETENCIAS}}} &
\multirow{2}{*}{\color{white}\scriptsize\textbf{\makecell{APRENDIZAJES}}} &
\multicolumn{3}{c|}{%
  \color{white}\makecell{\textbf{\footnotesize DESCRIPTORES DE DESEMPEÑO}\\[1pt]\textit{\scriptsize Metas del aprendizaje por periodo}}%
}\\
\cline{4-6}
\rowcolor{azulClaro}
\cellcolor{azulInst} & \cellcolor{azulInst} & \cellcolor{azulInst} &
\textbf{\scriptsize COGNITIVOS} &
\textbf{\scriptsize PROCEDIMENTALES} &
\textbf{\scriptsize ACTITUDINALES} \\
\hline
\endhead

% -- FILA 1: Naturaleza y evolución --
\textbf{Naturaleza y evolución de la tecnología} &
Comprender las 4 funciones básicas de toda máquina inteligente y reconocer la evolución del hardware. &
Identifica las 4 funciones (E/P/S/A) y las piezas del gabinete (S2–S3). &
Reconoce CPU, RAM, disco, fuente y tarjeta madre, comprendiendo el papel de cada pieza en el funcionamiento del sistema. &
Elabora un mapa anatómico del computador con 8 piezas identificadas y descripción propia. &
Valora la habilidad del relojero/costurera/panadero como antecedente del técnico actual. \\
\hline

% -- FILA 2: Apropiación y uso --
\textbf{Apropiación y uso de la tecnología} &
Manejar periféricos, sistema operativo y archivos digitales con criterio de organización profesional. &
Distingue periféricos de entrada, salida y mixtos (S4), comprende qué es el SO (S6) y organiza archivos en carpetas (S7). &
Diferencia hardware de software y conoce los principales sistemas operativos. &
Organiza sus archivos en estructura profesional (año/materia/periodo) con nomenclatura correcta. &
Aplica buenas prácticas en la sala de sistemas y respeta los equipos compartidos. \\
\hline

% -- FILA 3: Solución de problemas --
\textbf{Solución de problemas con tecnología} &
Diagnosticar problemas básicos de un equipo informático y aplicar mantenimiento preventivo. &
Aplica las 5 reglas de cuidado de la sala (S1, S8, S9) y reconoce señales de fallo (sobrecalentamiento, lentitud). &
Identifica las 3 acciones de mantenimiento básico: limpiar, actualizar, respaldar. &
Realiza limpieza física suave del equipo, actualiza el SO y crea respaldos de archivos. &
Cuida su cuerpo aplicando la regla 20-20-20 y la postura ergonómica. \\
\hline

% -- FILA 4: Tecnología y sociedad --
\textbf{Tecnología y sociedad} &
Elegir un equipo informático con criterio razonado, evitando el extremismo del precio. &
Construye una ficha técnica del computador ideal para estudio con presupuesto razonado (S10). &
Reconoce las especificaciones mínimas razonables para uso escolar (CPU i3/Ryzen3, RAM 8 GB, SSD 256 GB). &
Elabora una ficha técnica con marca, modelo, especificaciones y precio realista. &
Aplica criterio frente a la publicidad; reconoce que lo barato sale caro. \\
\hline

% -- FILA 5: PC (verde) --
\rowcolor{verdePCclaro}
\makecell[l]{\textbf{\color{verdePC}Pensamiento}\\\textbf{\color{verdePC}Computacional}\\
\footnotesize Wing (2006)\\\footnotesize MEN Or. 2022\\\footnotesize\textit{Descomposición · Abstracción}}\
 &
Aplicar la descomposición y la abstracción para comprender el sistema informático como capas funcionales interrelacionadas. &
Descompone el computador en sus 4 funciones (E/P/S/A) y abstrae el SO como capa de intermediación (S2–S6). &
Describe el computador como sistema jerárquico de capas (hardware / SO / aplicaciones / usuario). &
Elabora un diagrama de capas del sistema informático con funciones interrelacionadas. &
Valora la arquitectura del computador como diseño humano inteligente; entiende el sistema para usarlo con criterio. \\
\hline

% -- FILA 6: SEL (violeta) --
\rowcolor{moradoSELclaro}
\makecell[l]{\textbf{\color{moradoSEL}Habilidades}\\\textbf{\color{moradoSEL}Socioemocionales}\\
\footnotesize Ley 2383/2024\\\footnotesize CASEL 2020\\\footnotesize\textit{Autogestión · Conciencia social}}\
 &
Desarrollar autogestión de la frustración frente a equipos que fallan y conciencia del cuidado colectivo de los bienes tecnológicos compartidos. &
Identifica sus propias reacciones emocionales cuando un equipo falla y reconoce que los bienes tecnológicos de la sala son responsabilidad colectiva. &
Distingue entre frustración individual y responsabilidad colectiva en el uso de equipos compartidos. &
Aplica la regla de pausa-respira-reporta cuando un equipo presenta fallas, sin ocultarlas. &
Cuida los equipos de la sala como bienes comunes y reporta fallas con honestidad para el bien del grupo. \\
\hline

\end{longtable}
\medskip
\noindent\textbf{Nota:} Las filas 5 (\textcolor{verdePC}{verde}) y 6 (\textcolor{moradoSEL}{violeta}) son componentes propios de este plan, adicionales a los cuatro componentes de la Guía No.~30 del MEN.

\subsubsection{Periodo 3 · Escribir y buscar con criterio}
\noindent\begin{tabular}{ll}
\textbf{Grado:} Sexto & \textbf{Periodo:} Tercero \\
\textbf{Proyecto integrador:} Reporteros con criterio del barrio & \textbf{Anclaje:} Doña Mercedes, maestra rural (La Plata, Cartago) \\
\end{tabular}
\medskip

\begin{longtable}{|L{3.8cm}|L{3.3cm}|L{3.3cm}|L{2.5cm}|L{2.5cm}|L{2.5cm}|}
\caption{Estructura de malla curricular — G6° Periodo 3 · Escribir y buscar con criterio}
\label{tab:g6p3}\\
\hline
\rowcolor{azulInst}
\multirow{2}{*}{\color{white}\scriptsize\textbf{\makecell{ESTÁNDARES, LINEAMIENTOS\\U ORIENTACIONES}}} &
\multirow{2}{*}{\color{white}\scriptsize\textbf{\makecell{COMPETENCIAS}}} &
\multirow{2}{*}{\color{white}\scriptsize\textbf{\makecell{APRENDIZAJES}}} &
\multicolumn{3}{c|}{%
  \color{white}\makecell{\textbf{\footnotesize DESCRIPTORES DE DESEMPEÑO}\\[1pt]\textit{\scriptsize Metas del aprendizaje por periodo}}%
}\\
\cline{4-6}
\rowcolor{azulClaro}
\cellcolor{azulInst} & \cellcolor{azulInst} & \cellcolor{azulInst} &
\textbf{\scriptsize COGNITIVOS} &
\textbf{\scriptsize PROCEDIMENTALES} &
\textbf{\scriptsize ACTITUDINALES} \\
\hline
\endfirsthead
\hline
\rowcolor{azulInst}
\multirow{2}{*}{\color{white}\scriptsize\textbf{\makecell{ESTÁNDARES, LINEAMIENTOS\\U ORIENTACIONES}}} &
\multirow{2}{*}{\color{white}\scriptsize\textbf{\makecell{COMPETENCIAS}}} &
\multirow{2}{*}{\color{white}\scriptsize\textbf{\makecell{APRENDIZAJES}}} &
\multicolumn{3}{c|}{%
  \color{white}\makecell{\textbf{\footnotesize DESCRIPTORES DE DESEMPEÑO}\\[1pt]\textit{\scriptsize Metas del aprendizaje por periodo}}%
}\\
\cline{4-6}
\rowcolor{azulClaro}
\cellcolor{azulInst} & \cellcolor{azulInst} & \cellcolor{azulInst} &
\textbf{\scriptsize COGNITIVOS} &
\textbf{\scriptsize PROCEDIMENTALES} &
\textbf{\scriptsize ACTITUDINALES} \\
\hline
\endhead

% -- FILA 1: Naturaleza y evolución --
\textbf{Naturaleza y evolución de la tecnología} &
Comprender qué es internet, cómo viajan los datos y la diferencia entre internet y la World Wide Web. &
Reconoce la arquitectura cliente-servidor y los componentes de internet (S7). &
Describe internet como red de redes conectada por cables submarinos, satélites y antenas. &
Explica con ejemplos cómo viaja una petición desde su navegador hasta un servidor y de regreso. &
Aprecia internet como bien común global que requiere cuidado y mantenimiento colectivo. \\
\hline

% -- FILA 2: Apropiación y uso --
\textbf{Apropiación y uso de la tecnología} &
Producir documentos digitales con formato profesional usando un procesador de texto. &
Aplica formato básico, listas, tablas, imágenes y encabezados profesionales (S2–S6). &
Identifica los elementos de un documento profesional: fuente, tamaño, márgenes, alineación, listas, tablas, encabezado, pie y números de página. &
Produce un documento informativo de 4 páginas con todos los elementos de formato profesional. &
Cuida la presentación y valora la legibilidad como respeto al lector. \\
\hline

% -- FILA 3: Solución de problemas --
\textbf{Solución de problemas con tecnología} &
Buscar información con operadores avanzados y evaluar fuentes con criterio CRAAP para detectar fake news. &
Aplica operadores (site:, comillas, guion menos) y evalúa fuentes con las 5 letras CRAAP (S8–S9). &
Reconoce los 5 criterios CRAAP y las 5 señales típicas de fake news. &
Investiga un tema usando al menos 3 fuentes evaluadas con CRAAP; elabora guía visual anti-fake news. &
Antes de creer o reenviar, se detiene, respira y verifica; no comparte rumores. \\
\hline

% -- FILA 4: Tecnología y sociedad --
\textbf{Tecnología y sociedad} &
Producir y publicar información con responsabilidad ciudadana, citando fuentes y respetando autoría. &
Cita fuentes correctamente, evita el plagio y reconoce la responsabilidad de quien publica (S10). &
Identifica la importancia de citar fuentes, la honestidad académica y el rol del productor responsable en la infoesfera. &
Entrega un informe ciudadano del barrio con mínimo 2 fuentes citadas y declaración de uso de IA. &
Asume su rol como ciudadano digital productor y aporta claridad a su comunidad. \\
\hline

% -- FILA 5: PC (verde) --
\rowcolor{verdePCclaro}
\makecell[l]{\textbf{\color{verdePC}Pensamiento}\\\textbf{\color{verdePC}Computacional}\\
\footnotesize Wing (2006)\\\footnotesize MEN Or. 2022\\\footnotesize\textit{Diseño de algoritmos · Evaluación}}\
 &
Aplicar algoritmos de búsqueda y evaluación de fuentes como procesos lógicos secuenciales para tomar decisiones informadas. &
Aplica operadores de búsqueda como instrucciones secuenciales; usa CRAAP como algoritmo de verificación de 5 pasos. &
Reconoce los operadores de búsqueda como instrucciones precisas para un motor algorítmico; describe CRAAP como árbol de decisión. &
Diseña su propio flujo de búsqueda y evaluación como diagrama de pasos con condiciones visibles. &
Rechaza el pensamiento impulsivo y adopta la lógica de quien verifica antes de actuar. \\
\hline

% -- FILA 6: SEL (violeta) --
\rowcolor{moradoSELclaro}
\makecell[l]{\textbf{\color{moradoSEL}Habilidades}\\\textbf{\color{moradoSEL}Socioemocionales}\\
\footnotesize Ley 2383/2024\\\footnotesize CASEL 2020\\\footnotesize\textit{Autogestión · Toma de decisiones responsable}}\
 &
Regular el impulso de creer y compartir sin verificar, y tomar decisiones responsables como productor de información en la infoesfera. &
Identifica la emoción de urgencia que genera el contenido viral y reconoce cuándo esa emoción puede llevar a errores de juicio. &
Describe los efectos sociales de compartir información falsa y el papel del ciudadano digital en frenarlo. &
Aplica la secuencia pausa-verifica-decide antes de compartir cualquier contenido en redes. &
Asume responsabilidad por la información que produce y difunde, reconociendo su impacto en la comunidad. \\
\hline

\end{longtable}
\medskip
\noindent\textbf{Nota:} Las filas 5 (\textcolor{verdePC}{verde}) y 6 (\textcolor{moradoSEL}{violeta}) son componentes propios de este plan, adicionales a los cuatro componentes de la Guía No.~30 del MEN.

\subsection{Grado Séptimo}

\subsubsection{Periodo 1 · Trabajo colaborativo en la nube (Microsoft 365)}
\noindent\begin{tabular}{ll}
\textbf{Grado:} Séptimo & \textbf{Periodo:} Primero \\
\textbf{Proyecto integrador:} Minga digital: una obra colaborativa en la nube & \textbf{Anclaje:} La minga del Valle del Cauca \\
\end{tabular}
\medskip

\begin{longtable}{|L{3.8cm}|L{3.3cm}|L{3.3cm}|L{2.5cm}|L{2.5cm}|L{2.5cm}|}
\caption{Estructura de malla curricular — G7° Periodo 1 · Trabajo colaborativo en la nube (Microsoft 365)}
\label{tab:g7p1}\\
\hline
\rowcolor{azulInst}
\multirow{2}{*}{\color{white}\scriptsize\textbf{\makecell{ESTÁNDARES, LINEAMIENTOS\\U ORIENTACIONES}}} &
\multirow{2}{*}{\color{white}\scriptsize\textbf{\makecell{COMPETENCIAS}}} &
\multirow{2}{*}{\color{white}\scriptsize\textbf{\makecell{APRENDIZAJES}}} &
\multicolumn{3}{c|}{%
  \color{white}\makecell{\textbf{\footnotesize DESCRIPTORES DE DESEMPEÑO}\\[1pt]\textit{\scriptsize Metas del aprendizaje por periodo}}%
}\\
\cline{4-6}
\rowcolor{azulClaro}
\cellcolor{azulInst} & \cellcolor{azulInst} & \cellcolor{azulInst} &
\textbf{\scriptsize COGNITIVOS} &
\textbf{\scriptsize PROCEDIMENTALES} &
\textbf{\scriptsize ACTITUDINALES} \\
\hline
\endfirsthead
\hline
\rowcolor{azulInst}
\multirow{2}{*}{\color{white}\scriptsize\textbf{\makecell{ESTÁNDARES, LINEAMIENTOS\\U ORIENTACIONES}}} &
\multirow{2}{*}{\color{white}\scriptsize\textbf{\makecell{COMPETENCIAS}}} &
\multirow{2}{*}{\color{white}\scriptsize\textbf{\makecell{APRENDIZAJES}}} &
\multicolumn{3}{c|}{%
  \color{white}\makecell{\textbf{\footnotesize DESCRIPTORES DE DESEMPEÑO}\\[1pt]\textit{\scriptsize Metas del aprendizaje por periodo}}%
}\\
\cline{4-6}
\rowcolor{azulClaro}
\cellcolor{azulInst} & \cellcolor{azulInst} & \cellcolor{azulInst} &
\textbf{\scriptsize COGNITIVOS} &
\textbf{\scriptsize PROCEDIMENTALES} &
\textbf{\scriptsize ACTITUDINALES} \\
\hline
\endhead

% -- FILA 1: Naturaleza y evolución --
\textbf{Naturaleza y evolución de la tecnología} &
Comprender el paso de Office tradicional a Microsoft 365 y el concepto de computación en la nube. &
Identifica qué es la nube, diferencia Office tradicional de M365 y conoce OneDrive (S2–S3). &
Describe la arquitectura cliente-nube y el rol del navegador como portal universal de acceso. &
Accede a office.com con su cuenta institucional y explora OneDrive, Word, Excel y PowerPoint en línea. &
Reconoce la minga del Valle como ancestro cultural de la coautoría digital. \\
\hline

% -- FILA 2: Apropiación y uso --
\textbf{Apropiación y uso de la tecnología} &
Trabajar en equipo en archivos compartidos con coautoría asincrónica y comentarios constructivos. &
Comparte archivos con permisos, edita en coautoría asincrónica y usa comentarios (S5–S7). &
Distingue los 3 niveles de permisos (lectura, comentarios, edición) y comprende cuándo usar cada uno. &
Coedita un Word de 4 páginas con 2 compañeros, usa comentarios para sugerencias y resuelve revisiones. &
Respeta la autoría del compañero, hace comentarios constructivos y participa con reciprocidad como en la minga. \\
\hline

% -- FILA 3: Solución de problemas --
\textbf{Solución de problemas con tecnología} &
Recuperar versiones anteriores de un archivo y gestionar conflictos de edición con criterio. &
Aplica el historial de versiones cuando algo se daña en el archivo compartido (S8). &
Identifica que Microsoft 365 guarda automáticamente versiones y permite recuperar contenido perdido. &
Practica restaurar versiones anteriores cuando algo se daña, sin pánico. &
Asume la solución técnica y humana frente a errores del equipo. \\
\hline

% -- FILA 4: Tecnología y sociedad --
\textbf{Tecnología y sociedad} &
Gestionar reuniones virtuales y comunicación profesional con Outlook y Teams. &
Maneja Outlook y Teams en el contexto escolar (S9). &
Reconoce los protocolos de invitación, aceptación y cortesía en reuniones virtuales. &
Acepta invitaciones por Outlook, se une a reuniones por Teams y participa con criterio profesional. &
Aprecia la nube como bien común del equipo: cuida lo que pone allí porque es ambiente compartido. \\
\hline

% -- FILA 5: PC (verde) --
\rowcolor{verdePCclaro}
\makecell[l]{\textbf{\color{verdePC}Pensamiento}\\\textbf{\color{verdePC}Computacional}\\
\footnotesize Wing (2006)\\\footnotesize MEN Or. 2022\\\footnotesize\textit{Abstracción · Reconocimiento de patrones}}\
 &
Aplicar la abstracción para modelar procesos colaborativos en la nube y comprender la sincronización como algoritmo distribuido. &
Abstrae el flujo de coedición como proceso con reglas de concurrencia (S5–S7); modela conflictos de versiones como problema algorítmico. &
Describe la nube como sistema distribuido con reglas de sincronización, permisos y estados. &
Elabora un diagrama de flujo del proceso de coedición: estados (editando/guardado/en conflicto) y transiciones. &
Comprende que los sistemas digitales tienen lógica interna; actúa con criterio sobre los estados y permisos. \\
\hline

% -- FILA 6: SEL (violeta) --
\rowcolor{moradoSELclaro}
\makecell[l]{\textbf{\color{moradoSEL}Habilidades}\\\textbf{\color{moradoSEL}Socioemocionales}\\
\footnotesize Ley 2383/2024\\\footnotesize CASEL 2020\\\footnotesize\textit{Habilidades relacionales · Conciencia social}}\
 &
Desarrollar habilidades de colaboración auténtica en equipos digitales, respetando la autoría y aportando con reciprocidad como en la minga del Valle. &
Identifica las normas de colaboración digital justa y reconoce cómo sus acciones en el archivo compartido afectan al equipo. &
Describe la reciprocidad de la minga como modelo de trabajo colectivo aplicable a la coautoría digital. &
Participa activamente en la coedición del documento, hace aportes concretos y usa comentarios para dar retroalimentación respetuosa. &
Cuida la autoría de sus compañeros como cuida la propia; actúa con la reciprocidad de quien trabaja en minga. \\
\hline

\end{longtable}
\medskip
\noindent\textbf{Nota:} Las filas 5 (\textcolor{verdePC}{verde}) y 6 (\textcolor{moradoSEL}{violeta}) son componentes propios de este plan, adicionales a los cuatro componentes de la Guía No.~30 del MEN.

\subsubsection{Periodo 2 · Algoritmia y pensamiento computacional (Scratch)}
\noindent\begin{tabular}{ll}
\textbf{Grado:} Séptimo & \textbf{Periodo:} Segundo \\
\textbf{Proyecto integrador:} Tejedoras digitales: mi primer programa con criterio & \textbf{Anclaje:} Doña Sofía, la tejedora del Valle \\
\end{tabular}
\medskip

\begin{longtable}{|L{3.8cm}|L{3.3cm}|L{3.3cm}|L{2.5cm}|L{2.5cm}|L{2.5cm}|}
\caption{Estructura de malla curricular — G7° Periodo 2 · Algoritmia y pensamiento computacional (Scratch)}
\label{tab:g7p2}\\
\hline
\rowcolor{azulInst}
\multirow{2}{*}{\color{white}\scriptsize\textbf{\makecell{ESTÁNDARES, LINEAMIENTOS\\U ORIENTACIONES}}} &
\multirow{2}{*}{\color{white}\scriptsize\textbf{\makecell{COMPETENCIAS}}} &
\multirow{2}{*}{\color{white}\scriptsize\textbf{\makecell{APRENDIZAJES}}} &
\multicolumn{3}{c|}{%
  \color{white}\makecell{\textbf{\footnotesize DESCRIPTORES DE DESEMPEÑO}\\[1pt]\textit{\scriptsize Metas del aprendizaje por periodo}}%
}\\
\cline{4-6}
\rowcolor{azulClaro}
\cellcolor{azulInst} & \cellcolor{azulInst} & \cellcolor{azulInst} &
\textbf{\scriptsize COGNITIVOS} &
\textbf{\scriptsize PROCEDIMENTALES} &
\textbf{\scriptsize ACTITUDINALES} \\
\hline
\endfirsthead
\hline
\rowcolor{azulInst}
\multirow{2}{*}{\color{white}\scriptsize\textbf{\makecell{ESTÁNDARES, LINEAMIENTOS\\U ORIENTACIONES}}} &
\multirow{2}{*}{\color{white}\scriptsize\textbf{\makecell{COMPETENCIAS}}} &
\multirow{2}{*}{\color{white}\scriptsize\textbf{\makecell{APRENDIZAJES}}} &
\multicolumn{3}{c|}{%
  \color{white}\makecell{\textbf{\footnotesize DESCRIPTORES DE DESEMPEÑO}\\[1pt]\textit{\scriptsize Metas del aprendizaje por periodo}}%
}\\
\cline{4-6}
\rowcolor{azulClaro}
\cellcolor{azulInst} & \cellcolor{azulInst} & \cellcolor{azulInst} &
\textbf{\scriptsize COGNITIVOS} &
\textbf{\scriptsize PROCEDIMENTALES} &
\textbf{\scriptsize ACTITUDINALES} \\
\hline
\endhead

% -- FILA 1: Naturaleza y evolución --
\textbf{Naturaleza y evolución de la tecnología} &
Reconocer las raíces ancestrales del pensamiento computacional en oficios tradicionales (tejido, cocina, agricultura). &
Identifica las 4 técnicas del PC y reconoce que las tejedoras ya programaban (S1–S3). &
Define algoritmo, secuencia, repetición, condicional y variable, y los reconoce en oficios cotidianos. &
Descompone un problema cotidiano (cepillarse, cocinar, tejer) en pasos algorítmicos verificables. &
Honra a las tejedoras del Valle como primeras programadoras y reconoce la dignidad de los saberes ancestrales. \\
\hline

% -- FILA 2: Apropiación y uso --
\textbf{Apropiación y uso de la tecnología} &
Programar en Scratch usando variables, condicionales y bucles con propósito claro. &
Construye programas en Scratch con personajes que responden al usuario (S8–S9). &
Conoce los bloques principales de Scratch: eventos, control, sensores, variables, operadores. &
Programa un juego o narrativa interactiva con al menos 2 personajes, 2 escenas, 2 variables, 2 condicionales y 2 bucles. &
Programa con paciencia y claridad: variables con nombre, bucles con propósito, condiciones sin duplicar. \\
\hline

% -- FILA 3: Solución de problemas --
\textbf{Solución de problemas con tecnología} &
Aplicar diagramas de flujo y depuración para resolver problemas algorítmicos. &
Elabora diagramas de flujo con figuras correctas (S4) y depura programas paso a paso (S8–S9). &
Reconoce las figuras del diagrama de flujo (óvalo/rectángulo/rombo/flecha) y su significado. &
Diseña un diagrama de flujo en cuaderno antes de programar; prueba el programa con usuarios externos. &
Acepta los errores como parte del proceso y depura con paciencia. \\
\hline

% -- FILA 4: Tecnología y sociedad --
\textbf{Tecnología y sociedad} &
Producir un programa con instructivo de uso para que otros lo entiendan y usen. &
Documenta su programa con instructivo de uso de 1 página (S10). &
Comprende que un programa debe ser entendible por otros, no solo por su autor. &
Redacta un instructivo de uso con 7 partes y lo prueba con un compañero externo al equipo. &
Reconoce que el código que escribe hoy se ejecuta automáticamente después; es responsable de lo que su programa hace. \\
\hline

% -- FILA 5: PC (verde) --
\rowcolor{verdePCclaro}
\makecell[l]{\textbf{\color{verdePC}Pensamiento}\\\textbf{\color{verdePC}Computacional}\\
\footnotesize Wing (2006)\\\footnotesize MEN Or. 2022\\\footnotesize\textit{Descomposición · Patrones · Abstracción · Algoritmos · Depuración}}\
 &
Aplicar las 4 técnicas centrales del PC (descomposición, patrones, abstracción, algoritmos) para diseñar y depurar programas en Scratch. &
Diseña algoritmos con secuencia, repetición y condicional (S1–S4); programa variables y bucles (S5–S8); depura y documenta (S9–S10). &
Define descomposición, reconocimiento de patrones, abstracción y diseño de algoritmos; los aplica en Scratch. &
Crea un programa en Scratch (2+ personajes, 2 escenas, 2 variables, 2 condicionales, 2 bucles); traza diagrama de flujo antes de programar. &
Acepta los errores como parte de la depuración; honra las tejedoras del Valle como primeras programadoras. \\
\hline

% -- FILA 6: SEL (violeta) --
\rowcolor{moradoSELclaro}
\makecell[l]{\textbf{\color{moradoSEL}Habilidades}\\\textbf{\color{moradoSEL}Socioemocionales}\\
\footnotesize Ley 2383/2024\\\footnotesize CASEL 2020\\\footnotesize\textit{Autogestión · Toma de decisiones}}\
 &
Desarrollar perseverancia y autogestión emocional frente a los errores de programación, aprendiendo a depurar con calma y a tomar decisiones metódicas ante problemas complejos. &
Reconoce la frustración como señal de aprendizaje y aplica estrategias de calma para depurar sin rendirse. &
Distingue entre error de sintaxis, error de lógica y error de diseño; identifica la estrategia adecuada para cada tipo. &
Aplica la rutina depuración-paso-a-paso: aislar el error, formular hipótesis, probar, registrar resultado. &
Persiste con ecuanimidad frente a errores difíciles; reconoce que depurar bien es tan valioso como programar bien. \\
\hline

\end{longtable}
\medskip
\noindent\textbf{Nota:} Las filas 5 (\textcolor{verdePC}{verde}) y 6 (\textcolor{moradoSEL}{violeta}) son componentes propios de este plan, adicionales a los cuatro componentes de la Guía No.~30 del MEN.

\subsubsection{Periodo 3 · Inteligencia artificial con criterio}
\noindent\begin{tabular}{ll}
\textbf{Grado:} Séptimo & \textbf{Periodo:} Tercero \\
\textbf{Proyecto integrador:} Consejeros digitales: IA con criterio para mi salón & \textbf{Anclaje:} El consejero del barrio (don Lucho, doña Mercedes, doña Esperanza) \\
\end{tabular}
\medskip

\begin{longtable}{|L{3.8cm}|L{3.3cm}|L{3.3cm}|L{2.5cm}|L{2.5cm}|L{2.5cm}|}
\caption{Estructura de malla curricular — G7° Periodo 3 · Inteligencia artificial con criterio}
\label{tab:g7p3}\\
\hline
\rowcolor{azulInst}
\multirow{2}{*}{\color{white}\scriptsize\textbf{\makecell{ESTÁNDARES, LINEAMIENTOS\\U ORIENTACIONES}}} &
\multirow{2}{*}{\color{white}\scriptsize\textbf{\makecell{COMPETENCIAS}}} &
\multirow{2}{*}{\color{white}\scriptsize\textbf{\makecell{APRENDIZAJES}}} &
\multicolumn{3}{c|}{%
  \color{white}\makecell{\textbf{\footnotesize DESCRIPTORES DE DESEMPEÑO}\\[1pt]\textit{\scriptsize Metas del aprendizaje por periodo}}%
}\\
\cline{4-6}
\rowcolor{azulClaro}
\cellcolor{azulInst} & \cellcolor{azulInst} & \cellcolor{azulInst} &
\textbf{\scriptsize COGNITIVOS} &
\textbf{\scriptsize PROCEDIMENTALES} &
\textbf{\scriptsize ACTITUDINALES} \\
\hline
\endfirsthead
\hline
\rowcolor{azulInst}
\multirow{2}{*}{\color{white}\scriptsize\textbf{\makecell{ESTÁNDARES, LINEAMIENTOS\\U ORIENTACIONES}}} &
\multirow{2}{*}{\color{white}\scriptsize\textbf{\makecell{COMPETENCIAS}}} &
\multirow{2}{*}{\color{white}\scriptsize\textbf{\makecell{APRENDIZAJES}}} &
\multicolumn{3}{c|}{%
  \color{white}\makecell{\textbf{\footnotesize DESCRIPTORES DE DESEMPEÑO}\\[1pt]\textit{\scriptsize Metas del aprendizaje por periodo}}%
}\\
\cline{4-6}
\rowcolor{azulClaro}
\cellcolor{azulInst} & \cellcolor{azulInst} & \cellcolor{azulInst} &
\textbf{\scriptsize COGNITIVOS} &
\textbf{\scriptsize PROCEDIMENTALES} &
\textbf{\scriptsize ACTITUDINALES} \\
\hline
\endhead

% -- FILA 1: Naturaleza y evolución --
\textbf{Naturaleza y evolución de la tecnología} &
Comprender qué es la inteligencia artificial, sus tipos y cómo aprende. &
Identifica IA en su vida diaria (S2), distingue IA estrecha de general (S3) y comprende el entrenamiento (S4). &
Define IA, LLM, dataset y entrenamiento; distingue aprendizaje supervisado y no supervisado. &
Identifica al menos 5 ejemplos de IA en su vida diaria (TikTok, Siri, filtros, etc.). &
Reconoce que el consejero del barrio aconsejaba pero no decidía; la IA hoy ocupa ese rol con sus límites. \\
\hline

% -- FILA 2: Apropiación y uso --
\textbf{Apropiación y uso de la tecnología} &
Formular prompts efectivos a modelos de lenguaje con criterio profesional. &
Aplica los 4 elementos del prompt (rol, contexto, tarea, formato) y lo itera (S6–S7). &
Conoce los componentes de un prompt efectivo y las técnicas de iteración. &
Construye una guía de 6 prompts útiles para grado 7°, probados en una IA real. &
Pide bien hoy para no editar mañana; el tiempo invertido en formular se recupera al multiplicar. \\
\hline

% -- FILA 3: Solución de problemas --
\textbf{Solución de problemas con tecnología} &
Verificar la información que produce una IA y detectar alucinaciones. &
Reconoce alucinaciones y deshonestidad académica con IA (S8); verifica datos contra fuentes externas. &
Identifica las alucinaciones como respuestas estadísticamente probables pero no verdaderas. &
Verifica al menos 2 datos generados por IA contra fuentes oficiales antes de incluirlos en un trabajo. &
Mantiene escepticismo saludable: la IA habla rápido y seguro, pero requiere verificación humana. \\
\hline

% -- FILA 4: Tecnología y sociedad --
\textbf{Tecnología y sociedad} &
Declarar honestamente el uso de IA en cualquier trabajo académico y aplicar la política institucional. &
Aplica la declaración honesta de uso de IA (S9) y redacta el decálogo ético del salón (S10). &
Conoce la política institucional de IA por grados y los 5 principios transversales no negociables. &
Incluye sección "Cómo trabajamos con IA" en sus entregables especificando modelo, partes, edición y verificación. &
Habita la infoesfera como productor responsable: la IA amplifica su voz, no la silencia. \\
\hline

% -- FILA 5: PC (verde) --
\rowcolor{verdePCclaro}
\makecell[l]{\textbf{\color{verdePC}Pensamiento}\\\textbf{\color{verdePC}Computacional}\\
\footnotesize Wing (2006)\\\footnotesize MEN Or. 2022\\\footnotesize\textit{Reconocimiento de patrones · Abstracción}}\
 &
Aplicar el pensamiento computacional para comprender cómo los modelos de IA reconocen patrones y abstraen información para producir respuestas. &
Identifica el proceso de entrenamiento de una IA como reconocimiento de patrones a gran escala (S2–S4); diseña prompts como instrucciones algorítmicas (S6–S7). &
Describe el aprendizaje automático como proceso de abstracción estadística de patrones; reconoce que un LLM es una función probabilística, no un razonador con conciencia. &
Diseña 6 prompts con estructura algorítmica (rol + contexto + tarea + formato + restricciones + verificación); itera y evalúa contra fuentes externas. &
Mantiene escepticismo saludable: verifica antes de aceptar; comprende que la IA no piensa sino que procesa patrones estadísticos. \\
\hline

% -- FILA 6: SEL (violeta) --
\rowcolor{moradoSELclaro}
\makecell[l]{\textbf{\color{moradoSEL}Habilidades}\\\textbf{\color{moradoSEL}Socioemocionales}\\
\footnotesize Ley 2383/2024\\\footnotesize CASEL 2020\\\footnotesize\textit{Autoconciencia · Toma de decisiones responsable}}\
 &
Reconocer cómo los algoritmos de IA pueden influir en las propias creencias y emociones, y desarrollar criterio autónomo para tomar decisiones informadas frente a contenido generado por IA. &
Identifica situaciones en las que una IA puede generar sesgos en sus opiniones o emociones (burbujas de filtro, contenido personalizado). &
Describe el mecanismo de personalización algorítmica y su impacto en la formación de creencias. &
Analiza su propia experiencia con plataformas algorítmicas e identifica 3 casos en que el algoritmo influyó en lo que creyó o sintió. &
Ejerce su autonomía crítica frente al contenido generado por IA; decide con criterio propio, no por lo que la IA produce más fácil. \\
\hline

\end{longtable}
\medskip
\noindent\textbf{Nota:} Las filas 5 (\textcolor{verdePC}{verde}) y 6 (\textcolor{moradoSEL}{violeta}) son componentes propios de este plan, adicionales a los cuatro componentes de la Guía No.~30 del MEN.

\subsection{Grado Octavo}

\subsubsection{Periodo 1 · Lógica computacional avanzada}
\noindent\begin{tabular}{ll}
\textbf{Grado:} Octavo & \textbf{Periodo:} Primero \\
\textbf{Proyecto integrador:} Decisiones lógicas: del fogón al algoritmo & \textbf{Anclaje:} El gesto de la cocina del Valle \\
\end{tabular}
\medskip

\begin{longtable}{|L{3.8cm}|L{3.3cm}|L{3.3cm}|L{2.5cm}|L{2.5cm}|L{2.5cm}|}
\caption{Estructura de malla curricular — G8° Periodo 1 · Lógica computacional avanzada}
\label{tab:g8p1}\\
\hline
\rowcolor{azulInst}
\multirow{2}{*}{\color{white}\scriptsize\textbf{\makecell{ESTÁNDARES, LINEAMIENTOS\\U ORIENTACIONES}}} &
\multirow{2}{*}{\color{white}\scriptsize\textbf{\makecell{COMPETENCIAS}}} &
\multirow{2}{*}{\color{white}\scriptsize\textbf{\makecell{APRENDIZAJES}}} &
\multicolumn{3}{c|}{%
  \color{white}\makecell{\textbf{\footnotesize DESCRIPTORES DE DESEMPEÑO}\\[1pt]\textit{\scriptsize Metas del aprendizaje por periodo}}%
}\\
\cline{4-6}
\rowcolor{azulClaro}
\cellcolor{azulInst} & \cellcolor{azulInst} & \cellcolor{azulInst} &
\textbf{\scriptsize COGNITIVOS} &
\textbf{\scriptsize PROCEDIMENTALES} &
\textbf{\scriptsize ACTITUDINALES} \\
\hline
\endfirsthead
\hline
\rowcolor{azulInst}
\multirow{2}{*}{\color{white}\scriptsize\textbf{\makecell{ESTÁNDARES, LINEAMIENTOS\\U ORIENTACIONES}}} &
\multirow{2}{*}{\color{white}\scriptsize\textbf{\makecell{COMPETENCIAS}}} &
\multirow{2}{*}{\color{white}\scriptsize\textbf{\makecell{APRENDIZAJES}}} &
\multicolumn{3}{c|}{%
  \color{white}\makecell{\textbf{\footnotesize DESCRIPTORES DE DESEMPEÑO}\\[1pt]\textit{\scriptsize Metas del aprendizaje por periodo}}%
}\\
\cline{4-6}
\rowcolor{azulClaro}
\cellcolor{azulInst} & \cellcolor{azulInst} & \cellcolor{azulInst} &
\textbf{\scriptsize COGNITIVOS} &
\textbf{\scriptsize PROCEDIMENTALES} &
\textbf{\scriptsize ACTITUDINALES} \\
\hline
\endhead

% -- FILA 1: Naturaleza y evolución --
\textbf{Naturaleza y evolución de la tecnología} &
Comprender la lógica booleana como base del razonamiento computacional moderno. &
Identifica operadores lógicos AND, OR, NOT y su origen en el álgebra de Boole. &
Define los 3 operadores lógicos principales y construye tablas de verdad. &
Resuelve ejercicios de lógica booleana con tablas de verdad verificadas. &
Honra el gesto ancestral de cuidado y verificación cruzada de la abuela cocinera. \\
\hline

% -- FILA 2: Apropiación y uso --
\textbf{Apropiación y uso de la tecnología} &
Programar condicionales compuestas y estructuras anidadas en Scratch o entornos similares. &
Construye programas con condicionales múltiples y anidadas para resolver problemas reales. &
Conoce las estructuras: si-entonces-sino, anidamiento, casos múltiples. &
Programa una simulación de decisiones cotidianas (vestirse según clima, planear menú según ingredientes). &
Acepta la complejidad lógica de la vida cotidiana sin simplificarla. \\
\hline

% -- FILA 3: Solución de problemas --
\textbf{Solución de problemas con tecnología} &
Aplicar lógica condicional compuesta para resolver problemas de la vida real. &
Diseña algoritmos con condiciones AND/OR para evitar errores típicos. &
Identifica casos donde AND vs OR cambia el resultado del algoritmo. &
Crea un algoritmo verificador de decisiones cotidianas con condicionales compuestas. &
Verifica cruzando condiciones antes de actuar, como la abuela frente al fogón. \\
\hline

% -- FILA 4: Tecnología y sociedad --
\textbf{Tecnología y sociedad} &
Aplicar lógica compuesta para evitar errores socialmente costosos en sistemas automatizados. &
Reflexiona sobre cómo algoritmos mal diseñados pueden excluir o perjudicar a grupos. &
Reconoce ejemplos de sistemas automatizados con consecuencias sociales (filtros bancarios, sistemas de admisión). &
Analiza un caso real de decisión algorítmica que tuvo impacto social y propone mejoras. &
Asume responsabilidad por los efectos de los algoritmos que diseña en otros. \\
\hline

% -- FILA 5: PC (verde) --
\rowcolor{verdePCclaro}
\makecell[l]{\textbf{\color{verdePC}Pensamiento}\\\textbf{\color{verdePC}Computacional}\\
\footnotesize Wing (2006)\\\footnotesize MEN Or. 2022\\\footnotesize\textit{Diseño de algoritmos · Abstracción · Depuración}}\
 &
Aplicar operadores lógicos booleanos (AND, OR, NOT) como fundamento del pensamiento computacional y razonamiento algorítmico avanzado. &
Construye tablas de verdad (S1–S3); diseña algoritmos con condicionales múltiples y anidados (S4–S7); depura programas lógicos (S8–S9). &
Define álgebra de Boole y sus 3 operadores; reconoce que todo procesamiento digital se reduce a operaciones lógicas binarias; abstrae problemas cotidianos como expresiones booleanas. &
Diseña algoritmos con condiciones AND/OR; crea una simulación de decisiones con lógica booleana verificable y traza la tabla de verdad. &
Verifica cruzando condiciones antes de decidir; asume responsabilidad por los efectos de los algoritmos que diseña. \\
\hline

% -- FILA 6: SEL (violeta) --
\rowcolor{moradoSELclaro}
\makecell[l]{\textbf{\color{moradoSEL}Habilidades}\\\textbf{\color{moradoSEL}Socioemocionales}\\
\footnotesize Ley 2383/2024\\\footnotesize CASEL 2020\\\footnotesize\textit{Autogestión · Conciencia social}}\
 &
Desarrollar regulación emocional frente a la frustración de la lógica compleja y conciencia social frente a los efectos de decisiones algorítmicas en comunidades reales. &
Reconoce la frustración que genera la lógica compleja como señal de aprendizaje y la regula con estrategias concretas; identifica el impacto humano de los errores algorítmicos. &
Distingue entre error lógico personal y consecuencia social de un algoritmo defectuoso; describe casos de impacto real. &
Aplica la técnica de descomposición del problema lógico cuando se bloquea; analiza un caso de algoritmo con impacto social negativo y propone mejoras. &
Regula su frustración con paciencia; asume responsabilidad social por los algoritmos que diseña, más allá de la corrección técnica. \\
\hline

\end{longtable}
\medskip
\noindent\textbf{Nota:} Las filas 5 (\textcolor{verdePC}{verde}) y 6 (\textcolor{moradoSEL}{violeta}) son componentes propios de este plan, adicionales a los cuatro componentes de la Guía No.~30 del MEN.

\subsubsection{Periodo 2 · Datos como tablas: pensar con estructura}
\noindent\begin{tabular}{ll}
\textbf{Grado:} Octavo & \textbf{Periodo:} Segundo \\
\textbf{Proyecto integrador:} Datos del barrio: ver lo que no se ve a simple vista & \textbf{Anclaje:} Decisiones del campesino en la finca \\
\end{tabular}
\medskip

\begin{longtable}{|L{3.8cm}|L{3.3cm}|L{3.3cm}|L{2.5cm}|L{2.5cm}|L{2.5cm}|}
\caption{Estructura de malla curricular — G8° Periodo 2 · Datos como tablas: pensar con estructura}
\label{tab:g8p2}\\
\hline
\rowcolor{azulInst}
\multirow{2}{*}{\color{white}\scriptsize\textbf{\makecell{ESTÁNDARES, LINEAMIENTOS\\U ORIENTACIONES}}} &
\multirow{2}{*}{\color{white}\scriptsize\textbf{\makecell{COMPETENCIAS}}} &
\multirow{2}{*}{\color{white}\scriptsize\textbf{\makecell{APRENDIZAJES}}} &
\multicolumn{3}{c|}{%
  \color{white}\makecell{\textbf{\footnotesize DESCRIPTORES DE DESEMPEÑO}\\[1pt]\textit{\scriptsize Metas del aprendizaje por periodo}}%
}\\
\cline{4-6}
\rowcolor{azulClaro}
\cellcolor{azulInst} & \cellcolor{azulInst} & \cellcolor{azulInst} &
\textbf{\scriptsize COGNITIVOS} &
\textbf{\scriptsize PROCEDIMENTALES} &
\textbf{\scriptsize ACTITUDINALES} \\
\hline
\endfirsthead
\hline
\rowcolor{azulInst}
\multirow{2}{*}{\color{white}\scriptsize\textbf{\makecell{ESTÁNDARES, LINEAMIENTOS\\U ORIENTACIONES}}} &
\multirow{2}{*}{\color{white}\scriptsize\textbf{\makecell{COMPETENCIAS}}} &
\multirow{2}{*}{\color{white}\scriptsize\textbf{\makecell{APRENDIZAJES}}} &
\multicolumn{3}{c|}{%
  \color{white}\makecell{\textbf{\footnotesize DESCRIPTORES DE DESEMPEÑO}\\[1pt]\textit{\scriptsize Metas del aprendizaje por periodo}}%
}\\
\cline{4-6}
\rowcolor{azulClaro}
\cellcolor{azulInst} & \cellcolor{azulInst} & \cellcolor{azulInst} &
\textbf{\scriptsize COGNITIVOS} &
\textbf{\scriptsize PROCEDIMENTALES} &
\textbf{\scriptsize ACTITUDINALES} \\
\hline
\endhead

% -- FILA 1: Naturaleza y evolución --
\textbf{Naturaleza y evolución de la tecnología} &
Comprender el dato como representación digital del mundo. &
Identifica tipos de datos: numérico, texto, fecha, booleano. &
Diferencia los tipos de datos y sus operaciones permitidas. &
Clasifica datos en categorías y los organiza en tablas. &
Reconoce que el campesino del Valle ya hacía análisis de datos sin Excel. \\
\hline

% -- FILA 2: Apropiación y uso --
\textbf{Apropiación y uso de la tecnología} &
Manejar hojas de cálculo (Excel/Google Sheets) con criterio profesional. &
Construye tablas con encabezados, fórmulas básicas (SUMA, PROMEDIO) y gráficos simples. &
Conoce las funciones básicas de una hoja de cálculo. &
Construye una tabla de 20+ filas con datos reales y aplica fórmulas y gráficos. &
Cuida la calidad del dato: completo, verificado, fechado. \\
\hline

% -- FILA 3: Solución de problemas --
\textbf{Solución de problemas con tecnología} &
Resolver problemas analizando datos estructurados con criterio estadístico básico. &
Calcula promedios, mínimos, máximos y elabora conclusiones a partir de datos. &
Identifica las estadísticas básicas y su interpretación. &
Analiza un conjunto de datos del barrio y obtiene conclusiones. &
Verifica los datos antes de concluir; evita generalizar desde muestras pequeñas. \\
\hline

% -- FILA 4: Tecnología y sociedad --
\textbf{Tecnología y sociedad} &
Reconocer el rol ético del dato en decisiones sociales y políticas. &
Reflexiona sobre cómo datos mal usados pueden discriminar o manipular. &
Reconoce sesgos comunes: muestra pequeña, no representativa, sesgo de confirmación. &
Analiza un caso de manipulación estadística y propone correcciones. &
Asume responsabilidad por los datos que produce y publica. \\
\hline

% -- FILA 5: PC (verde) --
\rowcolor{verdePCclaro}
\makecell[l]{\textbf{\color{verdePC}Pensamiento}\\\textbf{\color{verdePC}Computacional}\\
\footnotesize Wing (2006)\\\footnotesize MEN Or. 2022\\\footnotesize\textit{Abstracción · Reconocimiento de patrones · Descomposición}}\
 &
Aplicar la abstracción del PC para modelar información real como datos estructurados y reconocer patrones estadísticos en conjuntos de datos del entorno. &
Abstrae datos del barrio como registros en tabla (S1–S3); aplica fórmulas estadísticas como algoritmos de procesamiento (S4–S7); visualiza patrones (S8–S9). &
Describe los tipos de datos como abstracciones del mundo real; reconoce que la estructura de la tabla determina las operaciones algorítmicas posibles. &
Diseña el modelo de datos de su investigación de barrio: tablas con nombre de columna, tipo de dato y validación; verifica calidad antes de analizar. &
Cuida la calidad del dato como ética de la información; rechaza generalizar desde muestras pequeñas. \\
\hline

% -- FILA 6: SEL (violeta) --
\rowcolor{moradoSELclaro}
\makecell[l]{\textbf{\color{moradoSEL}Habilidades}\\\textbf{\color{moradoSEL}Socioemocionales}\\
\footnotesize Ley 2383/2024\\\footnotesize CASEL 2020\\\footnotesize\textit{Conciencia social · Toma de decisiones responsable}}\
 &
Reconocer la responsabilidad ética de quien maneja datos de personas reales de la comunidad y tomar decisiones responsables que respeten la dignidad y privacidad de los involucrados. &
Identifica los principios de privacidad y dignidad que deben guiar el manejo de datos de personas reales de su entorno. &
Describe la diferencia entre dato anónimo y dato personal; reconoce las implicaciones éticas del manejo irresponsable de datos comunitarios. &
Diseña su instrumento de recolección de datos con consentimiento implícito y anonimización antes de publicar resultados. &
Trata los datos de su comunidad con la misma dignidad con que trataría a las personas; no publica datos que puedan perjudicar a alguien. \\
\hline

\end{longtable}
\medskip
\noindent\textbf{Nota:} Las filas 5 (\textcolor{verdePC}{verde}) y 6 (\textcolor{moradoSEL}{violeta}) son componentes propios de este plan, adicionales a los cuatro componentes de la Guía No.~30 del MEN.

\subsubsection{Periodo 3 · Diseño visual y comunicación digital}
\noindent\begin{tabular}{ll}
\textbf{Grado:} Octavo & \textbf{Periodo:} Tercero \\
\textbf{Proyecto integrador:} Diseño visual con raíces: mi pieza editorial del Pacífico & \textbf{Anclaje:} Diseño visual del Pacífico colombiano \\
\end{tabular}
\medskip

\begin{longtable}{|L{3.8cm}|L{3.3cm}|L{3.3cm}|L{2.5cm}|L{2.5cm}|L{2.5cm}|}
\caption{Estructura de malla curricular — G8° Periodo 3 · Diseño visual y comunicación digital}
\label{tab:g8p3}\\
\hline
\rowcolor{azulInst}
\multirow{2}{*}{\color{white}\scriptsize\textbf{\makecell{ESTÁNDARES, LINEAMIENTOS\\U ORIENTACIONES}}} &
\multirow{2}{*}{\color{white}\scriptsize\textbf{\makecell{COMPETENCIAS}}} &
\multirow{2}{*}{\color{white}\scriptsize\textbf{\makecell{APRENDIZAJES}}} &
\multicolumn{3}{c|}{%
  \color{white}\makecell{\textbf{\footnotesize DESCRIPTORES DE DESEMPEÑO}\\[1pt]\textit{\scriptsize Metas del aprendizaje por periodo}}%
}\\
\cline{4-6}
\rowcolor{azulClaro}
\cellcolor{azulInst} & \cellcolor{azulInst} & \cellcolor{azulInst} &
\textbf{\scriptsize COGNITIVOS} &
\textbf{\scriptsize PROCEDIMENTALES} &
\textbf{\scriptsize ACTITUDINALES} \\
\hline
\endfirsthead
\hline
\rowcolor{azulInst}
\multirow{2}{*}{\color{white}\scriptsize\textbf{\makecell{ESTÁNDARES, LINEAMIENTOS\\U ORIENTACIONES}}} &
\multirow{2}{*}{\color{white}\scriptsize\textbf{\makecell{COMPETENCIAS}}} &
\multirow{2}{*}{\color{white}\scriptsize\textbf{\makecell{APRENDIZAJES}}} &
\multicolumn{3}{c|}{%
  \color{white}\makecell{\textbf{\footnotesize DESCRIPTORES DE DESEMPEÑO}\\[1pt]\textit{\scriptsize Metas del aprendizaje por periodo}}%
}\\
\cline{4-6}
\rowcolor{azulClaro}
\cellcolor{azulInst} & \cellcolor{azulInst} & \cellcolor{azulInst} &
\textbf{\scriptsize COGNITIVOS} &
\textbf{\scriptsize PROCEDIMENTALES} &
\textbf{\scriptsize ACTITUDINALES} \\
\hline
\endhead

% -- FILA 1: Naturaleza y evolución --
\textbf{Naturaleza y evolución de la tecnología} &
Reconocer la historia del diseño visual desde los petroglifos hasta la era digital. &
Identifica hitos del diseño gráfico y la influencia ancestral en la estética contemporánea. &
Conoce los principios de diseño: jerarquía, equilibrio, contraste, ritmo, unidad. &
Analiza piezas de diseño popular (almanaques, esquelas) y contemporáneas. &
Valora el diseño visual del Pacífico como tradición digna y vigente. \\
\hline

% -- FILA 2: Apropiación y uso --
\textbf{Apropiación y uso de la tecnología} &
Producir piezas visuales con herramientas digitales (Canva, Figma, Google Slides). &
Diseña piezas comunicativas (cartel, infografía, presentación) con criterio profesional. &
Identifica los principios de tipografía, color y layout. &
Produce una pieza visual de 1 página integrando texto, imagen e identidad visual. &
Cuida la accesibilidad: contrastes legibles, fuentes claras, jerarquía evidente. \\
\hline

% -- FILA 3: Solución de problemas --
\textbf{Solución de problemas con tecnología} &
Resolver problemas de comunicación visual con criterio ético y estético. &
Identifica cuándo un diseño confunde y cómo corregirlo. &
Reconoce errores comunes: exceso de fuentes, color saturado, jerarquía confusa. &
Revisa diseños propios y ajenos aplicando una lista de chequeo de 5 criterios. &
Acepta retroalimentación con humildad y mejora iterativamente. \\
\hline

% -- FILA 4: Tecnología y sociedad --
\textbf{Tecnología y sociedad} &
Producir piezas visuales que dialoguen con saberes y estéticas locales. &
Integra elementos del diseño visual del Pacífico en producciones propias. &
Reconoce los aportes estéticos de comunidades del Chocó, Cauca y Valle. &
Diseña una pieza visual final que integre al menos un elemento ancestral identificable. &
Honra a las comunidades del Pacífico citando sus aportes en cada pieza. \\
\hline

% -- FILA 5: PC (verde) --
\rowcolor{verdePCclaro}
\makecell[l]{\textbf{\color{verdePC}Pensamiento}\\\textbf{\color{verdePC}Computacional}\\
\footnotesize Wing (2006)\\\footnotesize MEN Or. 2022\\\footnotesize\textit{Descomposición · Reconocimiento de patrones · Evaluación}}\
 &
Aplicar la descomposición y reconocimiento de patrones del PC al análisis y producción de piezas visuales, identificando principios de diseño como patrones reutilizables. &
Descompone piezas de diseño en sus principios (jerarquía, contraste, ritmo, equilibrio, unidad) como unidades analizables (S1–S3); aplica patrones visuales identificados (S4–S8). &
Reconoce los principios de diseño como patrones reutilizables; describe la jerarquía visual como algoritmo de lectura. &
Analiza 3 piezas de diseño descomponiéndolas en sus patrones; produce su pieza visual con lista de chequeo sistemática de 5 criterios. &
Acepta la retroalimentación con humildad; mejora iterativamente; reconoce que el buen diseño sigue lógica observable. \\
\hline

% -- FILA 6: SEL (violeta) --
\rowcolor{moradoSELclaro}
\makecell[l]{\textbf{\color{moradoSEL}Habilidades}\\\textbf{\color{moradoSEL}Socioemocionales}\\
\footnotesize Ley 2383/2024\\\footnotesize CASEL 2020\\\footnotesize\textit{Autoconciencia · Habilidades relacionales}}\
 &
Desarrollar autoconciencia estética sobre el propio gusto y las raíces culturales, y habilidades relacionales para recibir y dar retroalimentación constructiva sobre el trabajo creativo. &
Distingue entre preferencia estética personal y criterio de diseño fundado; identifica cómo sus raíces culturales informan su gusto visual. &
Describe la diferencia entre opinión estética y criterio técnico; reconoce cómo la cultura del Pacífico aporta elementos visuales propios. &
Da retroalimentación de diseño usando los 5 criterios técnicos (no "no me gusta"); recibe retroalimentación sin defensividad y aplica mejoras. &
Honra las raíces culturales propias en sus decisiones estéticas; recibe y da retroalimentación con generosidad y criterio. \\
\hline

\end{longtable}
\medskip
\noindent\textbf{Nota:} Las filas 5 (\textcolor{verdePC}{verde}) y 6 (\textcolor{moradoSEL}{violeta}) son componentes propios de este plan, adicionales a los cuatro componentes de la Guía No.~30 del MEN.

\subsection{Grado Noveno}

\subsubsection{Periodo 1 · Técnica, ciencia y tecnología}
\noindent\begin{tabular}{ll}
\textbf{Grado:} Noveno & \textbf{Periodo:} Primero \\
\textbf{Proyecto integrador:} Manifiesto del técnico crítico & \textbf{Anclaje:} La técnica en la mano (azuela, hueso tallado) \\
\end{tabular}
\medskip

\begin{longtable}{|L{3.8cm}|L{3.3cm}|L{3.3cm}|L{2.5cm}|L{2.5cm}|L{2.5cm}|}
\caption{Estructura de malla curricular — G9° Periodo 1 · Técnica, ciencia y tecnología}
\label{tab:g9p1}\\
\hline
\rowcolor{azulInst}
\multirow{2}{*}{\color{white}\scriptsize\textbf{\makecell{ESTÁNDARES, LINEAMIENTOS\\U ORIENTACIONES}}} &
\multirow{2}{*}{\color{white}\scriptsize\textbf{\makecell{COMPETENCIAS}}} &
\multirow{2}{*}{\color{white}\scriptsize\textbf{\makecell{APRENDIZAJES}}} &
\multicolumn{3}{c|}{%
  \color{white}\makecell{\textbf{\footnotesize DESCRIPTORES DE DESEMPEÑO}\\[1pt]\textit{\scriptsize Metas del aprendizaje por periodo}}%
}\\
\cline{4-6}
\rowcolor{azulClaro}
\cellcolor{azulInst} & \cellcolor{azulInst} & \cellcolor{azulInst} &
\textbf{\scriptsize COGNITIVOS} &
\textbf{\scriptsize PROCEDIMENTALES} &
\textbf{\scriptsize ACTITUDINALES} \\
\hline
\endfirsthead
\hline
\rowcolor{azulInst}
\multirow{2}{*}{\color{white}\scriptsize\textbf{\makecell{ESTÁNDARES, LINEAMIENTOS\\U ORIENTACIONES}}} &
\multirow{2}{*}{\color{white}\scriptsize\textbf{\makecell{COMPETENCIAS}}} &
\multirow{2}{*}{\color{white}\scriptsize\textbf{\makecell{APRENDIZAJES}}} &
\multicolumn{3}{c|}{%
  \color{white}\makecell{\textbf{\footnotesize DESCRIPTORES DE DESEMPEÑO}\\[1pt]\textit{\scriptsize Metas del aprendizaje por periodo}}%
}\\
\cline{4-6}
\rowcolor{azulClaro}
\cellcolor{azulInst} & \cellcolor{azulInst} & \cellcolor{azulInst} &
\textbf{\scriptsize COGNITIVOS} &
\textbf{\scriptsize PROCEDIMENTALES} &
\textbf{\scriptsize ACTITUDINALES} \\
\hline
\endhead

% -- FILA 1: Naturaleza y evolución --
\textbf{Naturaleza y evolución de la tecnología} &
Comprender la evolución histórica de la técnica desde los primeros artefactos hasta el chip. &
Construye una línea de tiempo de la técnica desde el fuego hasta la era digital. &
Distingue técnica, ciencia y tecnología; reconoce hitos: fuego, rueda, imprenta, electricidad, internet. &
Investiga 5 hitos tecnológicos y los presenta en línea de tiempo visual. &
Reconoce que la técnica nació en la mano humana antes que en la fábrica. \\
\hline

% -- FILA 2: Apropiación y uso --
\textbf{Apropiación y uso de la tecnología} &
Usar herramientas tecnológicas con criterio histórico y crítico. &
Maneja herramientas digitales valorando lo que cada una tiene de oficios anteriores. &
Conoce el origen y propósito de cada herramienta que usa cotidianamente. &
Documenta el uso de 5 herramientas digitales con su contexto histórico. &
Cuida las herramientas como extensión del cuerpo, no como descartables. \\
\hline

% -- FILA 3: Solución de problemas --
\textbf{Solución de problemas con tecnología} &
Resolver problemas tecnológicos con perspectiva histórica y crítica. &
Compara soluciones técnicas antiguas y modernas al mismo problema. &
Identifica problemas que la humanidad ha resuelto varias veces de formas distintas. &
Analiza un problema actual y propone solución integrando tradición y tecnología. &
No idolatra lo nuevo ni desprecia lo antiguo; juzga por criterio. \\
\hline

% -- FILA 4: Tecnología y sociedad --
\textbf{Tecnología y sociedad} &
Construir manifiesto del técnico crítico responsable. &
Redacta un manifiesto personal sobre su rol como técnico/usuario responsable. &
Reconoce las dimensiones éticas, ambientales y sociales del trabajo técnico. &
Elabora un manifiesto de 1 página con sus principios técnicos personales. &
Asume posición crítica frente al consumismo tecnológico. \\
\hline

% -- FILA 5: PC (verde) --
\rowcolor{verdePCclaro}
\makecell[l]{\textbf{\color{verdePC}Pensamiento}\\\textbf{\color{verdePC}Computacional}\\
\footnotesize Wing (2006)\\\footnotesize MEN Or. 2022\\\footnotesize\textit{Reconocimiento de patrones · Modelado · Evaluación}}\
 &
Aplicar el modelado del PC para comparar soluciones técnicas de distintas épocas como algoritmos alternativos de resolución de problemas humanos. &
Construye línea de tiempo tecnológica como modelo comparativo (S1–S3); compara soluciones técnicas como alternativas algorítmicas al mismo problema humano. &
Reconoce que la técnica es un proceso sistemático de pasos para transformar materia; comprende la evolución tecnológica como refinamiento iterativo de algoritmos de solución. &
Analiza un problema tecnológico actual y propone 2 soluciones (tradicional y digital) como algoritmos comparables con criterios explícitos de eficiencia, costo y contexto. &
No idolatra lo nuevo ni desprecia lo antiguo; evalúa con criterio algorítmico comparativo; reconoce la dignidad de la técnica ancestral. \\
\hline

% -- FILA 6: SEL (violeta) --
\rowcolor{moradoSELclaro}
\makecell[l]{\textbf{\color{moradoSEL}Habilidades}\\\textbf{\color{moradoSEL}Socioemocionales}\\
\footnotesize Ley 2383/2024\\\footnotesize CASEL 2020\\\footnotesize\textit{Autoconciencia · Toma de decisiones responsable}}\
 &
Desarrollar autoconciencia sobre el propio rol como técnico o usuario de tecnología, y tomar decisiones responsables frente al consumismo tecnológico y el impacto ambiental de los artefactos. &
Identifica sus propias actitudes frente a lo nuevo (fascinación, rechazo, criterio) y las examina con honestidad. &
Describe el impacto ambiental del ciclo de vida de los dispositivos tecnológicos y las alternativas de consumo responsable. &
Redacta un manifiesto personal de 10 principios como técnico responsable, incluyendo al menos 2 principios sobre impacto ambiental. &
Asume posición crítica frente al consumismo tecnológico; aplica sus principios en decisiones concretas de uso y descarte. \\
\hline

\end{longtable}
\medskip
\noindent\textbf{Nota:} Las filas 5 (\textcolor{verdePC}{verde}) y 6 (\textcolor{moradoSEL}{violeta}) son componentes propios de este plan, adicionales a los cuatro componentes de la Guía No.~30 del MEN.

\subsubsection{Periodo 2 · Diseño editorial y maquetación digital}
\noindent\begin{tabular}{ll}
\textbf{Grado:} Noveno & \textbf{Periodo:} Segundo \\
\textbf{Proyecto integrador:} Revista escolar con diseño editorial & \textbf{Anclaje:} Pasquines, almanaques y esquelas del Valle \\
\end{tabular}
\medskip

\begin{longtable}{|L{3.8cm}|L{3.3cm}|L{3.3cm}|L{2.5cm}|L{2.5cm}|L{2.5cm}|}
\caption{Estructura de malla curricular — G9° Periodo 2 · Diseño editorial y maquetación digital}
\label{tab:g9p2}\\
\hline
\rowcolor{azulInst}
\multirow{2}{*}{\color{white}\scriptsize\textbf{\makecell{ESTÁNDARES, LINEAMIENTOS\\U ORIENTACIONES}}} &
\multirow{2}{*}{\color{white}\scriptsize\textbf{\makecell{COMPETENCIAS}}} &
\multirow{2}{*}{\color{white}\scriptsize\textbf{\makecell{APRENDIZAJES}}} &
\multicolumn{3}{c|}{%
  \color{white}\makecell{\textbf{\footnotesize DESCRIPTORES DE DESEMPEÑO}\\[1pt]\textit{\scriptsize Metas del aprendizaje por periodo}}%
}\\
\cline{4-6}
\rowcolor{azulClaro}
\cellcolor{azulInst} & \cellcolor{azulInst} & \cellcolor{azulInst} &
\textbf{\scriptsize COGNITIVOS} &
\textbf{\scriptsize PROCEDIMENTALES} &
\textbf{\scriptsize ACTITUDINALES} \\
\hline
\endfirsthead
\hline
\rowcolor{azulInst}
\multirow{2}{*}{\color{white}\scriptsize\textbf{\makecell{ESTÁNDARES, LINEAMIENTOS\\U ORIENTACIONES}}} &
\multirow{2}{*}{\color{white}\scriptsize\textbf{\makecell{COMPETENCIAS}}} &
\multirow{2}{*}{\color{white}\scriptsize\textbf{\makecell{APRENDIZAJES}}} &
\multicolumn{3}{c|}{%
  \color{white}\makecell{\textbf{\footnotesize DESCRIPTORES DE DESEMPEÑO}\\[1pt]\textit{\scriptsize Metas del aprendizaje por periodo}}%
}\\
\cline{4-6}
\rowcolor{azulClaro}
\cellcolor{azulInst} & \cellcolor{azulInst} & \cellcolor{azulInst} &
\textbf{\scriptsize COGNITIVOS} &
\textbf{\scriptsize PROCEDIMENTALES} &
\textbf{\scriptsize ACTITUDINALES} \\
\hline
\endhead

% -- FILA 1: Naturaleza y evolución --
\textbf{Naturaleza y evolución de la tecnología} &
Comprender la evolución del diseño editorial desde la imprenta de Gutenberg hasta el diseño web. &
Identifica hitos del diseño editorial y su raíz popular en pasquines y almanaques. &
Conoce los componentes editoriales: portada, contraportada, sumario, secciones, cierre. &
Investiga 3 revistas y describe su estructura editorial. &
Honra el diseño popular del Valle como tradición editorial vigente. \\
\hline

% -- FILA 2: Apropiación y uso --
\textbf{Apropiación y uso de la tecnología} &
Maquetar piezas editoriales con herramientas digitales (Canva, Figma, Adobe Express). &
Produce una revista de 8–12 páginas con maquetación coherente. &
Reconoce reglas de cuadrícula, tipografía editorial y jerarquía visual. &
Maqueta una revista escolar de 8+ páginas con secciones, imágenes y texto. &
Cuida cada detalle como el editor de tradición. \\
\hline

% -- FILA 3: Solución de problemas --
\textbf{Solución de problemas con tecnología} &
Resolver problemas de maquetación con criterio profesional. &
Detecta y corrige problemas editoriales (jerarquía confusa, fuentes mal combinadas, accesibilidad). &
Conoce los 10 errores comunes en diseño editorial. &
Revisa la revista propia y ajena aplicando lista de chequeo. &
Itera con paciencia; el primer borrador nunca es el final. \\
\hline

% -- FILA 4: Tecnología y sociedad --
\textbf{Tecnología y sociedad} &
Publicar piezas editoriales que aporten valor a la comunidad escolar. &
Distribuye su revista a la comunidad y registra retroalimentación. &
Reconoce el rol social de la prensa escolar y su capacidad de informar y formar. &
Sustenta la revista ante el salón y recoge feedback de mínimo 5 lectores. &
Asume responsabilidad por lo que publica como editor. \\
\hline

% -- FILA 5: PC (verde) --
\rowcolor{verdePCclaro}
\makecell[l]{\textbf{\color{verdePC}Pensamiento}\\\textbf{\color{verdePC}Computacional}\\
\footnotesize Wing (2006)\\\footnotesize MEN Or. 2022\\\footnotesize\textit{Descomposición · Abstracción · Depuración iterativa}}\
 &
Aplicar la descomposición y abstracción del PC para analizar la estructura editorial como algoritmo de comunicación con reglas, módulos y componentes interdependientes. &
Descompone la estructura de una revista (portada, sumario, secciones) como módulos interdependientes (S1–S3); aplica reglas de cuadrícula como algoritmo de layout (S4–S7). &
Describe la maquetación como proceso algorítmico con reglas tipográficas, de cuadrícula y de jerarquía. &
Diseña la arquitectura de su revista con esquema de módulos, reglas tipográficas y sistema de cuadrícula documentados; depura con checklist. &
Itera con paciencia; el primer borrador nunca es el final; aplica depuración editorial como el programador depura su código. \\
\hline

% -- FILA 6: SEL (violeta) --
\rowcolor{moradoSELclaro}
\makecell[l]{\textbf{\color{moradoSEL}Habilidades}\\\textbf{\color{moradoSEL}Socioemocionales}\\
\footnotesize Ley 2383/2024\\\footnotesize CASEL 2020\\\footnotesize\textit{Autogestión · Habilidades relacionales}}\
 &
Desarrollar perseverancia en el proceso iterativo de diseño editorial y habilidades relacionales para incorporar retroalimentación de lectores y colaborar en la producción de la revista escolar. &
Identifica sus patrones de respuesta ante el error de diseño (abandono, perfeccionismo, iteración) y desarrolla la iteración como hábito. &
Describe la retroalimentación del lector como dato de diseño, no como crítica personal. &
Incorpora mínimo 3 sugerencias de lectores reales en la versión final de la revista; documenta qué cambió y por qué. &
Recibe feedback de lectores con apertura y gratitud; itera la revista con disposición colaborativa. \\
\hline

\end{longtable}
\medskip
\noindent\textbf{Nota:} Las filas 5 (\textcolor{verdePC}{verde}) y 6 (\textcolor{moradoSEL}{violeta}) son componentes propios de este plan, adicionales a los cuatro componentes de la Guía No.~30 del MEN.

\subsubsection{Periodo 3 · Datos, estadística y visualización}
\noindent\begin{tabular}{ll}
\textbf{Grado:} Noveno & \textbf{Periodo:} Tercero \\
\textbf{Proyecto integrador:} Diagnóstico de mi colegio con datos & \textbf{Anclaje:} La libreta de la abuela (conteo de cosecha) \\
\end{tabular}
\medskip

\begin{longtable}{|L{3.8cm}|L{3.3cm}|L{3.3cm}|L{2.5cm}|L{2.5cm}|L{2.5cm}|}
\caption{Estructura de malla curricular — G9° Periodo 3 · Datos, estadística y visualización}
\label{tab:g9p3}\\
\hline
\rowcolor{azulInst}
\multirow{2}{*}{\color{white}\scriptsize\textbf{\makecell{ESTÁNDARES, LINEAMIENTOS\\U ORIENTACIONES}}} &
\multirow{2}{*}{\color{white}\scriptsize\textbf{\makecell{COMPETENCIAS}}} &
\multirow{2}{*}{\color{white}\scriptsize\textbf{\makecell{APRENDIZAJES}}} &
\multicolumn{3}{c|}{%
  \color{white}\makecell{\textbf{\footnotesize DESCRIPTORES DE DESEMPEÑO}\\[1pt]\textit{\scriptsize Metas del aprendizaje por periodo}}%
}\\
\cline{4-6}
\rowcolor{azulClaro}
\cellcolor{azulInst} & \cellcolor{azulInst} & \cellcolor{azulInst} &
\textbf{\scriptsize COGNITIVOS} &
\textbf{\scriptsize PROCEDIMENTALES} &
\textbf{\scriptsize ACTITUDINALES} \\
\hline
\endfirsthead
\hline
\rowcolor{azulInst}
\multirow{2}{*}{\color{white}\scriptsize\textbf{\makecell{ESTÁNDARES, LINEAMIENTOS\\U ORIENTACIONES}}} &
\multirow{2}{*}{\color{white}\scriptsize\textbf{\makecell{COMPETENCIAS}}} &
\multirow{2}{*}{\color{white}\scriptsize\textbf{\makecell{APRENDIZAJES}}} &
\multicolumn{3}{c|}{%
  \color{white}\makecell{\textbf{\footnotesize DESCRIPTORES DE DESEMPEÑO}\\[1pt]\textit{\scriptsize Metas del aprendizaje por periodo}}%
}\\
\cline{4-6}
\rowcolor{azulClaro}
\cellcolor{azulInst} & \cellcolor{azulInst} & \cellcolor{azulInst} &
\textbf{\scriptsize COGNITIVOS} &
\textbf{\scriptsize PROCEDIMENTALES} &
\textbf{\scriptsize ACTITUDINALES} \\
\hline
\endhead

% -- FILA 1: Naturaleza y evolución --
\textbf{Naturaleza y evolución de la tecnología} &
Comprender la estadística como ciencia de la incertidumbre y la visualización como lenguaje universal. &
Identifica los tipos de gráficos y cuándo usar cada uno. &
Define media, mediana, moda, rango, desviación; distingue gráficas de barras, líneas, sectores. &
Calcula medidas de tendencia central a partir de un conjunto de datos. &
Honra a la abuela cosechera como estadista ancestral. \\
\hline

% -- FILA 2: Apropiación y uso --
\textbf{Apropiación y uso de la tecnología} &
Producir visualizaciones honestas con hojas de cálculo y herramientas de BI básico. &
Diseña gráficos en Excel/Sheets con criterio comunicativo. &
Conoce los principios de la visualización honesta: ejes correctos, escalas justas, sin distorsión. &
Produce 3 visualizaciones distintas de un mismo dataset escolar. &
Diseña gráficos para informar, no para impresionar ni engañar. \\
\hline

% -- FILA 3: Solución de problemas --
\textbf{Solución de problemas con tecnología} &
Resolver problemas comunitarios con análisis de datos. &
Identifica un problema del colegio, recolecta datos y propone solución basada en evidencia. &
Reconoce los pasos del método de investigación con datos. &
Realiza encuesta a mínimo 30 estudiantes, analiza resultados y presenta hallazgos. &
Aplica rigor estadístico básico; reconoce los límites de su muestra. \\
\hline

% -- FILA 4: Tecnología y sociedad --
\textbf{Tecnología y sociedad} &
Comunicar resultados de análisis de datos a la comunidad escolar con honestidad. &
Presenta el diagnóstico del colegio a directivos y comunidad. &
Reconoce el rol del dato en la toma de decisiones democráticas. &
Elabora informe ciudadano de 6 páginas con datos, gráficos y propuestas. &
Honra la estadística doméstica de la abuela como ética de cuidado comunitario. \\
\hline

% -- FILA 5: PC (verde) --
\rowcolor{verdePCclaro}
\makecell[l]{\textbf{\color{verdePC}Pensamiento}\\\textbf{\color{verdePC}Computacional}\\
\footnotesize Wing (2006)\\\footnotesize MEN Or. 2022\\\footnotesize\textit{Algoritmos estadísticos · Abstracción · Evaluación}}\
 &
Aplicar algoritmos estadísticos y principios de abstracción del PC para analizar datos del entorno y comunicar hallazgos con honestidad y rigor. &
Aplica algoritmos estadísticos (promedio, mediana, moda, desviación) como funciones sobre conjuntos de datos (S1–S5); diseña visualizaciones según el patrón estadístico adecuado. &
Reconoce las medidas de tendencia central como algoritmos de síntesis de datos; identifica sesgos comunes como errores de abstracción (muestra no representativa, escala engañosa). &
Recolecta datos de mínimo 30 estudiantes, aplica los algoritmos estadísticos y diseña 3 visualizaciones distintas del mismo dataset; valida honestidad de sus gráficos. &
Aplica rigor estadístico; reconoce los límites de su muestra; asume responsabilidad por los datos que publica y la interpretación que propone. \\
\hline

% -- FILA 6: SEL (violeta) --
\rowcolor{moradoSELclaro}
\makecell[l]{\textbf{\color{moradoSEL}Habilidades}\\\textbf{\color{moradoSEL}Socioemocionales}\\
\footnotesize Ley 2383/2024\\\footnotesize CASEL 2020\\\footnotesize\textit{Conciencia social · Toma de decisiones responsable}}\
 &
Desarrollar conciencia social sobre el impacto que tienen los datos y las estadísticas en la comunidad escolar, y tomar decisiones responsables sobre cómo comunicar hallazgos sin manipular ni distorsionar. &
Identifica cómo una visualización deshonesta puede generar alarma, estigma o decisiones erróneas en la comunidad. &
Describe al menos 3 técnicas de distorsión estadística (escala truncada, muestra sesgada, correlación-causalidad) y sus consecuencias sociales. &
Revisa sus visualizaciones con un compañero usando la lista de chequeo de honestidad estadística antes de presentar a la comunidad. &
Asume la comunicación de datos como acto político y ético; presenta hallazgos con rigor, humildad y respeto por su comunidad. \\
\hline

\end{longtable}
\medskip
\noindent\textbf{Nota:} Las filas 5 (\textcolor{verdePC}{verde}) y 6 (\textcolor{moradoSEL}{violeta}) son componentes propios de este plan, adicionales a los cuatro componentes de la Guía No.~30 del MEN.

\subsection{Grado Décimo}

\subsubsection{Periodo 1 · Edición digital y autoría con IA responsable}
\noindent\begin{tabular}{ll}
\textbf{Grado:} Décimo & \textbf{Periodo:} Primero \\
\textbf{Proyecto integrador:} Mi libro digital: del oficio del editor al sello propio & \textbf{Anclaje:} El personaje silencioso del barrio (intermediario) \\
\end{tabular}
\medskip

\begin{longtable}{|L{3.8cm}|L{3.3cm}|L{3.3cm}|L{2.5cm}|L{2.5cm}|L{2.5cm}|}
\caption{Estructura de malla curricular — G10° Periodo 1 · Edición digital y autoría con IA responsable}
\label{tab:g10p1}\\
\hline
\rowcolor{azulInst}
\multirow{2}{*}{\color{white}\scriptsize\textbf{\makecell{ESTÁNDARES, LINEAMIENTOS\\U ORIENTACIONES}}} &
\multirow{2}{*}{\color{white}\scriptsize\textbf{\makecell{COMPETENCIAS}}} &
\multirow{2}{*}{\color{white}\scriptsize\textbf{\makecell{APRENDIZAJES}}} &
\multicolumn{3}{c|}{%
  \color{white}\makecell{\textbf{\footnotesize DESCRIPTORES DE DESEMPEÑO}\\[1pt]\textit{\scriptsize Metas del aprendizaje por periodo}}%
}\\
\cline{4-6}
\rowcolor{azulClaro}
\cellcolor{azulInst} & \cellcolor{azulInst} & \cellcolor{azulInst} &
\textbf{\scriptsize COGNITIVOS} &
\textbf{\scriptsize PROCEDIMENTALES} &
\textbf{\scriptsize ACTITUDINALES} \\
\hline
\endfirsthead
\hline
\rowcolor{azulInst}
\multirow{2}{*}{\color{white}\scriptsize\textbf{\makecell{ESTÁNDARES, LINEAMIENTOS\\U ORIENTACIONES}}} &
\multirow{2}{*}{\color{white}\scriptsize\textbf{\makecell{COMPETENCIAS}}} &
\multirow{2}{*}{\color{white}\scriptsize\textbf{\makecell{APRENDIZAJES}}} &
\multicolumn{3}{c|}{%
  \color{white}\makecell{\textbf{\footnotesize DESCRIPTORES DE DESEMPEÑO}\\[1pt]\textit{\scriptsize Metas del aprendizaje por periodo}}%
}\\
\cline{4-6}
\rowcolor{azulClaro}
\cellcolor{azulInst} & \cellcolor{azulInst} & \cellcolor{azulInst} &
\textbf{\scriptsize COGNITIVOS} &
\textbf{\scriptsize PROCEDIMENTALES} &
\textbf{\scriptsize ACTITUDINALES} \\
\hline
\endhead

% -- FILA 1: Naturaleza y evolución --
\textbf{Naturaleza y evolución de la tecnología} &
Comprender la evolución de la autoría editorial desde la imprenta hasta los libros con IA. &
Distingue el rol histórico del editor y los desafíos contemporáneos con IA generativa. &
Define qué hace y qué no hace un editor; reconoce el ciclo editorial completo. &
Analiza 3 libros editoriales y describe su estructura, voz y producción. &
Honra al personaje silencioso del barrio como intermediario cultural ancestral. \\
\hline

% -- FILA 2: Apropiación y uso --
\textbf{Apropiación y uso de la tecnología} &
Usar IA con criterio profesional para escritura editorial. &
Aplica las 5 partes del prompt profesional para texto y lo itera. &
Conoce las técnicas de prompting técnico para escritura: rol, contexto, restricciones, ejemplos, formato. &
Produce capítulos de libro con IA como asistente, editando con voz propia. &
La IA es asistente, no autor; la voz del libro es la del estudiante. \\
\hline

% -- FILA 3: Solución de problemas --
\textbf{Solución de problemas con tecnología} &
Resolver problemas de derechos de autor y autoría con IA según legislación colombiana. &
Aplica la Ley 23 de 1982 y las consideraciones éticas sobre obras derivadas con IA. &
Reconoce la diferencia entre obra original, obra derivada y obra generada por IA. &
Declara explícitamente la participación de IA en cada capítulo del libro. &
Sostiene la honestidad académica frente a la presión del atajo. \\
\hline

% -- FILA 4: Tecnología y sociedad --
\textbf{Tecnología y sociedad} &
Publicar un libro escolar con criterio profesional y compromiso ético. &
Compila, sustenta y difunde su libro de 80 páginas en la feria del libro escolar (S10). &
Reconoce el rol del autor responsable en la infoesfera contemporánea. &
Sustenta su libro públicamente y recoge retroalimentación crítica. &
Asume el rol de autor con responsabilidad y orgullo. \\
\hline

% -- FILA 5: PC (verde) --
\rowcolor{verdePCclaro}
\makecell[l]{\textbf{\color{verdePC}Pensamiento}\\\textbf{\color{verdePC}Computacional}\\
\footnotesize Wing (2006)\\\footnotesize MEN Or. 2022\\\footnotesize\textit{Diseño de algoritmos (prompting) · Depuración iterativa}}\
 &
Aplicar el PC para diseñar flujos de trabajo editoriales asistidos por IA, reconociendo el prompting como programación en lenguaje natural con estructura algorítmica. &
Diseña prompts como instrucciones algorítmicas (rol, contexto, tarea, restricciones) (S3–S5); itera el flujo editorial como proceso de depuración sistemática (S6–S8). &
Reconoce el prompt como instrucción algorítmica en lenguaje natural; comprende el flujo editorial asistido por IA como proceso con etapas, condiciones y revisiones documentadas. &
Diseña su flujo editorial para el libro digital (borrador IA / edición propia / verificación / versión final) como diagrama de proceso con decisiones y declaraciones documentadas. &
Asume la IA como asistente; la voz del libro es la suya; declara el uso con honestidad y responsabilidad ante sus lectores. \\
\hline

% -- FILA 6: SEL (violeta) --
\rowcolor{moradoSELclaro}
\makecell[l]{\textbf{\color{moradoSEL}Habilidades}\\\textbf{\color{moradoSEL}Socioemocionales}\\
\footnotesize Ley 2383/2024\\\footnotesize CASEL 2020\\\footnotesize\textit{Autoconciencia · Toma de decisiones responsable}}\
 &
Desarrollar autoconciencia sobre la propia voz autoral y tomar decisiones responsables frente al uso de IA, manteniendo integridad académica y reconociendo los límites éticos de la autoría asistida. &
Distingue con claridad cuáles partes del libro son voz propia y cuáles son asistencia de IA; evalúa si la proporción es éticamente adecuada. &
Describe los límites éticos de la autoría asistida por IA según la legislación colombiana y los principios de honestidad académica. &
Revisa cada capítulo de su libro con la pregunta "¿esta voz es mía?" y documenta las decisiones editoriales que tomó. &
Defiende su voz autoral como algo valioso e irremplazable; usa la IA para amplificarla, nunca para sustituirla. \\
\hline

\end{longtable}
\medskip
\noindent\textbf{Nota:} Las filas 5 (\textcolor{verdePC}{verde}) y 6 (\textcolor{moradoSEL}{violeta}) son componentes propios de este plan, adicionales a los cuatro componentes de la Guía No.~30 del MEN.

\subsubsection{Periodo 2 · Comunicación profesional con IA}
\noindent\begin{tabular}{ll}
\textbf{Grado:} Décimo & \textbf{Periodo:} Segundo \\
\textbf{Proyecto integrador:} Informe profesional del oficio & \textbf{Anclaje:} Los 4 oficios de la palabra (barrio Obrero) \\
\end{tabular}
\medskip

\begin{longtable}{|L{3.8cm}|L{3.3cm}|L{3.3cm}|L{2.5cm}|L{2.5cm}|L{2.5cm}|}
\caption{Estructura de malla curricular — G10° Periodo 2 · Comunicación profesional con IA}
\label{tab:g10p2}\\
\hline
\rowcolor{azulInst}
\multirow{2}{*}{\color{white}\scriptsize\textbf{\makecell{ESTÁNDARES, LINEAMIENTOS\\U ORIENTACIONES}}} &
\multirow{2}{*}{\color{white}\scriptsize\textbf{\makecell{COMPETENCIAS}}} &
\multirow{2}{*}{\color{white}\scriptsize\textbf{\makecell{APRENDIZAJES}}} &
\multicolumn{3}{c|}{%
  \color{white}\makecell{\textbf{\footnotesize DESCRIPTORES DE DESEMPEÑO}\\[1pt]\textit{\scriptsize Metas del aprendizaje por periodo}}%
}\\
\cline{4-6}
\rowcolor{azulClaro}
\cellcolor{azulInst} & \cellcolor{azulInst} & \cellcolor{azulInst} &
\textbf{\scriptsize COGNITIVOS} &
\textbf{\scriptsize PROCEDIMENTALES} &
\textbf{\scriptsize ACTITUDINALES} \\
\hline
\endfirsthead
\hline
\rowcolor{azulInst}
\multirow{2}{*}{\color{white}\scriptsize\textbf{\makecell{ESTÁNDARES, LINEAMIENTOS\\U ORIENTACIONES}}} &
\multirow{2}{*}{\color{white}\scriptsize\textbf{\makecell{COMPETENCIAS}}} &
\multirow{2}{*}{\color{white}\scriptsize\textbf{\makecell{APRENDIZAJES}}} &
\multicolumn{3}{c|}{%
  \color{white}\makecell{\textbf{\footnotesize DESCRIPTORES DE DESEMPEÑO}\\[1pt]\textit{\scriptsize Metas del aprendizaje por periodo}}%
}\\
\cline{4-6}
\rowcolor{azulClaro}
\cellcolor{azulInst} & \cellcolor{azulInst} & \cellcolor{azulInst} &
\textbf{\scriptsize COGNITIVOS} &
\textbf{\scriptsize PROCEDIMENTALES} &
\textbf{\scriptsize ACTITUDINALES} \\
\hline
\endhead

% -- FILA 1: Naturaleza y evolución --
\textbf{Naturaleza y evolución de la tecnología} &
Comprender la evolución de los géneros de texto profesional y sus convenciones. &
Identifica los 4 tipos de texto profesional y sus convenciones. &
Distingue informe técnico, informe comercial, ensayo y estudio de mercado. &
Analiza 4 textos profesionales y los clasifica por género. &
Honra a los 4 oficios de la palabra del barrio Obrero. \\
\hline

% -- FILA 2: Apropiación y uso --
\textbf{Apropiación y uso de la tecnología} &
Producir documentos profesionales con Google Docs/Word y IA responsable. &
Maneja estilos, índices, citas y referencias en documentos profesionales. &
Reconoce los componentes de un documento profesional formal. &
Produce un informe técnico de 10+ páginas con todos los componentes formales. &
Cuida cada detalle como el cuentista cuida la palabra empeñada. \\
\hline

% -- FILA 3: Solución de problemas --
\textbf{Solución de problemas con tecnología} &
Resolver problemas comunicativos en contextos laborales reales. &
Redacta correos profesionales, propuestas y reportes para audiencias reales. &
Identifica el tono, registro y estructura adecuados según audiencia. &
Redacta y envía 3 correos profesionales reales (a docentes, directivos o empresas). &
Sostiene profesionalismo aún en correos cortos. \\
\hline

% -- FILA 4: Tecnología y sociedad --
\textbf{Tecnología y sociedad} &
Aplicar la palabra profesional como herramienta de cambio social y empleabilidad. &
Sustenta el oficio de la palabra como ciudadanía activa (S10). &
Reconoce el valor profesional y social de saber escribir con criterio. &
Sustenta su informe ante audiencia real con criterio y argumentos. &
Honra la palabra empeñada como los 4 oficios del barrio Obrero. \\
\hline

% -- FILA 5: PC (verde) --
\rowcolor{verdePCclaro}
\makecell[l]{\textbf{\color{verdePC}Pensamiento}\\\textbf{\color{verdePC}Computacional}\\
\footnotesize Wing (2006)\\\footnotesize MEN Or. 2022\\\footnotesize\textit{Abstracción · Reconocimiento de patrones}}\
 &
Aplicar la abstracción y reconocimiento de patrones del PC para identificar las estructuras de género textual profesional como plantillas reutilizables y adaptables. &
Abstrae los componentes de 4 géneros profesionales como plantillas con patrones reutilizables (S1–S4); aplica y adapta las plantillas con voz propia (S5–S8). &
Reconoce los géneros textuales como abstracciones de patrones comunicativos con estructura, tono y propósito predefinidos; los describe como plantillas algorítmicas adaptables. &
Analiza 4 textos profesionales identificando sus patrones estructurales; produce un informe técnico siguiendo la plantilla identificada y adaptándola con voz propia verificable. &
Honra la palabra empeñada con la precisión del cuentista del barrio Obrero; sostiene profesionalismo aun en documentos cortos. \\
\hline

% -- FILA 6: SEL (violeta) --
\rowcolor{moradoSELclaro}
\makecell[l]{\textbf{\color{moradoSEL}Habilidades}\\\textbf{\color{moradoSEL}Socioemocionales}\\
\footnotesize Ley 2383/2024\\\footnotesize CASEL 2020\\\footnotesize\textit{Habilidades relacionales · Autogestión}}\
 &
Desarrollar habilidades relacionales para comunicarse con criterio profesional en contextos laborales y académicos, y autogestión emocional para mantener el tono adecuado aún bajo presión. &
Identifica situaciones de comunicación profesional que generan tensión emocional (correo de reclamo, informe con errores, comunicación con directivos) y practica respuestas asertivas. &
Describe la asertividad como equilibrio entre firmeza y cortesía; reconoce los efectos de la comunicación agresiva o pasiva en contextos profesionales. &
Redacta 3 correos profesionales en situaciones de tensión (reclamar, corregir, solicitar) usando registro asertivo y tono respetuoso. &
Mantiene el profesionalismo aún cuando el mensaje es difícil; honra la palabra empeñada también en la adversidad. \\
\hline

\end{longtable}
\medskip
\noindent\textbf{Nota:} Las filas 5 (\textcolor{verdePC}{verde}) y 6 (\textcolor{moradoSEL}{violeta}) son componentes propios de este plan, adicionales a los cuatro componentes de la Guía No.~30 del MEN.

\subsubsection{Periodo 3 · Microemprendimiento y gestión digital}
\noindent\begin{tabular}{ll}
\textbf{Grado:} Décimo & \textbf{Periodo:} Tercero \\
\textbf{Proyecto integrador:} Microemprendimiento escolar con plan completo & \textbf{Anclaje:} Tiendas y panaderías del centro (Cartago) \\
\end{tabular}
\medskip

\begin{longtable}{|L{3.8cm}|L{3.3cm}|L{3.3cm}|L{2.5cm}|L{2.5cm}|L{2.5cm}|}
\caption{Estructura de malla curricular — G10° Periodo 3 · Microemprendimiento y gestión digital}
\label{tab:g10p3}\\
\hline
\rowcolor{azulInst}
\multirow{2}{*}{\color{white}\scriptsize\textbf{\makecell{ESTÁNDARES, LINEAMIENTOS\\U ORIENTACIONES}}} &
\multirow{2}{*}{\color{white}\scriptsize\textbf{\makecell{COMPETENCIAS}}} &
\multirow{2}{*}{\color{white}\scriptsize\textbf{\makecell{APRENDIZAJES}}} &
\multicolumn{3}{c|}{%
  \color{white}\makecell{\textbf{\footnotesize DESCRIPTORES DE DESEMPEÑO}\\[1pt]\textit{\scriptsize Metas del aprendizaje por periodo}}%
}\\
\cline{4-6}
\rowcolor{azulClaro}
\cellcolor{azulInst} & \cellcolor{azulInst} & \cellcolor{azulInst} &
\textbf{\scriptsize COGNITIVOS} &
\textbf{\scriptsize PROCEDIMENTALES} &
\textbf{\scriptsize ACTITUDINALES} \\
\hline
\endfirsthead
\hline
\rowcolor{azulInst}
\multirow{2}{*}{\color{white}\scriptsize\textbf{\makecell{ESTÁNDARES, LINEAMIENTOS\\U ORIENTACIONES}}} &
\multirow{2}{*}{\color{white}\scriptsize\textbf{\makecell{COMPETENCIAS}}} &
\multirow{2}{*}{\color{white}\scriptsize\textbf{\makecell{APRENDIZAJES}}} &
\multicolumn{3}{c|}{%
  \color{white}\makecell{\textbf{\footnotesize DESCRIPTORES DE DESEMPEÑO}\\[1pt]\textit{\scriptsize Metas del aprendizaje por periodo}}%
}\\
\cline{4-6}
\rowcolor{azulClaro}
\cellcolor{azulInst} & \cellcolor{azulInst} & \cellcolor{azulInst} &
\textbf{\scriptsize COGNITIVOS} &
\textbf{\scriptsize PROCEDIMENTALES} &
\textbf{\scriptsize ACTITUDINALES} \\
\hline
\endhead

% -- FILA 1: Naturaleza y evolución --
\textbf{Naturaleza y evolución de la tecnología} &
Comprender la evolución de la gestión empresarial desde la libreta hasta el ERP. &
Identifica las herramientas digitales clave para la gestión moderna. &
Distingue ofimática, contabilidad básica, gestión de inventarios y CRM. &
Analiza la gestión digital de una microempresa real. &
Honra a las tiendas y panaderías como microempresas ancestrales. \\
\hline

% -- FILA 2: Apropiación y uso --
\textbf{Apropiación y uso de la tecnología} &
Manejar hojas de cálculo avanzadas para contabilidad personal y empresarial. &
Aplica fórmulas avanzadas, formato condicional y gráficos dinámicos. &
Conoce funciones BUSCARV, SI, SUMAR.SI, gráficos dinámicos. &
Construye libro de cuentas simulado con flujo de caja y punto de equilibrio. &
Cuida la calidad del dato contable: completo, fechado, verificado. \\
\hline

% -- FILA 3: Solución de problemas --
\textbf{Solución de problemas con tecnología} &
Resolver problemas de microemprendimiento con plan de negocio y Business Model Canvas. &
Construye un plan de negocio con IA para microempresa escolar. &
Identifica los 9 componentes del Business Model Canvas. &
Elabora plan de negocio de 15 páginas y BMC visual para su proyecto. &
Aplica el saber del comerciante del centro: planear y mantener la palabra. \\
\hline

% -- FILA 4: Tecnología y sociedad --
\textbf{Tecnología y sociedad} &
Publicar el microemprendimiento como ciudadanía económica responsable. &
Sustenta en feria de microempresas escolares (S10). &
Reconoce el rol social del emprendimiento responsable. &
Sustenta ante jurado real (docentes + invitados) y registra mejoras. &
Honra al comerciante tradicional como referente ético. \\
\hline

% -- FILA 5: PC (verde) --
\rowcolor{verdePCclaro}
\makecell[l]{\textbf{\color{verdePC}Pensamiento}\\\textbf{\color{verdePC}Computacional}\\
\footnotesize Wing (2006)\\\footnotesize MEN Or. 2022\\\footnotesize\textit{Modelado · Abstracción · Algoritmos de gestión}}\
 &
Aplicar el modelado del PC para diseñar el sistema de un microemprendimiento usando el Business Model Canvas como abstracción funcional de 9 componentes interrelacionados. &
Modela su microempresa usando el BMC como abstracción de 9 componentes con relaciones y flujos (S3–S6); usa hoja de cálculo como sistema de procesamiento algorítmico de datos financieros. &
Reconoce el Business Model Canvas como modelo de abstracción con 9 componentes, relaciones y flujos interdependientes. &
Diseña su BMC digital con 9 componentes completos; crea un libro de cuentas con fórmulas de flujo de caja, inventario y punto de equilibrio verificado. &
Aplica el saber del comerciante del centro: planear y mantener la palabra; cuida la calidad del dato contable como ética empresarial. \\
\hline

% -- FILA 6: SEL (violeta) --
\rowcolor{moradoSELclaro}
\makecell[l]{\textbf{\color{moradoSEL}Habilidades}\\\textbf{\color{moradoSEL}Socioemocionales}\\
\footnotesize Ley 2383/2024\\\footnotesize CASEL 2020\\\footnotesize\textit{Conciencia social · Toma de decisiones responsable}}\
 &
Desarrollar conciencia del impacto social del microemprendimiento en la comunidad y tomar decisiones empresariales responsables que honren la dignidad de clientes, empleados y del entorno. &
Identifica a los grupos impactados por su microemprendimiento (clientes, proveedores, entorno) y evalúa el impacto social de sus decisiones empresariales. &
Describe la responsabilidad social empresarial como componente del emprendimiento responsable; reconoce la diferencia entre ganancia cortoplacista y sostenibilidad a largo plazo. &
Incluye en su plan de negocio una sección de impacto social y ambiental con al menos 3 compromisos concretos y verificables. &
Emprende con la ética del comerciante del barrio: palabra empeñada, trato justo y compromiso con la comunidad como legado. \\
\hline

\end{longtable}
\medskip
\noindent\textbf{Nota:} Las filas 5 (\textcolor{verdePC}{verde}) y 6 (\textcolor{moradoSEL}{violeta}) son componentes propios de este plan, adicionales a los cuatro componentes de la Guía No.~30 del MEN.

\subsection{Grado Undécimo}

\subsubsection{Periodo 1 · Presencia digital empresarial con IA}
\noindent\begin{tabular}{ll}
\textbf{Grado:} Undécimo & \textbf{Periodo:} Primero \\
\textbf{Proyecto integrador:} Presencia digital empresarial con IA & \textbf{Anclaje:} Oficio y palabra (identidad y reputación digital) \\
\end{tabular}
\medskip

\begin{longtable}{|L{3.8cm}|L{3.3cm}|L{3.3cm}|L{2.5cm}|L{2.5cm}|L{2.5cm}|}
\caption{Estructura de malla curricular — G11° Periodo 1 · Presencia digital empresarial con IA}
\label{tab:g11p1}\\
\hline
\rowcolor{azulInst}
\multirow{2}{*}{\color{white}\scriptsize\textbf{\makecell{ESTÁNDARES, LINEAMIENTOS\\U ORIENTACIONES}}} &
\multirow{2}{*}{\color{white}\scriptsize\textbf{\makecell{COMPETENCIAS}}} &
\multirow{2}{*}{\color{white}\scriptsize\textbf{\makecell{APRENDIZAJES}}} &
\multicolumn{3}{c|}{%
  \color{white}\makecell{\textbf{\footnotesize DESCRIPTORES DE DESEMPEÑO}\\[1pt]\textit{\scriptsize Metas del aprendizaje por periodo}}%
}\\
\cline{4-6}
\rowcolor{azulClaro}
\cellcolor{azulInst} & \cellcolor{azulInst} & \cellcolor{azulInst} &
\textbf{\scriptsize COGNITIVOS} &
\textbf{\scriptsize PROCEDIMENTALES} &
\textbf{\scriptsize ACTITUDINALES} \\
\hline
\endfirsthead
\hline
\rowcolor{azulInst}
\multirow{2}{*}{\color{white}\scriptsize\textbf{\makecell{ESTÁNDARES, LINEAMIENTOS\\U ORIENTACIONES}}} &
\multirow{2}{*}{\color{white}\scriptsize\textbf{\makecell{COMPETENCIAS}}} &
\multirow{2}{*}{\color{white}\scriptsize\textbf{\makecell{APRENDIZAJES}}} &
\multicolumn{3}{c|}{%
  \color{white}\makecell{\textbf{\footnotesize DESCRIPTORES DE DESEMPEÑO}\\[1pt]\textit{\scriptsize Metas del aprendizaje por periodo}}%
}\\
\cline{4-6}
\rowcolor{azulClaro}
\cellcolor{azulInst} & \cellcolor{azulInst} & \cellcolor{azulInst} &
\textbf{\scriptsize COGNITIVOS} &
\textbf{\scriptsize PROCEDIMENTALES} &
\textbf{\scriptsize ACTITUDINALES} \\
\hline
\endhead

% -- FILA 1: Naturaleza y evolución --
\textbf{Naturaleza y evolución de la tecnología} &
Comprender la evolución de la presencia digital desde el currículum impreso al portafolio en línea. &
Identifica los componentes de una presencia digital profesional contemporánea. &
Distingue marca personal, identidad visual, sitio web personal, portafolio. &
Analiza 3 presencias digitales profesionales (uno excelente, uno regular, uno pobre). &
Honra el binomio ancestral oficio + palabra como anclaje de la reputación digital. \\
\hline

% -- FILA 2: Apropiación y uso --
\textbf{Apropiación y uso de la tecnología} &
Construir un sitio web personal profesional usando plantillas y IA. &
Maneja plantillas web (HTML/CSS responsive), copywriting con IA y SEO básico. &
Conoce los principios del diseño web responsivo, accesibilidad y SEO. &
Publica un sitio web personal en Netlify/Vercel con 4 secciones mínimas. &
Sostiene voz propia aún usando plantillas y IA: 30\% mínimo de intervención humana visible. \\
\hline

% -- FILA 3: Solución de problemas --
\textbf{Solución de problemas con tecnología} &
Diagnosticar y resolver problemas de presencia digital con criterio profesional. &
Diagnostica problemas (foto poco profesional, biografía sin voz, perfiles inconsistentes) y los corrige. &
Identifica los 10 errores comunes en presencia digital. &
Revisa su presencia digital con lista de chequeo y aplica correcciones. &
Acepta retroalimentación honesta sobre su marca personal. \\
\hline

% -- FILA 4: Tecnología y sociedad --
\textbf{Tecnología y sociedad} &
Lanzar su presencia digital al mundo con criterio ético y profesional. &
Compila perfiles, sitio web y métricas para su lanzamiento P1 (S10). &
Reconoce el rol de la presencia digital en empleabilidad y educación superior. &
Lanza presencia digital pública y registra métricas básicas (visitas, fuente). &
Asume reputación digital como compromiso de largo plazo. \\
\hline

% -- FILA 5: PC (verde) --
\rowcolor{verdePCclaro}
\makecell[l]{\textbf{\color{verdePC}Pensamiento}\\\textbf{\color{verdePC}Computacional}\\
\footnotesize Wing (2006)\\\footnotesize MEN Or. 2022\\\footnotesize\textit{Abstracción · Algoritmos (SEO) · Modelado}}\
 &
Aplicar el PC para diseñar una presencia digital sistemática, reconociendo los algoritmos de visibilidad (SEO) como sistemas de reglas con lógica propia y verificable. &
Analiza el algoritmo de SEO básico como sistema de reglas para visibilidad (S3–S5); diseña su presencia digital como arquitectura de información sistemática con jerarquía clara. &
Describe el SEO como algoritmo de clasificación basado en reglas de relevancia, autoridad y experiencia de usuario; reconoce la estructura web como jerarquía de información modelable. &
Diseña la arquitectura de su sitio web personal: mapa de secciones, jerarquía de contenidos, metadatos SEO y criterios de accesibilidad; documenta las decisiones con criterio algorítmico. &
Sostiene voz propia aún usando plantillas; asume la reputación digital como compromiso de largo plazo construido sistemáticamente. \\
\hline

% -- FILA 6: SEL (violeta) --
\rowcolor{moradoSELclaro}
\makecell[l]{\textbf{\color{moradoSEL}Habilidades}\\\textbf{\color{moradoSEL}Socioemocionales}\\
\footnotesize Ley 2383/2024\\\footnotesize CASEL 2020\\\footnotesize\textit{Autoconciencia · Toma de decisiones responsable}}\
 &
Desarrollar autoconciencia sobre la coherencia entre los valores personales y la presencia digital, y tomar decisiones responsables sobre la imagen pública que se proyecta al mundo profesional. &
Evalúa si su presencia digital actual refleja sus valores reales y la imagen profesional que desea proyectar a futuro. &
Describe las consecuencias de largo plazo de la huella digital; distingue entre imagen auténtica e imagen gestionada con integridad. &
Realiza una auditoría de su huella digital actual (redes, fotos, publicaciones) y diseña un plan de mejora con 5 acciones concretas. &
Construye su presencia digital con coherencia entre quien es y quien proyecta; asume la reputación digital como legado, no como estrategia. \\
\hline

\end{longtable}
\medskip
\noindent\textbf{Nota:} Las filas 5 (\textcolor{verdePC}{verde}) y 6 (\textcolor{moradoSEL}{violeta}) son componentes propios de este plan, adicionales a los cuatro componentes de la Guía No.~30 del MEN.

\subsubsection{Periodo 2 · Automatización y gestión de procesos}
\noindent\begin{tabular}{ll}
\textbf{Grado:} Undécimo & \textbf{Periodo:} Segundo \\
\textbf{Proyecto integrador:} Automatiza un proceso real & \textbf{Anclaje:} El maestro carpintero (lectura del proceso) \\
\end{tabular}
\medskip

\begin{longtable}{|L{3.8cm}|L{3.3cm}|L{3.3cm}|L{2.5cm}|L{2.5cm}|L{2.5cm}|}
\caption{Estructura de malla curricular — G11° Periodo 2 · Automatización y gestión de procesos}
\label{tab:g11p2}\\
\hline
\rowcolor{azulInst}
\multirow{2}{*}{\color{white}\scriptsize\textbf{\makecell{ESTÁNDARES, LINEAMIENTOS\\U ORIENTACIONES}}} &
\multirow{2}{*}{\color{white}\scriptsize\textbf{\makecell{COMPETENCIAS}}} &
\multirow{2}{*}{\color{white}\scriptsize\textbf{\makecell{APRENDIZAJES}}} &
\multicolumn{3}{c|}{%
  \color{white}\makecell{\textbf{\footnotesize DESCRIPTORES DE DESEMPEÑO}\\[1pt]\textit{\scriptsize Metas del aprendizaje por periodo}}%
}\\
\cline{4-6}
\rowcolor{azulClaro}
\cellcolor{azulInst} & \cellcolor{azulInst} & \cellcolor{azulInst} &
\textbf{\scriptsize COGNITIVOS} &
\textbf{\scriptsize PROCEDIMENTALES} &
\textbf{\scriptsize ACTITUDINALES} \\
\hline
\endfirsthead
\hline
\rowcolor{azulInst}
\multirow{2}{*}{\color{white}\scriptsize\textbf{\makecell{ESTÁNDARES, LINEAMIENTOS\\U ORIENTACIONES}}} &
\multirow{2}{*}{\color{white}\scriptsize\textbf{\makecell{COMPETENCIAS}}} &
\multirow{2}{*}{\color{white}\scriptsize\textbf{\makecell{APRENDIZAJES}}} &
\multicolumn{3}{c|}{%
  \color{white}\makecell{\textbf{\footnotesize DESCRIPTORES DE DESEMPEÑO}\\[1pt]\textit{\scriptsize Metas del aprendizaje por periodo}}%
}\\
\cline{4-6}
\rowcolor{azulClaro}
\cellcolor{azulInst} & \cellcolor{azulInst} & \cellcolor{azulInst} &
\textbf{\scriptsize COGNITIVOS} &
\textbf{\scriptsize PROCEDIMENTALES} &
\textbf{\scriptsize ACTITUDINALES} \\
\hline
\endhead

% -- FILA 1: Naturaleza y evolución --
\textbf{Naturaleza y evolución de la tecnología} &
Comprender la mejora continua (Kaizen) y los principios de la automatización. &
Identifica procesos repetitivos automatizables y sus indicadores (KPIs). &
Define BPMN, KPI, trigger, flujo y mejora continua. &
Mapea un proceso real con diagrama BPMN básico. &
Honra al maestro carpintero como lector del proceso. \\
\hline

% -- FILA 2: Apropiación y uso --
\textbf{Apropiación y uso de la tecnología} &
Construir automatizaciones de bajo código con triggers. &
Maneja Zapier, Make o n8n para automatizar tareas con disparadores. &
Conoce los componentes de una automatización: trigger + acción + filtros + transformación. &
Implementa una automatización real con mínimo 3 pasos. &
Cuida que la automatización no excluya casos legítimos. \\
\hline

% -- FILA 3: Solución de problemas --
\textbf{Solución de problemas con tecnología} &
Diagnosticar y resolver cuellos de botella en procesos del entorno. &
Identifica dónde se pierde tiempo en un proceso y propone solución. &
Reconoce las 7 mudas (desperdicios) del pensamiento Lean. &
Analiza un proceso real y diseña 3 mejoras priorizadas. &
Aplica Kaizen: mejora pequeña, continua, sostenida. \\
\hline

% -- FILA 4: Tecnología y sociedad --
\textbf{Tecnología y sociedad} &
Aplicar IA en mensajería responsable (chatbots con criterio). &
Construye un chatbot sencillo con flujo de conversación claro. &
Reconoce las consideraciones éticas de mensajería con IA: claridad sobre qué es IA, privacidad. &
Implementa un chatbot real para un caso concreto del colegio o la casa. &
Honra la palabra: si es chatbot, lo declara abiertamente al usuario. \\
\hline

% -- FILA 5: PC (verde) --
\rowcolor{verdePCclaro}
\makecell[l]{\textbf{\color{verdePC}Pensamiento}\\\textbf{\color{verdePC}Computacional}\\
\footnotesize Wing (2006)\\\footnotesize MEN Or. 2022\\\footnotesize\textit{Algoritmos · Descomposición de procesos · Abstracción (BPMN)}}\
 &
Aplicar el PC para modelar procesos con BPMN como diagramas de flujo y diseñar automatizaciones como programación de bajo código con triggers, condiciones y acciones. &
Modela un proceso real con BPMN básico (S1–S3); implementa automatización en Zapier/Make como algoritmo de trigger-acción-filtro (S4–S7); aplica Kaizen como depuración iterativa. &
Reconoce BPMN como lenguaje de modelado algorítmico de procesos; comprende la automatización como programa con trigger (evento), condición (filtro) y acción (resultado) verificable. &
Mapea un proceso del colegio/casa con diagrama BPMN; implementa una automatización de mínimo 3 pasos con trigger, condición y acción documentados y verificados con usuarios. &
Honra al maestro carpintero: lee el proceso antes de intervenir; mejora pequeña, continua y sostenida; la automatización no excluye casos legítimos. \\
\hline

% -- FILA 6: SEL (violeta) --
\rowcolor{moradoSELclaro}
\makecell[l]{\textbf{\color{moradoSEL}Habilidades}\\\textbf{\color{moradoSEL}Socioemocionales}\\
\footnotesize Ley 2383/2024\\\footnotesize CASEL 2020\\\footnotesize\textit{Conciencia social · Toma de decisiones responsable}}\
 &
Desarrollar conciencia social sobre el impacto de la automatización en el empleo y la vida de personas reales, y tomar decisiones responsables que preserven la dignidad y la inclusión frente a los sistemas automatizados. &
Identifica los riesgos de desplazamiento laboral y exclusión que puede generar una automatización mal diseñada, y los evalúa antes de implementar. &
Describe los efectos históricos de la automatización en el mercado laboral; distingue entre automatización que libera y automatización que excluye. &
Incluye en el diseño de su automatización un análisis de impacto humano: quién se beneficia, quién podría quedar excluido y cómo mitigarlo. &
Diseña automatizaciones con criterio humano: la eficiencia nunca justifica la exclusión; asume la responsabilidad de los sistemas que pone en marcha. \\
\hline

\end{longtable}
\medskip
\noindent\textbf{Nota:} Las filas 5 (\textcolor{verdePC}{verde}) y 6 (\textcolor{moradoSEL}{violeta}) son componentes propios de este plan, adicionales a los cuatro componentes de la Guía No.~30 del MEN.

\subsubsection{Periodo 3 · Microempresa con MVP, datos y ética}
\noindent\begin{tabular}{ll}
\textbf{Grado:} Undécimo & \textbf{Periodo:} Tercero \\
\textbf{Proyecto integrador:} Microempresa con MVP, datos y ética & \textbf{Anclaje:} Oficios que escuchan al pueblo (validación de mercado ancestral) \\
\end{tabular}
\medskip

\begin{longtable}{|L{3.8cm}|L{3.3cm}|L{3.3cm}|L{2.5cm}|L{2.5cm}|L{2.5cm}|}
\caption{Estructura de malla curricular — G11° Periodo 3 · Microempresa con MVP, datos y ética}
\label{tab:g11p3}\\
\hline
\rowcolor{azulInst}
\multirow{2}{*}{\color{white}\scriptsize\textbf{\makecell{ESTÁNDARES, LINEAMIENTOS\\U ORIENTACIONES}}} &
\multirow{2}{*}{\color{white}\scriptsize\textbf{\makecell{COMPETENCIAS}}} &
\multirow{2}{*}{\color{white}\scriptsize\textbf{\makecell{APRENDIZAJES}}} &
\multicolumn{3}{c|}{%
  \color{white}\makecell{\textbf{\footnotesize DESCRIPTORES DE DESEMPEÑO}\\[1pt]\textit{\scriptsize Metas del aprendizaje por periodo}}%
}\\
\cline{4-6}
\rowcolor{azulClaro}
\cellcolor{azulInst} & \cellcolor{azulInst} & \cellcolor{azulInst} &
\textbf{\scriptsize COGNITIVOS} &
\textbf{\scriptsize PROCEDIMENTALES} &
\textbf{\scriptsize ACTITUDINALES} \\
\hline
\endfirsthead
\hline
\rowcolor{azulInst}
\multirow{2}{*}{\color{white}\scriptsize\textbf{\makecell{ESTÁNDARES, LINEAMIENTOS\\U ORIENTACIONES}}} &
\multirow{2}{*}{\color{white}\scriptsize\textbf{\makecell{COMPETENCIAS}}} &
\multirow{2}{*}{\color{white}\scriptsize\textbf{\makecell{APRENDIZAJES}}} &
\multicolumn{3}{c|}{%
  \color{white}\makecell{\textbf{\footnotesize DESCRIPTORES DE DESEMPEÑO}\\[1pt]\textit{\scriptsize Metas del aprendizaje por periodo}}%
}\\
\cline{4-6}
\rowcolor{azulClaro}
\cellcolor{azulInst} & \cellcolor{azulInst} & \cellcolor{azulInst} &
\textbf{\scriptsize COGNITIVOS} &
\textbf{\scriptsize PROCEDIMENTALES} &
\textbf{\scriptsize ACTITUDINALES} \\
\hline
\endhead

% -- FILA 1: Naturaleza y evolución --
\textbf{Naturaleza y evolución de la tecnología} &
Comprender el ciclo emprendedor moderno y su raíz en escuchar al pueblo. &
Identifica los 10 pasos del ciclo emprendedor desde la idea al lanzamiento. &
Define MVP, validación, pivote, traction, pitch. &
Documenta su recorrido emprendedor con bitácora. &
Honra a los oficios del Valle que nacían de escuchar al pueblo. \\
\hline

% -- FILA 2: Apropiación y uso --
\textbf{Apropiación y uso de la tecnología} &
Construir un MVP digital y recolectar datos de validación reales. &
Implementa un MVP en plataforma de bajo código y recolecta datos de mínimo 20 usuarios. &
Conoce las plataformas no-code para MVPs (Glide, Bubble, sitios web). &
Implementa un MVP funcional y recolecta datos de mínimo 20 usuarios reales. &
Acepta el MVP como imperfecto pero útil para aprender. \\
\hline

% -- FILA 3: Solución de problemas --
\textbf{Solución de problemas con tecnología} &
Resolver problemas reales de microemprendimiento con datos, presupuesto y ética. &
Construye plan completo: ideación, validación, modelo, MVP, datos, pitch, presupuesto, ética, versión final. &
Reconoce el ciclo completo del emprendimiento responsable. &
Entrega plan completo de 20+ páginas con todas las secciones. &
Sostiene la ética en cada decisión (palabra empeñada vs atajo deshonesto). \\
\hline

% -- FILA 4: Tecnología y sociedad --
\textbf{Tecnología y sociedad} &
Sustentar el microemprendimiento ante una comunidad real con compromiso ético. &
Sustenta ante jurado real (docentes + emprendedores) con declaración ética de uso de IA (S10). &
Reconoce su rol como emprendedor responsable en la economía y la infoesfera. &
Sustenta públicamente y registra compromiso de mejora a 6 meses. &
Honra al pueblo como destinatario y validador del oficio emprendedor. \\
\hline

% -- FILA 5: PC (verde) --
\rowcolor{verdePCclaro}
\makecell[l]{\textbf{\color{verdePC}Pensamiento}\\\textbf{\color{verdePC}Computacional}\\
\footnotesize Wing (2006)\\\footnotesize MEN Or. 2022\\\footnotesize\textit{Prototipado iterativo (MVP) · Depuración con datos · Modelado}}\
 &
Aplicar el PC para diseñar un ciclo emprendedor iterativo con MVP como prototipo mínimo verificable, validación de datos y depuración basada en evidencia real. &
Diseña el MVP como prototipo mínimo verificable (S3–S5); valida con datos de mínimo 20 usuarios reales; itera y pivota basado en evidencia (S6–S8). &
Reconoce el ciclo emprendedor (idea / modelo / MVP / validación / pivote / escala) como algoritmo iterativo de depuración; comprende el MVP como prototipo mínimo con criterio de verificación definido. &
Implementa un MVP funcional en plataforma no-code; recolecta datos de mínimo 20 usuarios; analiza resultados y documenta 3 pivotes o mejoras con evidencia verificable. &
Acepta el MVP como imperfecto pero útil para aprender; sostiene la ética en cada decisión; honra al pueblo como destinatario y validador real del oficio emprendedor. \\
\hline

% -- FILA 6: SEL (violeta) --
\rowcolor{moradoSELclaro}
\makecell[l]{\textbf{\color{moradoSEL}Habilidades}\\\textbf{\color{moradoSEL}Socioemocionales}\\
\footnotesize Ley 2383/2024\\\footnotesize CASEL 2020\\\footnotesize\textit{Autogestión · Conciencia social}}\
 &
Desarrollar resiliencia emprendedora para afrontar fracasos y pivotes con ecuanimidad, y conciencia social para emprender con responsabilidad hacia la comunidad y el entorno como legado del oficio ancestral. &
Identifica sus respuestas emocionales frente al fracaso del MVP o el rechazo de usuarios, y desarrolla estrategias de resiliencia y reencuadre productivo. &
Describe la resiliencia emprendedora como capacidad de aprender del fracaso sin perder la orientación al propósito; reconoce el pivote como decisión estratégica, no como derrota. &
Documenta en su bitácora al menos 2 momentos de fracaso o rechazo, la emoción que generó y la decisión que tomó a partir de ellos. &
Emprende con ecuanimidad y propósito social; honra al pueblo como destinatario del oficio; deja el legado del emprendedor que escucha antes de vender. \\
\hline

\end{longtable}
\medskip
\noindent\textbf{Nota:} Las filas 5 (\textcolor{verdePC}{verde}) y 6 (\textcolor{moradoSEL}{violeta}) son componentes propios de este plan, adicionales a los cuatro componentes de la Guía No.~30 del MEN.



\clearpage
\end{landscape}

\clearpage
% ============================================================
\section{Anexo B — Temario: títulos de sesión por grado y periodo}
% ============================================================

\noindent El siguiente temario lista los títulos de las \textbf{180 guías}
(6 grados $\times$ 3 periodos $\times$ 10 sesiones), en el orden en que el
estudiante las recorre. Es la ruta de aprendizaje completa publicada en la
Plataforma Conéctate; cada título corresponde a una guía navegable en línea y
descargable en PDF.

% ============================================================
% Temario — ruta de sesiones por grado y periodo
% Generado desde content/guias/*.yaml (180 guias)
% ============================================================

\needspace{6\baselineskip}
\bigskip
\noindent{\large\color{vinoInst}\textbf{Grado Sexto (6°)}}
\vspace{1pt}\par\noindent{\color{vinoInst}\rule{\textwidth}{1pt}}
\vspace{3pt}

\needspace{5\baselineskip}
\noindent\textbf{Periodo 1 · Comunicación e identidad digital}\hfill{\footnotesize\color{vinoInst}10 guías}\par
\begin{enumerate}[leftmargin=2.6em, itemsep=0.6pt, topsep=2pt, parsep=0pt, label=\footnotesize\arabic*.]
\footnotesize
  \item Bienvenida a Tecnología --- cómo aprenderemos este año
  \item Mi identidad digital --- cómo me presento en la red sin perderme
  \item Mi correo institucional --- la primera dirección que firmas como tuya
  \item Netiqueta --- los 10 acuerdos para llevarte bien en internet
  \item Historia de los medios --- del pregonero al chat de WhatsApp
  \item Mi huella digital --- el rastro que dejas sin darte cuenta
  \item Privacidad y datos personales --- cerrar las puertas con criterio
  \item Riesgos digitales --- reconocer la trampa antes de caer
  \item Comunicación responsable --- los 3 filtros antes de publicar
  \item Mi presentación digital --- 2 minutos para mostrar lo que aprendí
\end{enumerate}

\needspace{5\baselineskip}
\noindent\textbf{Periodo 2 · Hardware, software y la máquina por dentro}\hfill{\footnotesize\color{vinoInst}10 guías}\par
\begin{enumerate}[leftmargin=2.6em, itemsep=0.6pt, topsep=2pt, parsep=0pt, label=\footnotesize\arabic*.]
\footnotesize
  \item Apertura periodo 2 --- las máquinas que nos rodean por dentro
  \item ¿Qué es un computador? --- las 4 funciones básicas de toda máquina inteligente
  \item El gabinete por dentro --- conocer las piezas como las conoce el relojero
  \item Periféricos --- las manos, los ojos y la voz del computador
  \item Hardware vs software --- el cuerpo y las instrucciones
  \item El sistema operativo --- el gerente que coordina toda la máquina
  \item Archivos y carpetas --- ordenar mi biblioteca digital
  \item Mantenimiento básico --- cuidar el computador como cuida el relojero
  \item Postura y vida útil --- cuidarte tú mientras usas la máquina
  \item Mi computador ideal --- diseñar la máquina que yo necesito
\end{enumerate}

\needspace{5\baselineskip}
\noindent\textbf{Periodo 3 · Escribir y buscar con criterio}\hfill{\footnotesize\color{vinoInst}10 guías}\par
\begin{enumerate}[leftmargin=2.6em, itemsep=0.6pt, topsep=2pt, parsep=0pt, label=\footnotesize\arabic*.]
\footnotesize
  \item Apertura periodo 3 --- el oficio de escribir y buscar con criterio
  \item ¿Qué es un procesador de texto? --- mi primera página escrita en digital
  \item Formato básico --- darle voz visual a lo que escribo
  \item Listas, viñetas y numeración --- ordenar la información para que se entienda
  \item Tablas e imágenes --- mostrar datos y agregar voz visual
  \item Encabezado, pie de página y números --- terminar el documento como un profesional
  \item ¿Qué es internet? --- cómo viajan los datos para llegar a tu pantalla
  \item Buscar bien en internet --- encontrar lo que necesito en segundos
  \item Evaluar fuentes --- distinguir lo verdadero de lo inventado en internet
  \item Mi primer documento informativo --- producto que reúne todo lo aprendido
\end{enumerate}

\needspace{6\baselineskip}
\bigskip
\noindent{\large\color{vinoInst}\textbf{Grado Séptimo (7°)}}
\vspace{1pt}\par\noindent{\color{vinoInst}\rule{\textwidth}{1pt}}
\vspace{3pt}

\needspace{5\baselineskip}
\noindent\textbf{Periodo 1 · Trabajo colaborativo en la nube (Microsoft 365)}\hfill{\footnotesize\color{vinoInst}10 guías}\par
\begin{enumerate}[leftmargin=2.6em, itemsep=0.6pt, topsep=2pt, parsep=0pt, label=\footnotesize\arabic*.]
\footnotesize
  \item Apertura periodo 1 --- vamos a hacer minga digital
  \item Microsoft 365 --- conocer las herramientas del taller digital
  \item OneDrive --- mi nube personal y mi cuarto de archivos
  \item Word, Excel y PowerPoint en línea --- las 3 herramientas estrella
  \item Compartir y permisos --- abrir mi taller a invitados con criterio
  \item Coautoría asincrónica --- 3 personas escribiendo al mismo tiempo
  \item Comentarios y sugerencias --- revisar el trabajo ajeno con respeto
  \item Versiones e historial --- el sistema recuerda todo lo que escribiste
  \item Outlook y Teams --- comunicarme con voz profesional
  \item Cosecha del periodo --- mi primer proyecto colaborativo
\end{enumerate}

\needspace{5\baselineskip}
\noindent\textbf{Periodo 2 · Algoritmia y pensamiento computacional (Scratch)}\hfill{\footnotesize\color{vinoInst}10 guías}\par
\begin{enumerate}[leftmargin=2.6em, itemsep=0.6pt, topsep=2pt, parsep=0pt, label=\footnotesize\arabic*.]
\footnotesize
  \item Apertura periodo 2 --- pensar como teje la abuela
  \item ¿Qué es un algoritmo? --- la receta paso a paso
  \item Pensamiento computacional --- las 4 técnicas del oficio
  \item Diagramas de flujo --- dibujar el algoritmo
  \item Variables --- cajas con etiqueta que guardan información
  \item Condicionales --- el si... entonces de los algoritmos
  \item Bucles --- repeticiones inteligentes en los algoritmos
  \item Scratch --- mi primer programa con bloques
  \item Scratch · narrativa interactiva --- el cuento que el lector escoge
  \item Cosecha del periodo 2 --- mi primer programa Scratch completo
\end{enumerate}

\needspace{5\baselineskip}
\noindent\textbf{Periodo 3 · Inteligencia artificial con criterio}\hfill{\footnotesize\color{vinoInst}10 guías}\par
\begin{enumerate}[leftmargin=2.6em, itemsep=0.6pt, topsep=2pt, parsep=0pt, label=\footnotesize\arabic*.]
\footnotesize
  \item Apertura periodo 3 --- la IA como el nuevo consejero del barrio
  \item ¿Qué es la IA? --- definición, historia y casos cotidianos
  \item Tipos de IA --- estrecha, general, modelos especializados
  \item Cómo aprende una IA --- datos, entrenamiento y el problema del sesgo
  \item Modelos de lenguaje (LLMs) --- ChatGPT, Claude, Gemini
  \item Prompting básico --- cómo hablarle a la IA para que te responda bien
  \item Prompting avanzado --- técnicas que diferencian al usuario experto
  \item Ética de la IA --- sesgo, privacidad, derechos y deepfakes
  \item IA y mis tareas --- política personal de uso responsable en estudios
  \item Cosecha del periodo 3 --- mi proyecto con IA responsable
\end{enumerate}

\needspace{6\baselineskip}
\bigskip
\noindent{\large\color{vinoInst}\textbf{Grado Octavo (8°)}}
\vspace{1pt}\par\noindent{\color{vinoInst}\rule{\textwidth}{1pt}}
\vspace{3pt}

\needspace{5\baselineskip}
\noindent\textbf{Periodo 1 · Lógica computacional avanzada}\hfill{\footnotesize\color{vinoInst}10 guías}\par
\begin{enumerate}[leftmargin=2.6em, itemsep=0.6pt, topsep=2pt, parsep=0pt, label=\footnotesize\arabic*.]
\footnotesize
  \item Phronesis con datos --- preguntar antes de calcular
  \item Tipos de datos y formato de celdas en Excel
  \item Tablas de datos --- campos, registros y limpieza con bitácora
  \item SUMA, PROMEDIO, MAX, MIN --- 4 funciones que dirigen decisiones
  \item Referencias relativas y absolutas en Excel
  \item Fórmulas compuestas --- operadores, porcentajes y jerarquía PEMDAS
  \item Gráficos básicos --- barras, líneas y circulares con honestidad visual
  \item Formato condicional y validación de datos --- alertas visibles
  \item Mini-estudio de datos del entorno escolar
  \item Sustentación del mini-estudio --- 4 minutos de phronesis comunicativa
\end{enumerate}

\needspace{5\baselineskip}
\noindent\textbf{Periodo 2 · Datos como tablas: pensar con estructura}\hfill{\footnotesize\color{vinoInst}10 guías}\par
\begin{enumerate}[leftmargin=2.6em, itemsep=0.6pt, topsep=2pt, parsep=0pt, label=\footnotesize\arabic*.]
\footnotesize
  \item Lógica avanzada --- if-else con AND, OR y NOT
  \item Tablas de verdad --- AND, OR, NOT como reglas universales
  \item Algoritmos --- del pseudocódigo al diagrama de flujo
  \item Sensores en MakeCode --- el micro:bit que siente
  \item Actuadores --- LEDs, sonido y coreografía
  \item Sensores con variables, umbrales y calibración
  \item Depuración --- buscar el error como el relojero busca el tic-tac
  \item Alertas con lógica compuesta --- validación con escenarios
  \item Proyecto de monitoreo --- el micro:bit como vigía
  \item Sustentación del proyecto técnico --- honestidad del oficio
\end{enumerate}

\needspace{5\baselineskip}
\noindent\textbf{Periodo 3 · Diseño visual y comunicación digital}\hfill{\footnotesize\color{vinoInst}10 guías}\par
\begin{enumerate}[leftmargin=2.6em, itemsep=0.6pt, topsep=2pt, parsep=0pt, label=\footnotesize\arabic*.]
\footnotesize
  \item Multimedia --- diseño visual y jerarquía de información
  \item Animaciones e hipervínculos --- narrativa interactiva
  \item Narrativa de alto impacto --- estructura TED-style
  \item Edición de imagen --- ética visual y derechos de imagen
  \item Audio --- captura, edición y formatos
  \item Ciberbullying --- prevención y Línea 141
  \item Grooming y sexting --- riesgos digitales y Ley 1336
  \item Estética de la liberación --- cultura popular y producción ética
  \item Proyecto integrador comunitario --- voz para el cambio social
  \item Sustentación pública y portafolio digital del grado 8
\end{enumerate}

\needspace{6\baselineskip}
\bigskip
\noindent{\large\color{vinoInst}\textbf{Grado Noveno (9°)}}
\vspace{1pt}\par\noindent{\color{vinoInst}\rule{\textwidth}{1pt}}
\vspace{3pt}

\needspace{5\baselineskip}
\noindent\textbf{Periodo 1 · Técnica, ciencia y tecnología}\hfill{\footnotesize\color{vinoInst}10 guías}\par
\begin{enumerate}[leftmargin=2.6em, itemsep=0.6pt, topsep=2pt, parsep=0pt, label=\footnotesize\arabic*.]
\footnotesize
  \item ¿Qué es técnica? --- del cuerpo al instrumento
  \item La rueda y el agua --- máquinas simples del campo colombiano
  \item La balanza y el peso --- medir como acto político
  \item Línea del tiempo de la técnica --- del fuego a la imprenta
  \item Revolución industrial --- máquinas que cambiaron el trabajo
  \item La imprenta --- escribir muchas veces lo mismo
  \item La electricidad --- luz, motor, comunicación
  \item Era digital --- del ábaco al chip
  \item Tecnologías propias --- saberes que el mundo aprende ahora
  \item Cosecha P1 --- manifiesto del técnico crítico
\end{enumerate}

\needspace{5\baselineskip}
\noindent\textbf{Periodo 2 · Diseño editorial y maquetación digital}\hfill{\footnotesize\color{vinoInst}10 guías}\par
\begin{enumerate}[leftmargin=2.6em, itemsep=0.6pt, topsep=2pt, parsep=0pt, label=\footnotesize\arabic*.]
\footnotesize
  \item ¿Qué es diseñar? --- orden, jerarquía, intención
  \item La página y la cuadrícula --- primer principio del orden
  \item Tipografía --- la voz visible de las palabras
  \item Color en editorial --- contraste, jerarquía, accesibilidad
  \item Imagen y texto --- el ritmo de la lectura
  \item Herramientas digitales --- Figma, Canva o Adobe Express
  \item Accesibilidad lectora --- quién NO puede leer mi revista y por qué
  \item Edición y corrección --- el ojo crítico
  \item Maquetar la revista completa --- flujo de trabajo
  \item Cosecha P2 --- sustentación pública de la revista
\end{enumerate}

\needspace{5\baselineskip}
\noindent\textbf{Periodo 3 · Datos, estadística y visualización}\hfill{\footnotesize\color{vinoInst}10 guías}\par
\begin{enumerate}[leftmargin=2.6em, itemsep=0.6pt, topsep=2pt, parsep=0pt, label=\footnotesize\arabic*.]
\footnotesize
  \item ¿Qué es un dato? --- recolección con propósito
  \item Tablas y registros --- la columna como tipo, la fila como caso
  \item Fórmulas básicas --- SUMA, PROMEDIO, MAX, MIN, CONTAR
  \item Filtros y orden --- preguntar a los datos
  \item Tablas dinámicas --- agrupar y comparar
  \item Gráficos --- la imagen de los datos
  \item Limpieza de datos --- el oficio del archivero
  \item Visualización honesta --- gráficos que mienten
  \item Insights --- del dato a la decisión
  \item Cosecha P3 --- sustentación del mini-estudio
\end{enumerate}

\needspace{6\baselineskip}
\bigskip
\noindent{\large\color{vinoInst}\textbf{Grado Décimo (10°)}}
\vspace{1pt}\par\noindent{\color{vinoInst}\rule{\textwidth}{1pt}}
\vspace{3pt}

\needspace{5\baselineskip}
\noindent\textbf{Periodo 1 · Edición digital y autoría con IA responsable}\hfill{\footnotesize\color{vinoInst}10 guías}\par
\begin{enumerate}[leftmargin=2.6em, itemsep=0.6pt, topsep=2pt, parsep=0pt, label=\footnotesize\arabic*.]
\footnotesize
  \item El oficio del editor --- qué hace y qué no hace
  \item Concebir tu libro --- género, tema, audiencia y escaleta
  \item Prompting técnico para escritura --- las 5 partes del prompt profesional
  \item Iterar y refinar texto con IA --- del borrador al capítulo terminado
  \item Estructura editorial --- arco narrativo, ritmo y capítulos
  \item Derechos de autor y propiedad intelectual con IA --- Ley 23 de 1982 y obras derivadas
  \item Portada con IA generativa de imagen --- herramientas gratuitas
  \item Diagramación con herramientas gratuitas --- Canva, Google Docs, LibreOffice
  \item Producción final del libro --- 80 páginas compiladas, revisadas y firmadas
  \item Sustentación + feria del libro escolar
\end{enumerate}

\needspace{5\baselineskip}
\noindent\textbf{Periodo 2 · Comunicación profesional con IA}\hfill{\footnotesize\color{vinoInst}10 guías}\par
\begin{enumerate}[leftmargin=2.6em, itemsep=0.6pt, topsep=2pt, parsep=0pt, label=\footnotesize\arabic*.]
\footnotesize
  \item 4 tipos de texto profesional --- informe, comercial, técnico, estudio de mercado
  \item Estructura del informe técnico --- introducción, desarrollo, conclusiones
  \item Google Docs para informes profesionales --- estilos, índice y citas
  \item Markdown --- el formato del oficio digital portable
  \item Encuestas con phronesis --- Google Forms más IA para preguntar bien
  \item Análisis cualitativo con IA --- qué hace y qué no
  \item Informe comercial --- para audiencia real
  \item Correo profesional con IA --- asistente, no reemplazo
  \item Mini-proyecto --- informe técnico con datos reales y IA documentada
  \item Sustentación del periodo 2 --- defensa del oficio
\end{enumerate}

\needspace{5\baselineskip}
\noindent\textbf{Periodo 3 · Microemprendimiento y gestión digital}\hfill{\footnotesize\color{vinoInst}10 guías}\par
\begin{enumerate}[leftmargin=2.6em, itemsep=0.6pt, topsep=2pt, parsep=0pt, label=\footnotesize\arabic*.]
\footnotesize
  \item Ofimática con IA --- el copiloto del oficio digital cotidiano
  \item Contabilidad básica --- ingresos, egresos y balance personal o familiar
  \item Google Sheets + IA para contabilidad personal o microempresa
  \item Flujo de caja y punto de equilibrio --- proyectar con IA
  \item Visualización de datos contables con IA --- gráficos honestos
  \item Encuestas y análisis de mercado con IA
  \item Plan de negocio con IA --- Business Model Canvas asistido
  \item Comunicación empresarial profesional con IA
  \item Proyecto integrador --- micro-emprendimiento con plan completo
  \item Sustentación final + feria de microempresas escolares
\end{enumerate}

\needspace{6\baselineskip}
\bigskip
\noindent{\large\color{vinoInst}\textbf{Grado Undécimo (11°)}}
\vspace{1pt}\par\noindent{\color{vinoInst}\rule{\textwidth}{1pt}}
\vspace{3pt}

\needspace{5\baselineskip}
\noindent\textbf{Periodo 1 · Presencia digital empresarial con IA}\hfill{\footnotesize\color{vinoInst}10 guías}\par
\begin{enumerate}[leftmargin=2.6em, itemsep=0.6pt, topsep=2pt, parsep=0pt, label=\footnotesize\arabic*.]
\footnotesize
  \item Mi marca digital --- del oficio heredado a la huella propia
  \item Identidad visual --- del grafismo ancestral a la paleta propia
  \item Sitio web personal --- fundamentos de la casa digital
  \item Plantillas web --- elegir y adaptar sin perderse uno mismo
  \item Copywriting con IA --- escribir con criterio, no a ciegas
  \item SEO básico --- hacerme visible sin engaños
  \item Imagen propia --- cómo se ve mi marca en foto y video
  \item Perfiles profesionales --- LinkedIn y portafolio coherentes
  \item Métricas básicas --- quién entra a mi sitio y qué busca
  \item Lanzamiento P1 --- mi presencia digital sale al mundo
\end{enumerate}

\needspace{5\baselineskip}
\noindent\textbf{Periodo 2 · Automatización y gestión de procesos}\hfill{\footnotesize\color{vinoInst}10 guías}\par
\begin{enumerate}[leftmargin=2.6em, itemsep=0.6pt, topsep=2pt, parsep=0pt, label=\footnotesize\arabic*.]
\footnotesize
  \item Mapeo de procesos --- ver dónde se pierde el tiempo
  \item Diagramas de flujo --- lenguaje BPMN básico
  \item Formularios digitales --- captura de información sin papel
  \item Hojas de cálculo como base de datos
  \item Automatización con triggers --- Zapier, Make, n8n
  \item IA en mensajería --- clasificar y responder con criterio
  \item Chatbots sin código --- flujos de conversación
  \item KPIs --- métricas que importan en un proceso
  \item Mejora continua --- Kaizen aplicado a procesos digitales
  \item Caso real --- automatizar un proceso del colegio o la casa
\end{enumerate}

\needspace{5\baselineskip}
\noindent\textbf{Periodo 3 · Microempresa con MVP, datos y ética}\hfill{\footnotesize\color{vinoInst}10 guías}\par
\begin{enumerate}[leftmargin=2.6em, itemsep=0.6pt, topsep=2pt, parsep=0pt, label=\footnotesize\arabic*.]
\footnotesize
  \item Ideación --- un problema real de mi comunidad
  \item Validación --- hablar antes de construir
  \item Modelo de negocio --- el lienzo de tu proyecto
  \item MVP digital --- una versión imperfecta para probar
  \item Datos del MVP --- recoger evidencia con rigor
  \item Pitch --- contar la historia de tu proyecto
  \item Presupuesto --- cuánto siembras, cuánto cosechas
  \item Ética y legal --- la palabra empeñada en lo digital
  \item Versión final del MVP
  \item Sustentación pública --- el aprendiz frente a los mayores
\end{enumerate}



% ============================================================
% sección\,21 — APROBACIÓN Y FIRMA
% ============================================================
\newpage
\thispagestyle{empty}

% ---- COLOFÓN / CIERRE ----
\vspace*{1.2cm}
\begin{center}
{\large\itshape\color{vinoInst}«La tecnología no es neutral: puede liberar o\\
dominar según quién la controla.»}\\[0.25cm]
{\footnotesize\color{vinoInst}— principio rector del área, inspirado en la filosofía de la liberación}
\end{center}

\vspace{0.9cm}
\noindent Este Plan de Área no se cierra en esta página: continúa cada vez que un
estudiante de la I.E. Sor María Juliana abre una guía, escribe en su cuaderno y
decide, con criterio propio, qué hacer con la tecnología que tiene en las manos.
\textbf{Formar autores, no usuarios} es el compromiso que las firmas siguientes
respaldan.

\vspace{1.1cm}
\begin{center}
{\color{vinoInst}\rule{\linewidth}{2pt}}\\[0.6cm]
{\Large\bfseries\color{vinoInst}Aprobación y Firma del Plan de Área 2026}\\[0.3cm]
{\normalsize Área de Tecnología e Informática}\\[0.1cm]
{\normalsize I.E. Sor María Juliana · Cartago, Valle del Cauca}\\[0.3cm]
{\color{vinoInst}\rule{\linewidth}{0.5pt}}
\end{center}

\vspace{0.8cm}

Los abajo firmantes, docentes y directivos del Área de Tecnología e Informática
de la Institución Educativa Sor María Juliana, certificamos que el presente
\textbf{Plan de Área 2026} ha sido elaborado en concordancia con la Ley General
de Educación (Ley 115/1994), los lineamientos curriculares del MEN, la Guía
No. 30 (2008), las Orientaciones Curriculares de Tecnología e Informática (2022),
la Ley 2383/2024 sobre educación socioemocional, y el Proyecto Educativo
Institucional (PEI) de la I.E. Sor María Juliana. Manifestamos nuestro compromiso
con su implementación fiel, con la evaluación liberadora y con la formación
integral de los estudiantes de esta comunidad educativa.

\vspace{1.8cm}

\begin{center}
\begin{tabular}{M{5.5cm} M{2cm} M{5.5cm}}
\rule{5.5cm}{0.5pt} & & \rule{5.5cm}{0.5pt} \\[0.3cm]
{\small\textbf{Equipo del Área de Tecnología e Informática}} & &
{\small\textbf{Coordinador/a Académico/a}} \\[0.1cm]
{\small Docentes del área} & &
{\small Coordinación Académica} \\[0.1cm]
{\small Área de Tecnología e Informática} & &
{\small I.E. Sor María Juliana} \\[0.1cm]
{\small C.C. \underline{\hspace{4cm}}} & &
{\small C.C. \underline{\hspace{4cm}}} \\
\end{tabular}
\end{center}

\vspace{1.8cm}

\begin{center}
\begin{tabular}{M{5.5cm} M{2cm} M{5.5cm}}
\rule{5.5cm}{0.5pt} & & \rule{5.5cm}{0.5pt} \\[0.3cm]
{\small\textbf{Rector/a}} & &
{\small\textbf{Representante del Consejo Académico}} \\[0.1cm]
{\small Rectoría} & &
{\small Consejo Académico Institucional} \\[0.1cm]
{\small I.E. Sor María Juliana} & &
{\small I.E. Sor María Juliana} \\[0.1cm]
{\small C.C. \underline{\hspace{4cm}}} & &
{\small C.C. \underline{\hspace{4cm}}} \\
\end{tabular}
\end{center}

\vspace{2.5cm}

\begin{center}
\begin{tcolorbox}[
  enhanced,
  colback=vinoClaro,
  colframe=vinoInst,
  boxrule=1pt,
  arc=8pt,
  width=0.80\linewidth,
  halign=center
]
\centering
{\small\textbf{Aprobado en Consejo Académico, Cartago (Valle del Cauca), junio de 2026.}}\\[0.3cm]
{\small I.E. Sor María Juliana · NIT \_\_\_\_\_\_\_\_\_\_ · DANE \_\_\_\_\_\_\_\_\_\_}\\[0.2cm]
{\small Plataforma Conéctate: alguarito.github.io/plataformaconectate}\\
{\small Repositorio de guías: \texttt{make guia-status} (actualizado en tiempo real)}
\end{tcolorbox}
\end{center}

\end{document}
